% Options for packages loaded elsewhere
\PassOptionsToPackage{unicode}{hyperref}
\PassOptionsToPackage{hyphens}{url}
\documentclass[
  11pt,
]{article}
\usepackage{xcolor}
\usepackage[margin=1in]{geometry}
\usepackage{amsmath,amssymb}
\setcounter{secnumdepth}{-\maxdimen} % remove section numbering
\usepackage{iftex}
\ifPDFTeX
  \usepackage[T1]{fontenc}
  \usepackage[utf8]{inputenc}
  \usepackage{textcomp} % provide euro and other symbols
\else % if luatex or xetex
  \usepackage{unicode-math} % this also loads fontspec
  \defaultfontfeatures{Scale=MatchLowercase}
  \defaultfontfeatures[\rmfamily]{Ligatures=TeX,Scale=1}
\fi
\usepackage{lmodern}
\ifPDFTeX\else
  % xetex/luatex font selection
  \setmainfont[]{Cambria}
  \setmonofont[]{Consolas}
  \setmathfont[]{Cambria Math}
\fi
% Use upquote if available, for straight quotes in verbatim environments
\IfFileExists{upquote.sty}{\usepackage{upquote}}{}
\IfFileExists{microtype.sty}{% use microtype if available
  \usepackage[]{microtype}
  \UseMicrotypeSet[protrusion]{basicmath} % disable protrusion for tt fonts
}{}
\makeatletter
\@ifundefined{KOMAClassName}{% if non-KOMA class
  \IfFileExists{parskip.sty}{%
    \usepackage{parskip}
  }{% else
    \setlength{\parindent}{0pt}
    \setlength{\parskip}{6pt plus 2pt minus 1pt}}
}{% if KOMA class
  \KOMAoptions{parskip=half}}
\makeatother
\usepackage{longtable,booktabs,array}
\newcounter{none} % for unnumbered tables
\usepackage{calc} % for calculating minipage widths
% Correct order of tables after \paragraph or \subparagraph
\usepackage{etoolbox}
\makeatletter
\patchcmd\longtable{\par}{\if@noskipsec\mbox{}\fi\par}{}{}
\makeatother
% Allow footnotes in longtable head/foot
\IfFileExists{footnotehyper.sty}{\usepackage{footnotehyper}}{\usepackage{footnote}}
\makesavenoteenv{longtable}
\setlength{\emergencystretch}{3em} % prevent overfull lines
\providecommand{\tightlist}{%
  \setlength{\itemsep}{0pt}\setlength{\parskip}{0pt}}
% Center all longtables
\setlength{\LTleft}{0pt plus 1fill}
\setlength{\LTright}{0pt plus 1fill}
\usepackage{bookmark}
\IfFileExists{xurl.sty}{\usepackage{xurl}}{} % add URL line breaks if available
\urlstyle{same}
\hypersetup{
  hidelinks,
  pdfcreator={LaTeX via pandoc}}

\author{}
\date{}

\begin{document}

\begin{center}
{\LARGE\bfseries The Patient Compass: Domain Walls, Oxygen Gating, and the\\
3-D Survival Manifold of 198,862 Tumors}

\vspace{1.5em}

{\large Antonio P.\ Matos}\\[0.4em]
{\small ORCID: \href{https://orcid.org/0009-0002-0722-3752}{0009-0002-0722-3752}}\\[0.3em]
April 13, 2026\\[0.3em]
{\small DOI: \href{https://doi.org/10.5281/zenodo.19561701}{10.5281/zenodo.19561701}}\\[0.3em]
Independent Researcher; Fancyland LLC / Lattice OS\\[0.3em]
Preprint v1.0

\vspace{0.8em}

{\small\textbf{Keywords:} hallmarks of cancer, domain walls, frustrated proliferation, tissue oxygenation, forbidden topology, Cox proportional hazards, Ising model, pan-cancer survival, somatic mutations, cBioPortal}
\end{center}

\textbf{Companion papers:}

\begin{itemize}
\tightlist
\item
  {[}1{]} ``Active Transport on the Prime Gas: Flat-Band Condensation,
  the Rabi Phase Transition, and the Arithmetic Qubit,'' Matos (2026),
  DOI:
  \href{https://doi.org/10.5281/zenodo.19243258}{10.5281/zenodo.19243258}
\item
  {[}2{]} ``The Arithmetic Black Hole: Softmax Thermodynamics and the
  Four Eigenvalue Laws of the Prime Gas,'' Matos (2026), DOI:
  \href{https://doi.org/10.5281/zenodo.19442006}{10.5281/zenodo.19442006}
\item
  {[}3{]} ``The Unity Clock: Effective Dimensional Collapse of the
  Addition-Multiplication Commutator on Coprime Lattices,'' Matos
  (2026), DOI:
  \href{https://doi.org/10.5281/zenodo.19478727}{10.5281/zenodo.19478727}
\end{itemize}

\textbf{Companion patents:}

\begin{itemize}
\tightlist
\item
  {[}4{]} ``Fault-Injection-Immune Computational Unit Using Primorial
  Coprime Residue Topology,'' Matos (2026), U.S. Provisional Patent
  Application No.~64/031,440 (filed April 7, 2026)
\item
  {[}5{]} ``Holographic Eigen-Solver Using QM Boundary Projection on
  Coprime Residue Lattices,'' Matos (2026), U.S. Provisional Patent
  Application No.~64/033,689 (filed April 8, 2026)
\end{itemize}

\begin{center}\rule{0.5\linewidth}{0.5pt}\end{center}

\subsection{Abstract}\label{abstract}

We analyzed somatic mutation data from \textbf{198,862 tumors} across
\textbf{378 independent studies} (cBioPortal), mapping 45 cancer driver
genes to the eight hallmarks of cancer (Hanahan \& Weinberg, 2000;
2011). We began with a pair completeness hypothesis --- grouping the
eight hallmarks into four functional pairs and counting how many pairs
are simultaneously active --- and systematically tested, extended, and
ultimately falsified it as a survival framework. What survived this
process is a three-part model that is stronger than what we started
with.

\textbf{Domain walls predict survival.} Re-parameterizing the 8-bit
hallmark vector into domain-wall coordinates --- where each hallmark
pair is decomposed into \emph{coupled}, \emph{A-only}, and \emph{B-only}
states --- yields the dominant survival features. The frustrated
proliferator (proliferative signaling without invasion) is the single
strongest hallmark-derived predictor: HR \(= 1.81\) in hypoxic tissue
(\(p < 10^{-145}\)). TP53-coupled growth suppression + death evasion
increases hazard by 32\% (\(p < 10^{-83}\)). Hypoxic coupling
(simultaneous immortality + metabolism) is massively protective (HR
\(= 0.40\)). These wall features outperform both raw hallmark counts and
pair completeness (\(\Delta\mathrm{AIC} = -220\)).

\textbf{Tissue oxygenation gates all hallmark effects.} Classifying
patients by host tissue oxygen tension reveals that oxygenation is the
strongest independent survival predictor in the multivariate Cox model
(high-O\(_2\) HR \(= 0.59\); low-O\(_2\) HR \(= 1.32\)). The survival
effect of Genome Instability \emph{inverts} by oxygenation: harmful in
normoxic tissue (\(p = 0.001\)), protective in hypoxic tissue
(\(p = 8.3 \times 10^{-38}\)). Forbidden neighbors on the hallmark
hypercube --- the 94 states never observed in the data --- are
protective in hypoxic tissue but weakly lethal in normoxia, with the
sign flip confirmed across cancer types (\(p = 0.020\)).

\textbf{Eight hallmarks collapse to three dimensions.} Marchenko-Pastur
random matrix analysis finds exactly 3 eigenvalues above the noise
threshold in all oxygenation regimes. A 3-composite-score + O\(_2\)
model (\(C_\mathrm{CV} = 0.6101 \pm 0.0055\), 5 parameters) captures
96.3\% of the predictive power of the full 25-parameter interaction
model, with zero meaningful overfit. The effective cancer survival space
is a TP53 axis, an O\(_2\)-adaptive axis, and a Genome Instability
toggle.

The pair completeness hypothesis --- our starting point --- ranks 19th
of 105 possible hallmark pairings for survival separation. The
biological pairs are a human interpretive framework, not a thermodynamic
eigenvector. What the data actually support is a 7-parameter energy
function combining log-prevalence, forbidden-neighbor topology, hallmark
sum, and O\(_2\) interactions, which captures 87.9\% of hallmark model
discriminability. The model, gene mapping, O\(_2\) classification, and
an interactive patient-facing report are freely available.

\begin{center}\rule{0.5\linewidth}{0.5pt}\end{center}

\subsection{1. Introduction}\label{introduction}

\subsubsection{1.1 The Problem}\label{the-problem}

A patient diagnosed with cancer faces a paradox of information. Genomic
sequencing generates thousands of data points per tumor, yet the
patient's fundamental questions remain unanswered in accessible terms:
\emph{What has my cancer learned to do? What has it not learned? Does
that pattern mean anything for how I will do?}

The hallmarks of cancer framework (Hanahan \& Weinberg, 2000; updated
2011) provides a natural organizing principle. Rather than presenting
patients and clinicians with a list of mutated genes, we can ask:
\emph{which of the eight core cancer capabilities has this tumor
acquired, and which has it not?}

This reframing --- from genes to capabilities --- is not merely
pedagogical. If the \emph{pattern} of capability acquisition carries
survival information, then it is clinically actionable: it tells
oncologists which patients to monitor more closely, which drug targets
may matter more, and which patients may have tumors with exploitable
gaps.

\subsubsection{1.2 The Pair Hypothesis}\label{the-pair-hypothesis}

The eight hallmarks are not independent. Proliferative signaling and
invasion/metastasis represent two distinct growth strategies. Growth
suppression evasion and replicative immortality both involve overcoming
restraints on cell division. Resisting cell death and immune evasion
both involve surviving threats. Genome instability and deregulated
metabolism both involve altering the cell's internal machinery.

These functional tensions motivate the organization of the eight
hallmarks into four pairs:

{\def\LTcaptype{none} % do not increment counter
\begin{longtable}[]{@{}lll@{}}
\toprule\noalign{}
Pair & Hallmark A & Hallmark B \\
\midrule\noalign{}
\endhead
\bottomrule\noalign{}
\endlastfoot
1 & Proliferative Signaling & Invasion \& Metastasis \\
2 & Growth Suppression Evasion & Replicative Immortality \\
3 & Resisting Cell Death & Immune Evasion \\
4 & Genome Instability & Deregulated Metabolism \\
\end{longtable}
}

A pair is \textbf{complete} when a tumor has acquired mutations in
\emph{both} hallmarks of the pair. We count the number of complete pairs
(0 through 4) as an integer summary of hallmark balance.

The theoretical motivation for this specific pairing comes from the
spectral structure of hallmark co-occurrence matrices, described in
companion papers {[}1--3{]}, where these four pairs emerge as the
natural involution of a symmetry operator on the hallmark space. In this
paper, we set aside the spectral theory entirely and ask the purely
empirical question: \textbf{does pair completeness predict survival?}

\subsubsection{1.3 Scope}\label{scope}

We analyze every study on cBioPortal with mutation data and at least 50
samples, excluding cell lines and hematologic malignancies (whose
hallmark biology differs from solid tumors). The resulting dataset spans
198,862 tumors across 378 studies from dozens of institutions worldwide,
representing one of the largest pan-cancer mutation analyses conducted
to date.

\begin{center}\rule{0.5\linewidth}{0.5pt}\end{center}

\subsection{2. Methods}\label{methods}

\subsubsection{2.1 Data Source}\label{data-source}

All data were obtained from the cBioPortal for Cancer Genomics public
API (https://www.cbioportal.org/api), accessed in March 2026. cBioPortal
aggregates published cancer genomics datasets with standardized mutation
calls and clinical annotations. No patient-identifiable information was
accessed or stored.

\subsubsection{2.2 Study Selection}\label{study-selection}

We retrieved all studies from cBioPortal and applied the following
inclusion criteria:

\begin{itemize}
\tightlist
\item
  \textbf{Mutation data available:} Study must have at least one
  molecular profile of type \texttt{MUTATION\_EXTENDED}.
\item
  \textbf{Minimum sample size:} At least 50 samples.
\item
  \textbf{Solid tumors only:} Studies were excluded if the cancer type
  or study name matched hematologic terms (AML, ALL, CLL, lymphoma,
  leukemia, myeloid, myeloma, etc.).
\item
  \textbf{No cell lines:} Studies with ``cellline'' or ``ccle'' in the
  study identifier were excluded.
\end{itemize}

378 studies met all criteria. Patient deduplication was performed within
each study using patient identifiers to avoid counting multi-sample
patients more than once.

\subsubsection{2.3 Gene Panel}\label{gene-panel}

We selected 45 well-characterized cancer driver genes that collectively
cover the eight hallmarks of cancer. The gene-to-hallmark mapping was
constructed from established functional annotations in the cancer
biology literature:

\textbf{Proliferative Signaling (Hallmark 0):} KRAS, BRAF, NRAS, ERBB2,
ERBB3, MAP2K1, FGFR2, NF1, PIK3CA\\
\textbf{Growth Suppression Evasion (Hallmark 1):} TP53, APC, RB1,
CDKN2A, SMAD4, SMAD2, SMAD3, TGFBR2, FBXW7, ACVR2A, NOTCH1, SOX9,
TCF7L2, RNF43, AMER1, CTNNB1\\
\textbf{Resisting Cell Death (Hallmark 2):} PIK3CA, TP53, PTEN, CREBBP,
TSC2\\
\textbf{Replicative Immortality (Hallmark 3):} ATRX\\
\textbf{Invasion \& Metastasis (Hallmark 4):} CTNNB1, CDH1, FAT1, FAT4\\
\textbf{Immune Evasion (Hallmark 5):} JAK1\\
\textbf{Genome Instability (Hallmark 6):} MLH1, MSH2, MSH6, PMS2, POLE,
BRCA2, ATM, ARID1A, KMT2C, KMT2D, ATRX\\
\textbf{Deregulated Metabolism (Hallmark 7):} IDH1, GNAS

Some genes (PIK3CA, TP53, CTNNB1, ATRX) map to multiple hallmarks,
reflecting their pleiotropic roles.

A hallmark is considered \textbf{active} in a tumor if any gene mapped
to that hallmark carries a somatic mutation in that tumor's sequencing
data.

\subsubsection{2.4 Pair Completeness}\label{pair-completeness}

For each tumor, we compute the 8-bit hallmark activation vector
\(\mathbf{h} = (h_0, h_1, \ldots, h_7) \in \{0,1\}^8\) and count the
number of complete pairs:

\[\text{PairCount} = \sum_{(a,b) \in \mathcal{P}} h_a \cdot h_b \tag{1}\]

where \(\mathcal{P} = \{(0,4), (1,3), (2,5), (6,7)\}\) is the set of
four hallmark pairs. This yields an integer in \(\{0, 1, 2, 3, 4\}\).

\subsubsection{2.5 Symmetry Measure}\label{symmetry-measure}

For each study independently, we construct the hallmark co-occurrence
matrix \(D_h\) (log-odds ratios of pairwise hallmark co-occurrence) and
derive a spectral symmetry measure \(\sigma_z\) for each patient,
following the method described in {[}1{]}. This continuous measure
captures how a tumor's hallmark pattern aligns with the dominant
spectral modes of the study's hallmark co-occurrence structure.

For survival analysis, patients within each study are split at the
median \(\sigma_z\) value into ``high'' and ``low'' symmetry groups.

\subsubsection{2.6 Survival Analysis}\label{survival-analysis}

Overall survival (OS) data were obtained from cBioPortal clinical
annotations where available (OS\_STATUS and OS\_MONTHS fields). A study
was included in survival analysis if it had at least 30 patients with
survival data and at least 10 death events.

For each eligible study, we performed a Mann-Whitney U test comparing OS
between the high and low \(\sigma_z\) groups. The Mann-Whitney test is
non-parametric and does not assume proportional hazards or any
distributional form, making it robust for the heterogeneous datasets in
this analysis.

We report: - The Mann-Whitney z-statistic and two-sided p-value for each
study - The median OS difference (Δ) between groups in months - The
\textbf{direction}: whether higher symmetry is associated with longer or
shorter survival

For multivariate survival modeling (§3.19 onward), we used Cox
proportional hazards regression with Efron's method for tied event times
and an L2 (Ridge) penalty (penalizer \(= 0.01\)), implemented in
lifelines 0.30. The domain wall Cox model (§3.23) uses 10 wall features
+ 2 tissue O\(_2\) indicators = 12 covariates, fitted on 48,675 patients
with survival data and an assigned O\(_2\) tier. Model comparison uses
Harrell's C-index for discrimination and AIC for parsimony. All
concordance values are verified by 5-fold cross-validation (§3.29).

\subsubsection{2.7 Meta-Analysis}\label{meta-analysis}

To combine evidence across studies, we applied Stouffer's method in
three forms:

\begin{enumerate}
\def\labelenumi{\arabic{enumi}.}
\tightlist
\item
  \textbf{Unweighted Stouffer:} \(Z = \sum z_i / \sqrt{k}\), where \(k\)
  is the number of studies tested.
\item
  \textbf{Weighted Stouffer:}
  \(Z_w = \sum w_i z_i / \sqrt{\sum w_i^2}\), with weights
  \(w_i = \sqrt{N_i}\).
\item
  \textbf{Direction-agnostic Stouffer:}
  \(Z_{|z|} = \sum |z_i| / \sqrt{k}\). This tests whether pair
  completeness carries \emph{any} survival information, regardless of
  direction.
\end{enumerate}

The direction-agnostic test is the most appropriate for the pan-cancer
setting, because we hypothesize that pair completeness is informative
but do not claim it has a universal direction.

\begin{center}\rule{0.5\linewidth}{0.5pt}\end{center}

\subsection{3. Results}\label{results}

\emph{The first five sections describe the raw structure of the data:
how 198,862 tumors distribute across hallmark space, which states are
forbidden, and what thermodynamic law governs the distribution.}

\subsubsection{3.1 The Lopsidedness of
Cancer}\label{the-lopsidedness-of-cancer}

\textbf{Finding: 85.8\% of tumors have zero complete hallmark pairs.}

The distribution of pair completeness across all 198,862 tumors is
dramatically skewed toward zero:

{\def\LTcaptype{none} % do not increment counter
\begin{longtable}[]{@{}lll@{}}
\toprule\noalign{}
Complete Pairs & N & \% \\
\midrule\noalign{}
\endhead
\bottomrule\noalign{}
\endlastfoot
0 & 170,622 & 85.8\% \\
1 & 22,412 & 11.3\% \\
2 & 4,711 & 2.4\% \\
3 & 864 & 0.4\% \\
4 & 253 & 0.1\% \\
\end{longtable}
}

This is a structural finding about how cancer develops. Despite having
eight capabilities to acquire and a theoretical maximum of four complete
pairs, the overwhelming majority of tumors are \emph{lopsided} --- they
acquire capabilities asymmetrically, activating one side of a functional
pair without activating the other.

This has clinical implications. A lopsided tumor --- one with zero or
one complete pairs --- has \textbf{biological blind spots}: capabilities
it has not acquired. These blind spots may represent therapeutic
vulnerabilities. A tumor that has acquired proliferative signaling but
not invasion capacity, for example, may be locally dangerous but less
likely to metastasize, and treatments targeting proliferation may be
more effective precisely because the tumor lacks the compensatory
capability of its paired hallmark.

\subsubsection{3.2 State Space
Concentration}\label{state-space-concentration}

Each tumor's hallmark profile is an 8-bit binary vector
\(\sigma \in \{0,1\}^8\), yielding \(2^8 = 256\) possible states. We
treat these states as vertices of an 8-dimensional binary hypercube,
where Hamming-adjacent states (differing in exactly one hallmark) are
connected by edges. This lattice formulation --- analyzing observed
versus forbidden vertices, Hamming neighborhoods, energy landscapes, and
spectral properties of the adjacency structure --- follows the coprime
residue lattice methodology developed in the companion papers
{[}1--3{]}, where the same framework revealed universal spectral
concentration on algebraic lattices. Here we apply it to the biological
lattice of cancer hallmarks.

Only 162 of the 256 possible hallmark states were observed in the data.
The distribution is highly concentrated:

{\def\LTcaptype{none} % do not increment counter
\begin{longtable}[]{@{}
  >{\raggedright\arraybackslash}p{(\linewidth - 8\tabcolsep) * \real{0.2000}}
  >{\raggedright\arraybackslash}p{(\linewidth - 8\tabcolsep) * \real{0.2000}}
  >{\raggedright\arraybackslash}p{(\linewidth - 8\tabcolsep) * \real{0.2000}}
  >{\raggedright\arraybackslash}p{(\linewidth - 8\tabcolsep) * \real{0.2000}}
  >{\raggedright\arraybackslash}p{(\linewidth - 8\tabcolsep) * \real{0.2000}}@{}}
\toprule\noalign{}
\begin{minipage}[b]{\linewidth}\raggedright
Rank
\end{minipage} & \begin{minipage}[b]{\linewidth}\raggedright
State
\end{minipage} & \begin{minipage}[b]{\linewidth}\raggedright
Hallmarks Active
\end{minipage} & \begin{minipage}[b]{\linewidth}\raggedright
N
\end{minipage} & \begin{minipage}[b]{\linewidth}\raggedright
\%
\end{minipage} \\
\midrule\noalign{}
\endhead
\bottomrule\noalign{}
\endlastfoot
1 & 01100000 & Growth Suppression Evasion, Resisting Cell Death & 41,198
& 20.7\% \\
2 & 11100000 & Proliferative Signaling, Growth Suppression Evasion,
Resisting Cell Death & 29,939 & 15.1\% \\
3 & 10000000 & Proliferative Signaling only & 12,504 & 6.3\% \\
4 & 11100010 & Proliferative Signaling, Growth Suppression Evasion,
Resisting Cell Death, Genome Instability & 11,373 & 5.7\% \\
5 & 01100010 & Growth Suppression Evasion, Resisting Cell Death, Genome
Instability & 11,040 & 5.6\% \\
6 & 00000010 & Genome Instability only & 10,421 & 5.2\% \\
7 & 10100000 & Proliferative Signaling, Resisting Cell Death & 7,609 &
3.8\% \\
8 & 01000000 & Growth Suppression Evasion only & 6,233 & 3.1\% \\
9 & 11101010 & Proliferative Signaling, Growth Suppression Evasion,
Resisting Cell Death, Invasion \& Metastasis, Genome Instability & 4,875
& 2.5\% \\
10 & 01101000 & Growth Suppression Evasion, Resisting Cell Death,
Replicative Immortality & 4,089 & 2.1\% \\
\end{longtable}
}

The top 10 states account for 70.0\% of all patients (139,281 /
198,862). Patients in states with N ≥ 20 account for 99.8\% of the
total.

\textbf{The most common single state} --- Growth Suppression Evasion +
Resisting Cell Death (01100000) --- represents one in five tumors across
all cancer types. This pattern is driven by TP53 mutations, which map to
both hallmarks and are the most common somatic mutation in human cancer.

Note that none of the top 3 states contain \emph{any} complete pair. The
most common state with a complete pair is rank 9 (11101010, with pair 1
complete: Proliferative Signaling + Invasion \& Metastasis),
representing only 2.5\% of patients.

\subsubsection{3.3 Hallmark Prevalence Across Cancer
Types}\label{hallmark-prevalence-across-cancer-types}

The eight hallmarks vary widely in prevalence across the 378 studies:

{\def\LTcaptype{none} % do not increment counter
\begin{longtable}[]{@{}llll@{}}
\toprule\noalign{}
Hallmark & Mean Prevalence & Median & Range \\
\midrule\noalign{}
\endhead
\bottomrule\noalign{}
\endlastfoot
Growth Suppression Evasion & 44.8\% & 42.6\% & 0--98.8\% \\
Resisting Cell Death & 43.0\% & 41.3\% & 0--98.8\% \\
Proliferative Signaling & 29.6\% & 23.8\% & 0--100\% \\
Genome Instability & 25.6\% & 22.2\% & 0--97.2\% \\
Invasion \& Metastasis & 13.0\% & 9.1\% & 0--61.5\% \\
Deregulated Metabolism & 4.2\% & 1.8\% & 0--100\% \\
Replicative Immortality & 4.0\% & 2.2\% & 0--82.0\% \\
Immune Evasion & 1.8\% & 1.1\% & 0--42.0\% \\
\end{longtable}
}

The extreme asymmetry in hallmark prevalence explains the lopsidedness
finding. Growth Suppression Evasion and Resisting Cell Death --- both
driven heavily by TP53 --- are active in nearly half of all tumors. But
their paired hallmarks (Replicative Immortality at 4.0\% and Immune
Evasion at 1.8\%) are rarely active. Pairs 2 and 3 are therefore almost
never complete, not because tumors avoid completing them, but because
the individual hallmarks on one side of the pair are simply rare in
mutation data.

This observation is important for interpretation: \textbf{the
lopsidedness of cancer, as measured here, is partly a property of the
gene panel.} A broader panel --- including epigenetic, copy number, and
expression data --- might detect hallmark activations that mutation data
alone misses. This is a limitation we address in §4.3.

\subsubsection{3.4 Thermodynamic Distribution of Pair
Completeness}\label{thermodynamic-distribution-of-pair-completeness}

The pair completeness distribution (§3.1) follows an exponential decay
with remarkable precision. Fitting \(\ln P(k) = -\beta k + C\) by
ordinary least squares gives:

\[\beta = 1.6283, \quad R^2 = 0.9935 \tag{2}\]

The probability of a tumor having \(k\) complete pairs is well described
by a Boltzmann distribution:

\[P(k) \propto e^{-\beta k} \tag{3}\]

with effective inverse temperature \(\beta = 1.63\) and effective
temperature \(T = 1/\beta = 0.614\).

Each additional complete pair costs approximately 1.63 nats of selective
free energy. The quality of the fit (\(R^2 > 0.99\)) establishes that
pair acquisition is not a random combinatorial process --- if it were,
the distribution would follow a binomial, not an exponential. Instead,
the distribution obeys thermodynamic statistics, where each successive
pair faces a free-energy barrier of similar magnitude.

This Boltzmann law governs the distribution of 198,862 tumors across
cancer types and institutions. It is the single strongest quantitative
regularity in our dataset.

\subsubsection{3.5 Selection Pressure Against Pair
Completion}\label{selection-pressure-against-pair-completion}

Of 256 possible hallmark states, 94 (36.7\%) are never observed in
198,862 tumors. These ``forbidden'' states are not uniformly distributed
across the phase space --- they are concentrated where specific pairs
are present.

\textbf{Forbidden rate by pair-completeness level:}

{\def\LTcaptype{none} % do not increment counter
\begin{longtable}[]{@{}
  >{\raggedright\arraybackslash}p{(\linewidth - 10\tabcolsep) * \real{0.1667}}
  >{\raggedright\arraybackslash}p{(\linewidth - 10\tabcolsep) * \real{0.1667}}
  >{\raggedright\arraybackslash}p{(\linewidth - 10\tabcolsep) * \real{0.1667}}
  >{\raggedright\arraybackslash}p{(\linewidth - 10\tabcolsep) * \real{0.1667}}
  >{\raggedright\arraybackslash}p{(\linewidth - 10\tabcolsep) * \real{0.1667}}
  >{\raggedright\arraybackslash}p{(\linewidth - 10\tabcolsep) * \real{0.1667}}@{}}
\toprule\noalign{}
\begin{minipage}[b]{\linewidth}\raggedright
Complete Pairs
\end{minipage} & \begin{minipage}[b]{\linewidth}\raggedright
Total States
\end{minipage} & \begin{minipage}[b]{\linewidth}\raggedright
Observed
\end{minipage} & \begin{minipage}[b]{\linewidth}\raggedright
Forbidden
\end{minipage} & \begin{minipage}[b]{\linewidth}\raggedright
\% Forbidden
\end{minipage} & \begin{minipage}[b]{\linewidth}\raggedright
\% of Patients
\end{minipage} \\
\midrule\noalign{}
\endhead
\bottomrule\noalign{}
\endlastfoot
0 & 81 & 59 & 22 & 27.2\% & 85.8\% \\
1 & 108 & 64 & 44 & 40.7\% & 11.3\% \\
2 & 54 & 29 & 25 & 46.3\% & 2.4\% \\
3 & 12 & 9 & 3 & 25.0\% & 0.4\% \\
4 & 1 & 1 & 0 & 0.0\% & 0.1\% \\
\end{longtable}
}

The forbidden rate increases from 27\% at zero pairs to 46\% at two
pairs, confirming that higher pair completeness faces stronger selection
pressure.

More revealing is the pair-specific decomposition. For each pair, we
compute the fraction of states that are forbidden when that pair is
complete versus when it is not:

\textbf{Pair-specific forbidden rate differential:}

{\def\LTcaptype{none} % do not increment counter
\begin{longtable}[]{@{}
  >{\raggedright\arraybackslash}p{(\linewidth - 6\tabcolsep) * \real{0.2500}}
  >{\raggedright\arraybackslash}p{(\linewidth - 6\tabcolsep) * \real{0.2500}}
  >{\raggedright\arraybackslash}p{(\linewidth - 6\tabcolsep) * \real{0.2500}}
  >{\raggedright\arraybackslash}p{(\linewidth - 6\tabcolsep) * \real{0.2500}}@{}}
\toprule\noalign{}
\begin{minipage}[b]{\linewidth}\raggedright
Pair
\end{minipage} & \begin{minipage}[b]{\linewidth}\raggedright
Description
\end{minipage} & \begin{minipage}[b]{\linewidth}\raggedright
Δ Forbidden Rate
\end{minipage} & \begin{minipage}[b]{\linewidth}\raggedright
Role
\end{minipage} \\
\midrule\noalign{}
\endhead
\bottomrule\noalign{}
\endlastfoot
2 & Growth Suppression Evasion + Replicative Immortality & +24.0 pp &
\textbf{Barrier} --- highest thermodynamic cost \\
3 & Resisting Cell Death + Immune Evasion & +9.4 pp & Moderate
barrier \\
1 & Proliferative Signaling + Invasion \& Metastasis & −1.0 pp &
Neutral \\
4 & Genome Instability + Deregulated Metabolism & −13.5 pp &
\textbf{Enabler} --- lowers system energy \\
\end{longtable}
}

Pair 2 is the hardest to achieve: states containing it are 24 percentage
points more likely to be forbidden. Biologically, this reflects the
mechanistic independence of its components --- evading growth
suppression (TP53/RB pathway) and achieving replicative immortality
(telomerase/ATRX) require two unrelated evolutionary steps.

Pair 4 is an enabler: states containing it are 13.5 percentage points
\emph{less} likely to be forbidden. Biologically, this reflects
mechanistic coupling --- genome instability (e.g., mismatch repair
deficiency) directly drives metabolic reprogramming, and vice versa.
IDH1 mutations in glioma exemplify this: a single mutation affects both
hallmarks simultaneously.

This pair-specific asymmetry extends to the symmetric states --- the 16
states where each hallmark pair is either fully complete or fully
absent. Six of these 16 states are forbidden. All 6 forbidden symmetric
states require the hard pairs (Pair 2 or Pair 3) without the support of
enabler pairs. The 10 surviving symmetric states all contain Pair 4 or
Pair 1. Symmetric exclusion is driven by pair difficulty, not by
symmetry per se.

\begin{center}\rule{0.5\linewidth}{0.5pt}\end{center}

\emph{The pair completeness hypothesis groups the eight hallmarks into
four functional pairs. This act builds the Ising model, tests it
exhaustively, and follows the evidence to its breaking point.}

\subsubsection{3.6 Pan-Cancer Survival
Stratification}\label{pan-cancer-survival-stratification}

Of 378 studies, 181 had sufficient survival data for analysis (≥30
patients with OS data, ≥10 death events).

\textbf{29 studies (16.0\%) showed statistically significant survival
stratification at p \textless{} 0.05} (Appendix A). An additional 11
studies reached p \textless{} 0.10.

Under the null hypothesis (pair completeness carries no survival
information), we would expect approximately 9 studies (5\% of 181) to
reach p \textless{} 0.05 by chance. Observing 29 represents a 3.2-fold
enrichment.

The direction-agnostic Stouffer meta-analysis, which tests whether pair
completeness carries \emph{any} survival signal across all 181 studies
regardless of direction, yields:

\[Z_{|z|} = +15.29, \quad p = 8.9 \times 10^{-53} \tag{4}\]

This is not marginal. Pair completeness carries survival information
across the full breadth of solid tumor types in this dataset with
overwhelming statistical confidence.

The unweighted Stouffer test, which accounts for direction, yields z =
+2.25, p = 0.025, indicating a net positive direction (higher symmetry
associated with longer survival) but weaker due to cancellation between
cancer types with opposite directions.

\subsubsection{3.7 Glioma: Six Independent
Replications}\label{glioma-six-independent-replications}

The strongest and most consistent signal comes from diffuse glioma,
where \textbf{six independent cohorts} all reach p \textless{} 0.05 in
the same direction:

{\def\LTcaptype{none} % do not increment counter
\begin{longtable}[]{@{}
  >{\raggedright\arraybackslash}p{(\linewidth - 10\tabcolsep) * \real{0.1667}}
  >{\raggedright\arraybackslash}p{(\linewidth - 10\tabcolsep) * \real{0.1667}}
  >{\raggedright\arraybackslash}p{(\linewidth - 10\tabcolsep) * \real{0.1667}}
  >{\raggedright\arraybackslash}p{(\linewidth - 10\tabcolsep) * \real{0.1667}}
  >{\raggedright\arraybackslash}p{(\linewidth - 10\tabcolsep) * \real{0.1667}}
  >{\raggedright\arraybackslash}p{(\linewidth - 10\tabcolsep) * \real{0.1667}}@{}}
\toprule\noalign{}
\begin{minipage}[b]{\linewidth}\raggedright
Study
\end{minipage} & \begin{minipage}[b]{\linewidth}\raggedright
Source
\end{minipage} & \begin{minipage}[b]{\linewidth}\raggedright
N
\end{minipage} & \begin{minipage}[b]{\linewidth}\raggedright
Deaths
\end{minipage} & \begin{minipage}[b]{\linewidth}\raggedright
Δ OS
\end{minipage} & \begin{minipage}[b]{\linewidth}\raggedright
p-value
\end{minipage} \\
\midrule\noalign{}
\endhead
\bottomrule\noalign{}
\endlastfoot
Diffuse Glioma & GLASS Consortium & 256 & 192 & +29.0 mo & 3.4 × 10⁻⁷ \\
Diffuse Glioma & GLASS Consortium, Nature 2019 & 189 & 141 & +31.5 mo &
3.0 × 10⁻⁷ \\
Glioma & MSK, Clin Cancer Res 2019 & 784 & 287 & +10.6 mo & 1.3 ×
10⁻⁵ \\
Diffuse Glioma & TCGA GDC, 2025 & 462 & 104 & +7.7 mo & 2.5 × 10⁻³ \\
Brain Lower Grade Glioma & TCGA, PanCancer Atlas & 462 & 106 & +7.5 mo &
3.6 × 10⁻³ \\
Brain Lower Grade Glioma & TCGA, Firehose Legacy & 262 & 63 & +10.1 mo &
4.4 × 10⁻³ \\
\end{longtable}
}

In every glioma cohort, higher hallmark balance is associated with
\textbf{longer} survival. The GLASS Consortium cohorts show the most
dramatic separation: approximately 30 months of median survival
difference. The MSK cohort, the largest (N=784), shows 10.6 months.

This pattern replicates across three independent institutions (GLASS,
MSK, TCGA), three time periods, and both lower-grade and diffuse
histologies. To our knowledge, this is the first demonstration that a
mutation-derived hallmark balance measure stratifies glioma survival
with this degree of consistency.

\subsubsection{3.8 The Direction Split}\label{the-direction-split}

Among the 29 significant studies, the direction of the survival
association splits by cancer type:

\textbf{Higher balance → Longer survival (18 studies):} - Diffuse Glioma
(6 cohorts): +7.5 to +31.5 months - Melanoma (2 cohorts): +13.8 to +14.9
months - Breast Cancer (2 cohorts): +7.2 to +9.8 months - Prostate
Cancer (1 cohort): +5.2 months - Uterine Endometrial (1 cohort): +8.4
months - Uterine Carcinosarcoma (1 cohort): +13.5 months - Thyroid,
Anaplastic (1 cohort): +40.6 months - Bladder Urothelial (1 cohort):
+2.6 months - TMB/Immunotherapy mixed (1 cohort): +3.0 months -
MSK-IMPACT 50K mixed (1 cohort): +0.6 months

\textbf{Higher balance → Shorter survival (11 studies):} - MSK-CHORD
mixed (1 cohort, N=20,914): −2.6 months - Colorectal Cancer (2 cohorts):
−4.6 to −29.6 months - Ovarian Cancer (2 cohorts): −7.6 to −14.5 months
- Pancreatic Cancer (2 cohorts): −7.5 to −9.6 months -
Cholangiocarcinoma (1 cohort): −11.9 months - Adrenocortical Carcinoma
(1 cohort): −29.8 months - Prostate Adenocarcinoma (1 cohort): −4.5
months - IDH-mutant Glioma (1 cohort): −11.2 months

This directional split is not noise --- it is biology. The cancers where
balance predicts worse outcomes (colorectal, ovarian, pancreatic) are
predominantly gastrointestinal and gynecologic malignancies where
activating multiple hallmark pathways simultaneously may reflect a more
advanced, treatment-resistant disease state. The cancers where balance
predicts better outcomes (glioma, melanoma, breast) may represent tumors
where symmetric capability acquisition reflects a more differentiated,
less aggressive phenotype --- or where the immune system is more
effective against tumors that present a broader antigenic surface.

The single exception within the glioma cohorts --- the IDH-mutant
diffuse glioma study (MSK, p = 0.036, Δ = −11.2 months) --- is itself
biologically informative. IDH mutations fundamentally alter glioma
metabolism and epigenetics, creating a distinct disease entity. That the
compass direction flips specifically for IDH-mutant gliomas, while being
consistent across six other glioma cohorts, suggests the pair measure is
capturing real molecular subtype differences.

\subsubsection{3.9 Pair-Specific Survival}\label{pair-specific-survival}

Each of the four hallmark pairs individually carries survival
information. We performed pooled Mann-Whitney U tests comparing median
overall survival between patients with complete versus incomplete pairs
across all studies with survival data:

\footnotesize

{\def\LTcaptype{none} % do not increment counter
\begin{longtable}[]{@{}
  >{\raggedright\arraybackslash}p{(\linewidth - 14\tabcolsep) * \real{0.1250}}
  >{\raggedright\arraybackslash}p{(\linewidth - 14\tabcolsep) * \real{0.1250}}
  >{\raggedright\arraybackslash}p{(\linewidth - 14\tabcolsep) * \real{0.1250}}
  >{\raggedright\arraybackslash}p{(\linewidth - 14\tabcolsep) * \real{0.1250}}
  >{\raggedright\arraybackslash}p{(\linewidth - 14\tabcolsep) * \real{0.1250}}
  >{\raggedright\arraybackslash}p{(\linewidth - 14\tabcolsep) * \real{0.1250}}
  >{\raggedright\arraybackslash}p{(\linewidth - 14\tabcolsep) * \real{0.1250}}
  >{\raggedright\arraybackslash}p{(\linewidth - 14\tabcolsep) * \real{0.1250}}@{}}
\toprule\noalign{}
\begin{minipage}[b]{\linewidth}\raggedright
Pair
\end{minipage} & \begin{minipage}[b]{\linewidth}\raggedright
Complete N
\end{minipage} & \begin{minipage}[b]{\linewidth}\raggedright
Incomplete N
\end{minipage} & \begin{minipage}[b]{\linewidth}\raggedright
Median OS (Complete)
\end{minipage} & \begin{minipage}[b]{\linewidth}\raggedright
Median OS (Incomplete)
\end{minipage} & \begin{minipage}[b]{\linewidth}\raggedright
Δ OS
\end{minipage} & \begin{minipage}[b]{\linewidth}\raggedright
z
\end{minipage} & \begin{minipage}[b]{\linewidth}\raggedright
p
\end{minipage} \\
\midrule\noalign{}
\endhead
\bottomrule\noalign{}
\endlastfoot
1 (Prolif + Invade) & 13,462 & 132,000 & 22.8 mo & 19.2 mo & +3.6 &
+14.88 & \(4.4 \times 10^{-50}\) \\
2 (GrowSup + Immort) & 6,482 & 138,980 & 22.4 mo & 19.4 mo & +2.9 &
+8.28 & \(1.3 \times 10^{-16}\) \\
3 (Death + Immune) & 3,075 & 142,387 & 22.7 mo & 19.5 mo & +3.2 & +8.68
& \(4.1 \times 10^{-18}\) \\
4 (Genome + Metab) & 4,076 & 141,386 & 23.5 mo & 19.4 mo & +4.1 & +11.78
& \(4.7 \times 10^{-32}\) \\
\end{longtable}
}

\normalsize

All four pairs are independently significant at \(p < 10^{-6}\). The
survival advantage associated with pair completion ranges from +2.9
months (Pair 2) to +4.1 months (Pair 4).

Notably, Pair 4 --- the enabler identified by forbidden state analysis
--- shows the largest survival differential (+4.1 months), while Pair 2
--- the barrier --- shows the smallest (+2.9 months). The thermodynamic
difficulty of a pair and its survival impact are inversely related: the
``easier'' pairs carry the stronger survival signal, possibly because
they represent more common and therefore more statistically powered
configurations.

These pooled results represent the net direction across all cancer
types. As established in §3.8, the direction varies by cancer type. The
pair-specific results confirm that the survival signal is distributed
across all four pairs, consistent with the hallmark pair framework
capturing a multi-dimensional biological axis --- not an artifact of any
single gene or pathway.

\subsubsection{3.10 The Cooperative
Effect}\label{the-cooperative-effect}

If pair completions were independent events (a ``non-interacting''
Boltzmann model), the transition ratio \(P(k{+}1)/P(k)\) would be a
constant equal to \(e^{-\beta} = 0.196\). The observed ratios deviate
systematically:

{\def\LTcaptype{none} % do not increment counter
\begin{longtable}[]{@{}llll@{}}
\toprule\noalign{}
Transition & Observed Ratio & Predicted (Boltzmann) & Deviation \\
\midrule\noalign{}
\endhead
\bottomrule\noalign{}
\endlastfoot
0 → 1 pair & 0.131 & 0.196 & −33.1\% \\
1 → 2 pairs & 0.210 & 0.196 & +7.1\% \\
2 → 3 pairs & 0.183 & 0.196 & −6.5\% \\
3 → 4 pairs & 0.293 & 0.196 & +49.2\% \\
\end{longtable}
}

The first transition (0 → 1) faces a \emph{higher} barrier than the
Boltzmann average, while the final transition (3 → 4) faces a
\emph{lower} barrier. The ratio increases from 0.131 to 0.293 across the
four transitions.

This is a cooperative effect: once a tumor has completed one pair,
subsequent completions become progressively easier. In condensed matter
physics, this pattern is characteristic of nucleation --- an initial
barrier followed by accelerating growth once a critical nucleus forms.
In cancer biology, this suggests that the first pair completion grants
the tumor capabilities (e.g., genomic instability via Pair 4) that
catalyze the acquisition of additional pairs.

The cooperative effect produces detectable deviations from pure
Boltzmann statistics at \(k = 0\) (33\% suppressed) and \(k = 4\) (49\%
enhanced). Tumors reaching all four complete pairs are 49\% more common
than a non-interacting model predicts, consistent with the
all-or-nothing character of advanced disease: tumors that begin to
accumulate pairs may undergo a phase transition-like acceleration of
progression.

\subsubsection{3.11 The Pair Interaction
Matrix}\label{the-pair-interaction-matrix}

The cooperative effect (§3.10) establishes that pair completions
interact, but does not resolve the structure of those interactions. To
extract the pair-pair coupling, we model the system as four binary spins
\(\sigma_i \in \{0,1\}\) (pair \(i\) complete or incomplete) governed
by:

\[H = \sum_i \varepsilon_i \sigma_i + \sum_{i<j} J_{ij} \sigma_i \sigma_j \tag{5}\]

where \(\varepsilon_i\) is the intrinsic field (pair difficulty) and
\(J_{ij}\) is the spin-spin coupling. The coupling \(J_{ij}\) is
estimated as the log-odds ratio of co-occurrence:
\(J_{ij} = \ln\bigl(P(\sigma_i{=}1, \sigma_j{=}1) / P(\sigma_i{=}1) P(\sigma_j{=}1)\bigr)\).

The resulting coupling matrix, computed from all 198,862 tumors:

{\def\LTcaptype{none} % do not increment counter
\begin{longtable}[]{@{}lllll@{}}
\toprule\noalign{}
& Pair 1 & Pair 2 & Pair 3 & Pair 4 \\
\midrule\noalign{}
\endhead
\bottomrule\noalign{}
\endlastfoot
\textbf{Pair 1} (Prolif + Invade) & --- & +1.05 & +1.55 & +1.16 \\
\textbf{Pair 2} (GrowSup + Immort) & +1.05 & --- & +1.52 & +2.29 \\
\textbf{Pair 3} (Death + Immune) & +1.55 & +1.52 & --- & +1.74 \\
\textbf{Pair 4} (Genome + Metab) & +1.16 & +2.29 & +1.74 & --- \\
\end{longtable}
}

All six couplings are positive (ferromagnetic), meaning that completing
any pair increases the probability of completing any other. This
confirms the cooperative effect is universal across pair combinations,
not driven by a single pair interaction.

The strongest coupling is \(J_{24} = +2.29\) --- the interaction between
Pair 2 (the barrier) and Pair 4 (the enabler). This is the signature of
the enabler mechanism: Pair 4 completion most strongly facilitates the
pair that is hardest to achieve alone. Without the enabler, Pair 2 is
nearly inaccessible (only 8 of 27 possible Pair-2-only microstates are
observed --- a 70\% forbidden rate, the highest of any single-pair
macrostate). With Pair 4 present, Pair 2 becomes accessible (7 of 9
Pair-2+4 microstates observed --- only 22\% forbidden).

The weakest coupling is \(J_{12} = +1.05\) --- the interaction between
Pair 1 (neutral) and Pair 2 (barrier). Pair 1 completion provides the
least assistance in overcoming the Pair 2 barrier.

The eigenvalues of \(J\) are
\(\lambda = \{+4.70, -0.69, -1.70, -2.31\}\). The leading eigenvalue is
isolated by a gap of 5.40 from the second eigenvalue, indicating that
72\% of the interaction variance is captured by a single cooperative
mode. The system is not four independent interactions --- it is
dominated by a single collective tendency toward simultaneous pair
completion.

The per-macrostate forbidden rates reveal the full landscape of
thermodynamic accessibility across all 16 pair-states:

{\def\LTcaptype{none} % do not increment counter
\begin{longtable}[]{@{}
  >{\raggedright\arraybackslash}p{(\linewidth - 10\tabcolsep) * \real{0.1667}}
  >{\raggedright\arraybackslash}p{(\linewidth - 10\tabcolsep) * \real{0.1667}}
  >{\raggedright\arraybackslash}p{(\linewidth - 10\tabcolsep) * \real{0.1667}}
  >{\raggedright\arraybackslash}p{(\linewidth - 10\tabcolsep) * \real{0.1667}}
  >{\raggedright\arraybackslash}p{(\linewidth - 10\tabcolsep) * \real{0.1667}}
  >{\raggedright\arraybackslash}p{(\linewidth - 10\tabcolsep) * \real{0.1667}}@{}}
\toprule\noalign{}
\begin{minipage}[b]{\linewidth}\raggedright
Macrostate (P1,P2,P3,P4)
\end{minipage} & \begin{minipage}[b]{\linewidth}\raggedright
k
\end{minipage} & \begin{minipage}[b]{\linewidth}\raggedright
Microstates
\end{minipage} & \begin{minipage}[b]{\linewidth}\raggedright
Observed
\end{minipage} & \begin{minipage}[b]{\linewidth}\raggedright
Forbidden
\end{minipage} & \begin{minipage}[b]{\linewidth}\raggedright
N
\end{minipage} \\
\midrule\noalign{}
\endhead
\bottomrule\noalign{}
\endlastfoot
(0,0,0,0) & 0 & 81 & 59 & 22 (27\%) & 170,622 \\
(1,0,0,0) & 1 & 27 & 19 & 8 (30\%) & 13,906 \\
(0,1,0,0) & 1 & 27 & 8 & 19 (70\%) & 4,511 \\
(0,0,1,0) & 1 & 27 & 16 & 11 (41\%) & 2,086 \\
(0,0,0,1) & 1 & 27 & 21 & 6 (22\%) & 1,909 \\
(1,1,0,0) & 2 & 9 & 3 & 6 (67\%) & 1,228 \\
(0,1,0,1) & 2 & 9 & 7 & 2 (22\%) & 1,500 \\
(1,0,1,0) & 2 & 9 & 6 & 3 (33\%) & 976 \\
(1,0,0,1) & 2 & 9 & 6 & 3 (33\%) & 698 \\
(0,0,1,1) & 2 & 9 & 4 & 5 (56\%) & 142 \\
(0,1,1,0) & 2 & 9 & 3 & 6 (67\%) & 167 \\
(1,1,1,0) & 3 & 3 & 1 & 2 (67\%) & 359 \\
(1,1,0,1) & 3 & 3 & 3 & 0 (0\%) & 301 \\
(1,0,1,1) & 3 & 3 & 2 & 1 (33\%) & 179 \\
(0,1,1,1) & 3 & 3 & 3 & 0 (0\%) & 25 \\
(1,1,1,1) & 4 & 1 & 1 & 0 (0\%) & 253 \\
\end{longtable}
}

The pattern is striking. Macrostates containing Pair 2 \emph{without}
Pair 4 have the highest forbidden rates: (0,1,0,0) at 70\%, (1,1,0,0) at
67\%, (0,1,1,0) at 67\%. But macrostates containing Pair 2 \emph{with}
Pair 4 drop dramatically: (0,1,0,1) at 22\%, (1,1,0,1) at 0\%. Pair 4
rescues Pair 2 from inaccessibility.

This interaction matrix provides the complete empirical specification of
the tumor's selective landscape at the pair level. The diagonal elements
(pair difficulties \(\varepsilon_i\)) describe the isolated cost of each
pair. The off-diagonal elements (\(J_{ij}\)) describe how pairs
interact. Together, they define a 10-parameter effective Hamiltonian
that governs the distribution of 198,862 tumors across \(2^4 = 16\)
macrostates with five degrees of freedom.

\subsubsection{3.12 The IDH1 Bifurcation: A Phase Transition in
Glioma}\label{the-idh1-bifurcation-a-phase-transition-in-glioma}

The Ising model predicts that the enabler coupling \(J_{24} = +2.29\)
should produce its most dramatic effects in cancer types where Pair 4
completion is common. Diffuse glioma provides the critical test: 27.0\%
of glioma patients have Pair 4 complete (vs.~2.5\% globally --- a
10.7-fold enrichment), driven almost entirely by IDH1 mutation.

The survival stratification by Pair 4 status in 4,351 glioma patients
yields the strongest single signal in the entire dataset:

{\def\LTcaptype{none} % do not increment counter
\begin{longtable}[]{@{}lllll@{}}
\toprule\noalign{}
Group & N (deceased) & Median OS (mo) & z & p \\
\midrule\noalign{}
\endhead
\bottomrule\noalign{}
\endlastfoot
Pair 4 complete & 326 & 51.4 & +20.66 & 7.8 × 10⁻⁹⁵ \\
Pair 4 incomplete & 1,546 & 15.0 & --- & --- \\
\textbf{Δ} & --- & \textbf{+36.4} & --- & --- \\
\end{longtable}
}

Patients with Pair 4 complete (IDH-mutant) live a median \textbf{36.4
months longer} than those without. This is the IDH-mutant favorable
subclass, recovered purely from the pair framework without using any
molecular subtype labels.

The mechanism is visible in the macrostate distribution. The \(J_{24}\)
coupling drags Pair 2 along: 25.0\% of glioma patients have Pair 2
complete (vs.~4.2\% globally --- a 6.0-fold enrichment). Of the 1,173
glioma patients with Pair 4 complete, 981 (83.6\%) also have Pair 2
complete. The dominant glioma macrostate is (0,1,0,1) --- Pair 2 + Pair
4 complete, all others incomplete --- containing 937 patients (21.5\% of
all glioma). This is the IDH-mutant astrocytoma (ATRX + TP53 + IDH1),
expressed as a single thermodynamic ground state.

The glioma pair-completeness distribution is \textbf{bimodal}:

{\def\LTcaptype{none} % do not increment counter
\begin{longtable}[]{@{}llll@{}}
\toprule\noalign{}
k & N & \% & Phase \\
\midrule\noalign{}
\endhead
\bottomrule\noalign{}
\endlastfoot
0 & 3,016 & 69.3\% & Paramagnetic (IDH-wildtype) \\
1 & 315 & 7.2\% & Boundary (depleted) \\
2 & 967 & 22.2\% & Ferromagnetic (IDH-mutant) \\
3--4 & 53 & 1.2\% & Saturated \\
\end{longtable}
}

The depletion at \(k = 1\) (7.2\%, vs.~the monotonic decay predicted by
a single Boltzmann distribution) is the hallmark of a first-order phase
transition: two coexisting phases with a depleted boundary. Glioma
tumors are either paramagnetic (k = 0, IDH-wildtype GBM) or
ferromagnetic (k = 2, IDH-mutant astrocytoma), with few tumors at the
transition point. The Ising model identifies this as a nucleation event:
IDH1 mutation simultaneously completes Pair 4 and, via the dominant
\(J_{24}\) coupling, drags Pair 2 across the accessibility threshold.
The tumor jumps from \(k = 0\) to \(k = 2\) in a single molecular event.

In breast cancer, the same Pair 4 completion produces the
\textbf{opposite} survival effect (Δ = −18.1 months, z = −2.20, p =
0.028). The enabler is protective in glioma and destructive in breast
--- confirming that survival direction is determined not by temperature
alone, but by which pairs are completed and the tissue-specific biology
of those completions.

\subsubsection{3.13 Per-Cancer Effective
Temperature}\label{per-cancer-effective-temperature}

The global Boltzmann fit yields \(\beta = 1.63\) across 198,862 tumors.
Fitting \(\beta\) independently for each cancer type with \$\geq\$500
patients reveals a three-fold range of effective temperatures:

{\def\LTcaptype{none} % do not increment counter
\begin{longtable}[]{@{}lllll@{}}
\toprule\noalign{}
Cancer Type & N & β & T = 1/β & k = 0 (\%) \\
\midrule\noalign{}
\endhead
\bottomrule\noalign{}
\endlastfoot
Clear Cell Renal & 740 & 2.94 & 0.34 & 96.1\% \\
Soft Tissue Sarcoma & 5,474 & 2.91 & 0.34 & 90.3\% \\
Pancreatic & 2,862 & 2.69 & 0.37 & 96.5\% \\
Hepatocellular & 2,906 & 2.61 & 0.38 & 90.5\% \\
Cholangiocarcinoma & 540 & 2.55 & 0.39 & 90.6\% \\
Esophageal & 661 & 2.50 & 0.40 & 89.7\% \\
High-Grade Serous Ovarian & 1,372 & 2.32 & 0.43 & 97.0\% \\
Lung Adenocarcinoma & 5,272 & 2.18 & 0.46 & 83.5\% \\
Breast (targeted) & 2,474 & 2.17 & 0.46 & 82.2\% \\
Lung Squamous & 1,453 & 2.12 & 0.47 & 85.8\% \\
Bladder & 3,007 & 2.01 & 0.50 & 78.8\% \\
Breast (BRCA) & 9,763 & 1.96 & 0.51 & 85.1\% \\
Prostate (PRAD) & 5,619 & 1.88 & 0.53 & 95.2\% \\
Melanoma & 2,627 & 1.53 & 0.66 & 57.0\% \\
Colorectal & 4,596 & 1.42 & 0.70 & 80.7\% \\
\textbf{Diffuse Glioma} & \textbf{4,351} & \textbf{1.32} & \textbf{0.76}
& \textbf{69.3\%} \\
Gastric (STAD) & 1,612 & 1.25 & 0.80 & 75.1\% \\
Uterine Endometrial & 3,875 & 1.04 & 0.96 & 66.8\% \\
\end{longtable}
}

The coldest cancer (clear cell renal, \(\beta = 2.94\)) has 2.8× higher
selective barriers than the hottest (uterine endometrial,
\(\beta = 1.04\)). Glioma (\(\beta = 1.32\)) is among the hottest ---
not because glioma tumors are more aggressive, but because the IDH1
pathway provides a low-energy route to pair completion that is
unavailable in colder cancers. The high temperature reflects the
thermodynamic accessibility of the (0,1,0,1) ground state, not the
clinical aggressiveness of the disease.

The correlation between per-cancer \(\beta\) and survival direction is
weak (Pearson \(r = -0.13\), \(p = 0.49\), \(n = 30\) cancer types),
indicating that temperature alone does not determine whether balance is
protective or destructive. The survival direction depends on which
specific pairs are completed and the tissue-specific biology of those
completions --- the full coupling structure \(J_{ij}\), not merely the
scalar temperature.

\subsubsection{3.14 The Periodic Table of Oncology --- Per-Cancer
Coupling
Matrices}\label{the-periodic-table-of-oncology-per-cancer-coupling-matrices}

Having measured the global coupling matrix (§3.11) and per-cancer
effective temperatures (§3.13), we now compute the full \(4 \times 4\)
coupling matrix \(J_{ij}\) independently for each of 37 cancer types
with \(\geq 200\) patients, using Laplace-smoothed log-odds ratios. This
yields the complete Ising parameterization per cancer type: four
marginal pair probabilities, six couplings, the eigenvalue spectrum, and
the three candidate universal constants. The key findings follow.

\textbf{The three constants are emergent, not intrinsic.} The products
\(\beta \cdot \alpha_c^2\), \(\lambda_1/\pi\), and
\((1 - f_\mathrm{forb}) \cdot \beta\) hold to high precision for the
global 198,862-patient mixture (deviations of 0.08\%, 5.6\%, and 3.2\%,
respectively), but no individual cancer type satisfies all three
simultaneously. Per-cancer \(\beta \cdot \alpha_c^2\) ranges from 1.30
(BCC) to 4.51 (clear cell renal); \(\lambda_1/\pi\) ranges from 0.24
(BCC) to 3.02 (HGSOC); \((1 - f_\mathrm{forb}) \cdot \beta\) ranges from
0.17 (uterine endometrioid) to 0.97 (mixed). The constants characterize
the \emph{mixture}, not the individual components. In statistical
mechanical terms, each cancer type is a thermodynamic subsystem with its
own temperature and coupling structure; the apparent constants emerge
when these subsystems are observed in aggregate --- analogous to how the
ideal gas law holds for a gas mixture even when individual molecular
species have different interaction potentials.

\textbf{Coupling matrices are cancer-specific fingerprints.} The
coupling matrices are not perturbations around a universal shape ---
they are qualitatively different by cancer type:

{\def\LTcaptype{none} % do not increment counter
\begin{longtable}[]{@{}
  >{\raggedright\arraybackslash}p{(\linewidth - 12\tabcolsep) * \real{0.1429}}
  >{\raggedright\arraybackslash}p{(\linewidth - 12\tabcolsep) * \real{0.1429}}
  >{\raggedright\arraybackslash}p{(\linewidth - 12\tabcolsep) * \real{0.1429}}
  >{\raggedright\arraybackslash}p{(\linewidth - 12\tabcolsep) * \real{0.1429}}
  >{\raggedright\arraybackslash}p{(\linewidth - 12\tabcolsep) * \real{0.1429}}
  >{\raggedright\arraybackslash}p{(\linewidth - 12\tabcolsep) * \real{0.1429}}
  >{\raggedright\arraybackslash}p{(\linewidth - 12\tabcolsep) * \real{0.1429}}@{}}
\toprule\noalign{}
\begin{minipage}[b]{\linewidth}\raggedright
Cancer
\end{minipage} & \begin{minipage}[b]{\linewidth}\raggedright
N
\end{minipage} & \begin{minipage}[b]{\linewidth}\raggedright
Strongest Coupling
\end{minipage} & \begin{minipage}[b]{\linewidth}\raggedright
J value
\end{minipage} & \begin{minipage}[b]{\linewidth}\raggedright
Weakest Coupling
\end{minipage} & \begin{minipage}[b]{\linewidth}\raggedright
J value
\end{minipage} & \begin{minipage}[b]{\linewidth}\raggedright
Phase
\end{minipage} \\
\midrule\noalign{}
\endhead
\bottomrule\noalign{}
\endlastfoot
Breast (BRCA) & 9,763 & J(2,4) & +2.78 & P1--P2 & +0.96 & Critical \\
Prostate (PRAD) & 5,619 & J(3,4) & +2.76 & P1--P2 & +2.00 & Critical \\
Soft Tissue Sarcoma & 5,474 & J(1,3) & +3.33 & P1--P2 & −0.14 & Mixed \\
Lung Adenocarcinoma & 5,272 & J(3,4) & +1.05 & P2--P3 & +0.48 &
Critical \\
Colorectal & 4,596 & J(3,4) & +2.06 & P1--P2 & +1.18 & Disordered \\
Diffuse Glioma & 4,351 & J(1,3) & +3.12 & P3--P4 & +0.58 & Disordered \\
Uterine Endometrial & 3,875 & J(2,4) & +1.95 & P1--P3 & +0.68 &
Disordered \\
Bladder & 3,007 & J(2,3) & +1.38 & P3--P4 & −0.17 & Mixed \\
\end{longtable}
}

In breast cancer, the enabler--barrier coupling J(2,4) dominates,
consistent with the global picture. In glioma, the strongest coupling is
J(1,3) = +3.12, linking Pair 1 (Proliferation + Invasion) to Pair 3
(Death + Immune) --- a fundamentally different topology. Soft tissue
sarcoma shows an \emph{antiferromagnetic} coupling between Pair 1 and
Pair 2 (J = −0.14), meaning proliferative-invasive and
growth-suppression-immortality signatures are mutually exclusive in
sarcoma tumors.

\textbf{Per-pair survival signatures.} Instead of a single z-score
(balance vs lopsided), we compute a 4-element survival vector
\((z_{P1}, z_{P2}, z_{P3}, z_{P4})\) for each cancer type --- the
z-score for survival when each pair is individually complete vs
incomplete:

{\def\LTcaptype{none} % do not increment counter
\begin{longtable}[]{@{}llllll@{}}
\toprule\noalign{}
Cancer & z(P1) & z(P2) & z(P3) & z(P4) & Pattern \\
\midrule\noalign{}
\endhead
\bottomrule\noalign{}
\endlastfoot
Diffuse Glioma & +4.3 & +17.1 & +2.3 & +20.7 & All pairs → longer \\
Breast (BRCA) & −2.7 & −1.8 & −1.0 & −2.2 & All pairs → shorter \\
Bladder & +3.0 & +3.0 & −1.3 & +2.2 & P3 reversal \\
Uterine Endometrial & −3.2 & +0.5 & −1.7 & −0.2 & P1 dominant \\
Melanoma (SKCM) & +1.5 & −0.8 & +1.1 & +3.1 & P4 dominant \\
Soft Tissue & +1.4 & +2.7 & −0.9 & +2.4 & P3 reversal \\
HGSOC & +0.6 & +3.5 & --- & --- & P2 dominant \\
\end{longtable}
}

Glioma and breast cancer occupy opposite poles: in glioma, every pair
completion predicts longer survival; in breast, every pair completion
predicts shorter survival. Several cancers (bladder, soft tissue) show
\emph{mixed} signatures where most pairs predict longer survival but
Pair 3 (Death + Immune) reverses direction. This 4-vector fingerprint
provides more information than the scalar balance z-score and may
resolve cases where the overall balance signal is weak.

\textbf{No coupling feature predicts survival direction.} We tested 17
features derived from the per-cancer coupling matrix (including all six
\(J_{ij}\), eigenvalues, pair marginals, forbidden fraction, bimodality
index) for correlation with the balance survival z-score across 34
cancer types. The strongest Pearson correlation is P(Pair 4)
(\(r = +0.77\), \(p = 7.3 \times 10^{-8}\)), but this result is driven
entirely by glioma. Excluding glioma: P(Pair 4) drops to \(r = -0.11\),
\(p = 0.54\). All 17 features become non-significant (\(p > 0.27\)) when
glioma is excluded. The survival direction is not predictable from any
single feature of the coupling matrix.

\textbf{Phase classification.} Based on the coupling matrix topology and
effective temperature, we classify each cancer into thermodynamic
phases:

\begin{itemize}
\tightlist
\item
  \textbf{Disordered} (\(\beta < 1.5\), all \(J_{ij} > 0\)): difg, ucec,
  coadread, stad, egc, coad, mel, bcc, rcc, uec --- high temperature,
  all couplings ferromagnetic, most tumors at \(k = 0\).
\item
  \textbf{Critical} (\(1.5 \leq \beta \leq 2.5\), all \(J_{ij} > 0\)):
  brca, luad, skcm, breast, bowel, prostate, bladder, nsclc, mbn, brain,
  stomach, cesc, hgsoc --- intermediate temperature, strong cooperative
  mode.
\item
  \textbf{Ordered} (\(\beta > 2.5\), all \(J_{ij} > 0\)): pancreas,
  paad, ccrcc, chol --- low temperature, very high \(k = 0\) fraction
  (\(>90\%\)), pair completion is thermodynamically expensive.
\item
  \textbf{Frustrated} (some \(J_{ij} < -1\)): lusc --- competing
  ferromagnetic and antiferromagnetic interactions.
\item
  \textbf{Mixed} (some \(J_{ij}\) slightly negative): soft tissue, blca,
  hcc, hnsc, esca, lung, apad --- local antiferromagnetic couplings
  reflecting mutually exclusive hallmark pathways.
\end{itemize}

Glioma is the only cancer with unambiguous first-order phase transition
behavior (bimodality index 0.18, where values below 0.5 indicate k=1
depletion). Prostate shows a secondary candidate (index 0.45). All other
cancers have unimodal k-distributions consistent with thermal
equilibrium in a single phase.

\textbf{The leading eigenvector is nearly uniform.} Across 37 cancer
types, the leading eigenvector \(\mathbf{v}_1\) of the coupling matrix
is approximately \(\frac{1}{2}(1, 1, 1, 1)\) --- all four components
within 0.3--0.65 for most types. The system behaves as a single
collective mode where all pairs activate together. The notable exception
is glioma, where \(\mathbf{v}_1 = (0.63, 0.32, 0.63, 0.31)\): Pairs 1
and 3 (Prolif/Invade and Death/Immune) carry double the weight of Pairs
2 and 4 in the dominant mode, reflecting the IDH1-driven pathway that
simultaneously activates metabolism and growth suppression while leaving
proliferation and immune evasion relatively uncoupled.

\subsubsection{3.15 Exhaustive Hypothesis Testing --- Burning the
Seams}\label{exhaustive-hypothesis-testing-burning-the-seams}

The open question from §3.14 --- what, if anything, predicts survival
direction across cancer types? --- demands exhaustive testing. We
computed 37 candidate features for each of the 34 cancer types with ≥30
deaths and tested each for Pearson and Spearman correlation with the
balance survival z-score under three nested robustness conditions: (1)
full dataset, (2) excluding glioma, and (3) excluding glioma and breast
cancer.

\textbf{Features tested.} The 37 features fall into six categories:

\begin{itemize}
\tightlist
\item
  \emph{Frustration indices} (2): the Pearson anti-correlation
  \(F = -\mathrm{corr}(\mathbf{P}, \mathbf{R})\) between pair marginals
  and coupling row sums; the cosine-based frustration
  \(1 - \cos(\hat{\mathbf{P}}, \hat{\mathbf{R}})\).
\item
  \emph{Marginal probabilities} (4): \(P(\mathrm{Pair}\ i)\) for
  \(i = 1, \ldots, 4\).
\item
  \emph{Eigenvalue features} (6): \(\lambda_1\), \(\lambda_2\), spectral
  gap, \(\lambda_1/|\lambda_2|\), variance fraction, participation
  ratio.
\item
  \emph{Coupling topology} (7): \(J_{13}\), \(J_{24}\), \(J_{34}\),
  \(\bar{J}\), \(J_\mathrm{max}\), \(J_\mathrm{range}\),
  \(J_{13} - J_{24}\) dominance, count of negative couplings.
\item
  \emph{Compound features} (12): \(P_4 \cdot \beta\), \(P_4 / \beta\),
  effective field \(\mathbf{P} \cdot \mathbf{R}\), \(\det(J)\),
  \(\mathrm{tr}(J^2)\), population-eigenvector alignment
  \(\cos(\hat{\mathbf{P}}, \hat{\mathbf{v}}_1)\),
  \(P_\mathrm{max}/P_\mathrm{min}\), pair entropy, ground state energy,
  Pair 4 share, rank inversion, enabler/barrier ratio.
\item
  \emph{PCA modes} (4): projections of the 6-dimensional coupling vector
  onto the first four principal components.
\end{itemize}

\textbf{Results: what burned.}

The frustration index \(F = -\mathrm{corr}(\mathbf{P}, \mathbf{R})\) ---
which would predict that cancers with misaligned marginals and couplings
have stronger survival signals --- showed no significant correlation
with any outcome: \(r = +0.18\) with \(|z_\mathrm{balance}|\)
(\(p = 0.24\)), \(r = +0.13\) with bimodality (\(p = 0.46\)),
\(r = -0.07\) with \(\lambda_2\) (\(p = 0.70\)). All three predictions
failed. Excluding glioma made all correlations weaker. The frustration
index is dead.

The marginal P(Pair 4) showed \(r = +0.77\) (\(p = 8.6 \times 10^{-8}\))
with \(z_\mathrm{balance}\) in the full dataset --- an apparently
powerful result. However, the Spearman rank correlation was
\(\rho = +0.06\) (\(p = 0.73\)), revealing that the Pearson value is
driven entirely by the glioma outlier (P(Pair 4) = 27\% vs.~median
1.4\%). Excluding glioma: \(r = -0.11\), \(p = 0.54\). Binary threshold
analysis at every cutoff from 1\% to 10\% showed no significant split
without glioma (all \(p > 0.29\)). P(Pair 4) as a continuous phase
coordinate is dead.

Temperature (\(\beta\)), all six coupling constants (\(J_{ij}\)), all
eigenvalue features (\(\lambda_1\), \(\lambda_2\), spectral gap),
determinant, trace of \(J^2\), ground state energy, PCA modes PC1--PC4,
and all compound features (torque, rank inversion, enabler/barrier
ratio, effective field) --- all showed \(p > 0.10\) in every robustness
tier. Twenty-seven features are dead across all three conditions.

\textbf{Results: what survived.}

One feature passed Pearson significance (\(p < 0.05\)) in all three
robustness conditions:

\[\cos(\hat{\mathbf{P}}, \hat{\mathbf{v}}_1) \quad \text{vs.} \quad z_\mathrm{balance}: \quad r = -0.38\ (p = 0.027), \quad -0.37\ (p = 0.032), \quad -0.44\ (p = 0.012) \tag{6}\]

where \(\hat{\mathbf{P}} = \mathbf{P}/\|\mathbf{P}\|\) is the unit
population-marginal vector and \(\hat{\mathbf{v}}_1\) is the unit
dominant eigenvector of the coupling matrix \(J\). The effect is
negative: greater alignment of the population state with the coupling
ground state predicts \emph{worse} survival from balance.

\textbf{Physical interpretation.} The eigenvector \(\mathbf{v}_1\)
defines the ``easy axis'' of the coupling landscape --- the direction of
minimum energy in the Ising Hamiltonian. When the population's
pair-probability vector \(\mathbf{P}\) is aligned with \(\mathbf{v}_1\),
the population sits in a deep potential well. Perturbation toward
balance (\(k = 2\)) costs energy: it moves the system uphill, and
shorter survival follows. When \(\mathbf{P}\) is misaligned (population
displaced from the minimum), balance can \emph{lower} energy by moving
toward a more coupled configuration, and longer survival follows.

In concrete terms: soft tissue sarcoma (\(\cos\theta = 0.32\),
\(z = +3.07\)) and breast carcinoma (\(\cos\theta = 0.42\),
\(z = +3.64\)) are maximally misaligned and show the strongest positive
balance effect. Clear cell renal (\(\cos\theta = 0.94\), \(z = -0.96\))
and prostate (\(\cos\theta = 0.90\), \(z = -1.30\)) are maximally
aligned with their coupling ground state, and balance hurts.

\textbf{Robustness characterization:}

\begin{itemize}
\tightlist
\item
  \emph{Leave-one-out:} all 34/34 iterations yield \(p < 0.05\) (\(r\)
  range: \(-0.35\) to \(-0.41\), \(\sigma = 0.012\)). No single cancer
  type drives the effect.
\item
  \emph{Partial correlation:} controlling for \(\log(N)\) and \(\beta\)
  jointly: \(r = -0.39\) (\(p = 0.024\)). The effect is not confounded
  by sample size or temperature.
\item
  \emph{Large-sample cancers only} (\(N > 1{,}000\), \(n = 23\)):
  \(r = -0.46\) (\(p = 0.028\)), Spearman \(\rho = -0.45\)
  (\(p = 0.030\)). Both Pearson and Spearman agree in well-powered
  types.
\item
  \emph{Bootstrap 95\% CI:} \([-0.59, -0.10]\) (full);
  \([-0.63, +0.01]\) (excluding glioma). The full-dataset CI excludes
  zero; the glioma-excluded CI barely includes it.
\end{itemize}

\textbf{Limitations.} The effect explains \(R^2 = 0.14\) to \(0.19\) of
variance (14--19\%). With 37 features tested, it does not survive
Bonferroni correction (threshold \(p < 0.0014\); observed \(p = 0.027\))
or Benjamini-Hochberg FDR at \(q = 0.05\). Spearman correlation is
marginal in the full dataset (\(\rho = -0.30\), \(p = 0.083\)) though it
reaches significance in \(N > 1{,}000\) types (\(\rho = -0.45\),
\(p = 0.030\)). The population-eigenvector alignment is the best
available scalar predictor but is not definitive.

\textbf{The pair entropy confound.} The cosine alignment
\(\cos(\hat{\mathbf{P}}, \hat{\mathbf{v}}_1)\) correlates at
\(r = 0.95\) with the Shannon entropy of the pair-probability vector.
This means the alignment effect is largely equivalent to: cancer types
with low pair-entropy (one pair dominates) tend to benefit from balance;
cancer types with high pair-entropy (many pairs active) tend not to
benefit. The two features are mathematically near-collinear and cannot
be distinguished at \(n = 34\).

\textbf{Conclusion.} From 37 features, 36 are dead. One is wounded but
standing. The data's clearest message is negative: no single scalar
parameter governs survival direction across 37 cancer types. The
survival direction is an emergent property of the full \(4 \times 4\)
interaction between the population state \(\mathbf{P}\) and the coupling
landscape \(J\) --- it cannot be compressed to a scalar without losing
81--86\% of the information. The critical point described by the three
universal constants (\(\beta \cdot \alpha_c^2 = 5/2\),
\(\lambda_1/\pi = 3/2\), \((1 - f_\mathrm{forb}) \cdot \beta = 1\)) is
the centroid of a 37-point cloud in coupling space, not a single state
of matter.

\subsubsection{3.16 Study-Level Decomposition --- Unblending the
Mixture}\label{study-level-decomposition-unblending-the-mixture}

The preceding analyses (§3.13--3.15) aggregate all patients of a given
cancer type into a single data point, yielding \(n = 34\) cancer types.
This averaging may mask within-tissue variation: different studies of
the same cancer type (e.g., 11 lung adenocarcinoma cohorts) sample
different patient populations, sequencing eras, and treatment regimes.
If the control signal hides in within-tissue variance, aggregation would
destroy it.

We decomposed the analysis to the study level: for each of the 155
studies with ≥10 deaths and a single dominant cancer type, we computed a
study-level pair-probability vector \(\mathbf{P}_\mathrm{study}\) from
that study's patients, then measured the cosine alignment
\(\cos(\hat{\mathbf{P}}_\mathrm{study}, \hat{\mathbf{v}}_1)\) against
the cancer-type's dominant coupling eigenvector \(\hat{\mathbf{v}}_1\)
(from the per-cancer \(J\) matrix, §3.14). This decomposition treats
\(J\) as a tissue property (coupling landscape) and \(\mathbf{P}\) as a
cohort property (where this particular study's population sits in that
landscape).

\textbf{Result: the signal collapses.} Study-level
\(\cos(\hat{\mathbf{P}}, \hat{\mathbf{v}}_1)\) vs.~study-level
\(z_\mathrm{balance}\) yields \(r = -0.096\) (\(p = 0.23\)) across 155
studies --- not significant. The within-tissue partial correlation,
controlling for cancer type via fixed effects (\(n = 137\) studies from
types with \(\geq 3\) cohorts), is \(r = -0.034\) (\(p = 0.70\)). No
individual cancer type with \(\geq 3\) studies shows a robust
correlation. The lone survivor from §3.15 (\(r = -0.38\), \(p = 0.027\)
at the cancer-type level) does not survive decomposition to the study
level.

Study-level pair entropy (\(r = +0.053\), \(p = 0.52\)) and study-level
mean broken pair count (\(r = -0.124\), \(p = 0.13\)) are similarly
non-significant. The wounded survivor from §3.15 is now dead at the
study level. The population-eigenvector alignment captured a
between-tissue signal that does not replicate within tissues.

\subsubsection{3.17 The Survival Hessian --- Second-Order Interactions
and Model-Free
Curvature}\label{the-survival-hessian-second-order-interactions-and-model-free-curvature}

The pair-level survival analysis (§3.9) treats each pair independently:
does completing Pair \(i\) predict survival? But pairs may interact: the
survival effect of completing Pair \(i\) may depend on whether Pair
\(j\) is also complete. We measured this interaction by computing the
synergy:

\[S_{ij} = z_\mathrm{both} - z_i - z_j \tag{7}\]

where \(z_i\) is the survival \(z\)-score for patients with Pair \(i\)
complete vs.~not (globally), and \(z_\mathrm{both}\) is the \(z\)-score
for patients with both Pair \(i\) and Pair \(j\) complete vs.~neither
complete.

\textbf{Global second-order tensor.} The individual pair \(z\)-scores
are \(z_1 = +5.26\), \(z_2 = +6.52\), \(z_3 = +0.25\), \(z_4 = +11.13\).
All six pairwise synergies are negative:

{\def\LTcaptype{none} % do not increment counter
\begin{longtable}[]{@{}lrrrr@{}}
\toprule\noalign{}
& P1 & P2 & P3 & P4 \\
\midrule\noalign{}
\endhead
\bottomrule\noalign{}
\endlastfoot
\textbf{P1} & --- & \(-10.1\) & \(-4.5\) & \(-14.0\) \\
\textbf{P2} & --- & --- & \(-6.0\) & \(-1.2\) \\
\textbf{P3} & --- & --- & --- & \(-9.2\) \\
\textbf{P4} & --- & --- & --- & --- \\
\end{longtable}
}

Every entry is negative. The pair interactions are universally
\emph{antagonistic}: completing two pairs simultaneously produces a
survival benefit that is weaker than the sum of completing each alone.
This is consistent with diminishing returns --- the first pair
completion captures the largest survival gain, and subsequent
completions add progressively less, because the most informative
biological capability has already been acquired.

The strongest antagonism (\(S_{14} = -14.0\)) occurs between the neutral
pair (Prolif + Invade) and the enabler (Genome + Metab), suggesting that
these pairs provide overlapping rather than complementary survival
information. The weakest antagonism (\(S_{24} = -1.2\)) links the
barrier to the enabler --- the biologically most distant pair, whose
joint completion provides the most independent information.

\textbf{Per-cancer synergy correlates with survival direction.} Among
cancer types with sufficient data to compute the synergy tensor (11
types for \(S_{14}\), 9 for \(S_{24}\), 6 for \(S_{34}\)), the synergy
magnitude correlates near-perfectly with the balance survival
\(z\)-score: \(r = -0.996\) (\(S_{14}\)), \(r = -0.979\) (\(S_{24}\)),
\(r = -0.983\) (\(S_{34}\)). These correlations are striking but should
be interpreted cautiously: synergy \(S_{ij}\) and \(z_\mathrm{balance}\)
are both derived from survival comparisons in the same data, so partial
mathematical coupling is expected.

\textbf{From crude synergy to model-free curvature.} The crude synergy
tensor suggested near-perfect prediction of survival direction
(\(r = -0.996\)), but was computed from synergy differences in only
6--11 cancer types with no standard errors. We replaced it with the
survival Hessian: the matrix of second partial derivatives
\(H_{ij} = \partial^2 S / \partial \sigma_i \partial \sigma_j\),
estimated via OLS regression on individual patient survival with
interaction terms. This model-free approach does not assume a
Hamiltonian --- it measures curvature directly from the data with proper
standard errors.

\textbf{Global Hessian.} For 60,988 deceased patients, we fitted the
full 16-parameter Taylor expansion through fourth order:
\(S(\sigma) = a + \sum g_i \sigma_i + \sum H_{ij} \sigma_i \sigma_j + \sum T_{ijk} \sigma_i \sigma_j \sigma_k + K_{1234} \sigma_1 \sigma_2 \sigma_3 \sigma_4\).
The complete expansion is identifiable because the 16 pair states span
\(2^4\) degrees of freedom. Results by order:

\emph{Gradient (first order).} Pair 1 (Prolif + Invade) completion adds
\(+4.94\) months (\(t = +11.4\), \(p < 10^{-29}\)). Pair 4 (Genome +
Metab) adds \(+2.90\) months (\(t = +2.7\), \(p = 0.007\)). Pairs 2 and
3 are not individually significant (\(|t| < 1\)).

\emph{Hessian (second order).} The dominant element is
\(H_{24} = +26.1\) months (\(t = +14.2\), \(p < 10^{-45}\)): patients
who complete both the barrier (Pair 2) and the enabler (Pair 4)
simultaneously gain 26 months beyond the sum of individual pair effects.
The 384 patients with pair state \((0,1,0,1)\) --- Pair 2 + Pair 4 only
--- have mean OS 49.5 months vs.~21.0 months for the ground state
\((0,0,0,0)\). This is the IDH-mutant astrocytoma signal (§3.12)
expressed in the curvature language. The remaining significant elements
are \(H_{14} = -5.3\) months (\(t = -2.4\)) and \(H_{34} = -8.3\) months
(\(t = -2.1\)) --- both antagonistic.

\emph{Third-order tensor.} One element is significant:
\(T_{124} = -11.3\) months (\(t = -2.6\), \(p = 0.009\)), indicating a
three-body interaction among Pairs 1, 2, and 4 that is not captured by
pairwise terms alone. The Ising model's pairwise assumption is violated
at third order.

\emph{Fourth-order coupling.} \(K_{1234} = +8.5\) months (\(t = +0.5\),
\(p = 0.62\), not significant). The four-body interaction does not
contribute detectably.

The Hessian eigenvalues are \(\{-26.8, -3.1, +0.5, +29.5\}\) --- a
saddle structure with one strongly positive and one strongly negative
principal curvature. The energy landscape is neither uniformly convex
nor concave; it has a ridge.

\textbf{Per-cancer Hessians.} For 30 cancer types with \(\geq 50\)
deceased patients, we fitted the second-order model
(\(S \sim \sigma_i + \sigma_i \sigma_j\), 11 parameters) and
eigendecomposed each Hessian.

\textbf{Test: does the Hessian predict survival direction?} The total
interaction \(TI = \sum_{i < j} H_{ij}\) across 27 cancer types with
nonzero \(z_\mathrm{balance}\) yields \(r = -0.15\) (\(p = 0.45\)). The
crude synergy result of \(r = -0.996\) does not survive proper
estimation across the full set of cancer types. The apparent
near-perfect prediction was a small-sample artifact (\(n = 6\)--\(11\))
amplified by mathematical coupling: both the synergy and
\(z_\mathrm{balance}\) are derived from survival stratification by pair
state. With proper regression-based Hessian estimation and \(n = 27\),
the total-interaction signal vanishes.

However, the Hessian reveals two real signals:

\begin{enumerate}
\def\labelenumi{\arabic{enumi}.}
\item
  The specific element \(H_{34}\) (Death/Immune \(\times\) Genome/Metab
  interaction) correlates with survival direction at \(r = -0.52\)
  (\(p = 0.005\)). Cancer types where this interaction is more
  antagonistic tend to benefit from balance. This is the only Hessian
  element that individually predicts direction.
\item
  The gradient magnitude \(\|\nabla S\|\) --- how strongly individual
  pair completions affect survival within a cancer type --- correlates
  with \(z_\mathrm{balance}\) at \(r = +0.49\) (\(p = 0.009\)). Cancer
  types where pair completions matter more individually also show
  stronger balance effects.
\item
  The mean ratio of survival Hessian to thermodynamic coupling
  \(\langle H_{ij} / J_{ij} \rangle\) correlates with survival direction
  at \(r = -0.40\) (\(p = 0.041\)). The survival-to-coupling conversion
  factor --- how efficiently the Ising structure translates to clinical
  outcome --- varies systematically across cancer types.
\end{enumerate}

\textbf{Patient-level composite predictor.} For 19,646 deceased patients
in cancer types with computed Hessians, we fitted
\(\mathrm{OS} \sim n_\mathrm{walls} + d_\mathrm{forb} + \sigma^T H \sigma\).
The domain wall count is the dominant predictor: \(\beta = -4.94\)
months per wall (\(t = -15.1\), \(p < 10^{-50}\)). Forbidden boundary
distance adds independently: \(\beta = +5.81\) months per unit
(\(t = +8.6\), \(p < 10^{-17}\)). Hessian curvature exposure
(\(\sigma^T H \sigma\)) is not significant as a patient-level feature
(\(t = +1.0\), \(p = 0.30\)). The composite model explains
\(R^2 = 1.2\%\) of survival variance --- small in absolute terms but
highly significant, and consistent with the expectation that state-level
features capture only a fraction of the variance in a quantity dominated
by stage, treatment, and tumor biology.

\textbf{Verdict.} The Hessian approach burned the crude synergy tensor's
\(r = -0.996\) --- that signal was noise, not the organizing principle.
What survives: the barrier-enabler synergy \(H_{24} = +26.1\) months is
a real global phenomenon driven by the IDH-mutant glioma axis; the
\(H_{34}\) element specifically predicts cross-cancer survival
direction; domain wall count and forbidden boundary distance are
significant patient-level survival features; and the Ising model's
pairwise assumption is violated at third order (\(T_{124}\),
\(p = 0.009\)). The survival landscape has genuine curvature, but that
curvature does not collapse to a single scalar predictor of direction.
The landscape is a saddle --- and the full \(4 \times 4\)
eigenstructure, not any scalar summary, is required to navigate it.

\subsubsection{3.18 Simpson's Paradox --- Direction Flips at the
Aggregation
Boundary}\label{simpsons-paradox-direction-flips-at-the-aggregation-boundary}

The study-level decomposition (§3.16) permits a direct test for
Simpson's paradox: does the survival direction \emph{reverse} between
the study level and the cancer-type level? For each cancer type with
\(\geq 2\) studies, we compared the cancer-type \(z_\mathrm{balance}\)
(all patients aggregated) with the mean of study-level \(z\)-scores.

Of 28 cancer types tested, \textbf{12 show direction flips (43\%)}.
Notable examples:

{\def\LTcaptype{none} % do not increment counter
\begin{longtable}[]{@{}lrrc@{}}
\toprule\noalign{}
Cancer type & Cancer-level \(z\) & Mean study-level \(z\) & Flip? \\
\midrule\noalign{}
\endhead
\bottomrule\noalign{}
\endlastfoot
Breast carcinoma & \(-2.31\) & \(+0.71\) & YES \\
Lung adenocarcinoma & \(+1.59\) & \(-0.28\) & YES \\
High-grade serous ovarian & \(+3.37\) & \(-1.54\) & YES \\
Soft tissue sarcoma & \(+3.07\) & \(-0.57\) & YES \\
Esophageal & \(-2.40\) & \(+0.89\) & YES \\
Uterine endometrial & \(-3.15\) & \(+1.35\) & YES \\
\end{longtable}
}

In breast cancer, aggregating all patients shows balance is harmful
(\(z = -2.31\)), but the average \emph{individual study} shows balance
is beneficial (\(z = +0.71\)). In soft tissue sarcoma, the aggregated
signal is strongly positive (\(z = +3.07\)), but the average study-level
signal is negative (\(z = -0.57\)).

A 43\% direction-flip rate means that for nearly half of cancer types,
the relationship between pair completeness and survival depends on
\emph{which patients are pooled}. This is a classic Simpson's paradox: a
trend present in aggregated data reverses in disaggregated subgroups.
The confounding variable is likely the study composition --- studies
within the same cancer type differ in stage, treatment, sequencing
technology, and institutional patient mix, all of which correlate with
both pair completeness and survival.

\textbf{Implications.} The Simpson's paradox result has two
consequences. First, it explains why scalar predictors fail (§3.15): the
cancer-type-level signal is an average over study-level signals that
often point in opposite directions. Averaging conflicting signals
produces a weak, noisy aggregate. Second, it suggests that the true
survival effect of pair completeness is study-specific rather than
cancer-type-specific --- the relevant covariate for predicting direction
is not the tissue of origin but the specific patient population,
treatment context, and sequencing cohort. This motivates prospective,
study-stratified validation rather than pan-cancer aggregation.

\subsubsection{3.19 The Rebinning Offensive --- Cox Hazards, State
Rebinning, and Pair
Discovery}\label{the-rebinning-offensive-cox-hazards-state-rebinning-and-pair-discovery}

The preceding analyses used OLS regression on deceased patients only,
discarding 83,925 censored (alive) observations. They also binned
exclusively by cancer type, leaving open the possibility that the
hallmark state itself --- the 8-bit molecular signature --- is the more
fundamental unit. Finally, the four biological pairs {[}(0,4), (1,3),
(2,5), (6,7){]} were assumed, not discovered. We attacked these three
assumptions simultaneously.

\textbf{Cox Proportional Hazards Hessian (log-hazard space).} We
refitted the full \(2^4\)-state interaction model as a Cox PH regression
on all 145,462 patients with survival data (61,537 events, 83,925
censored). The Cox model outputs log-hazard ratios, which are
dimensionally commensurable with the log-odds coupling matrix \(J\) ---
both live in logarithmic probability space. The Cox Hessian is:

\[H^{\text{Cox}} = \begin{pmatrix} 0 & -0.131 & -0.142 & +0.040 \\ -0.131 & 0 & -0.040 & -0.528 \\ -0.142 & -0.040 & 0 & +0.310 \\ +0.040 & -0.528 & +0.310 & 0 \end{pmatrix} \tag{8}\]

The dominant element is \(H^{\text{Cox}}_{24} = -0.528\) log-HR
(\(z = -7.46\), \(p < 10^{-13}\)): simultaneous completion of the
barrier (Pair 2) and enabler (Pair 4) reduces the hazard rate by
\(1 - e^{-0.528} = 41\%\). This is the Cox analog of the OLS finding
\(H_{24} = +26.1\) months; the sign reversal reflects the hazard-vs-time
inversion (lower hazard = longer survival). The only hazard-increasing
interaction is \(H^{\text{Cox}}_{34} = +0.310\) (\(z = +2.01\),
\(p = 0.045\)): Death/Immune \(\times\) Genome/Metab completion
\emph{increases} hazard by 36\%. Third-order terms remain significant
(\(T_{124} = -0.601\), \(z = -2.37\); \(T_{124} = +0.486\),
\(z = +2.88\)), confirming the pairwise Ising violation. Overall
concordance: \(C = 0.5187\) --- the pair interaction structure, despite
its statistical significance, has essentially zero predictive power at
the individual patient level.

\textbf{Phase space: probability \(\times\) lethality.} Computing the
empirical hazard rate (deaths / person-months) and observation
probability for all 16 macrostates reveals that pair count \(k\) is the
single strongest predictor of hazard across macrostates: \(r = -0.832\)
(\(p = 6.4 \times 10^{-5}\)). The ground state (0000, 85.8\% of
patients) has hazard 0.0154 deaths/month; the fully completed state
(1111, 0.13\% of patients) has hazard 0.0032 --- a \(4.8\times\)
reduction. Probability and lethality are only marginally correlated
(\(r = +0.45\), \(p = 0.078\)), suggesting that the thermodynamic
fitness landscape (which states exist) partially decouples from the
clinical outcome landscape (which states kill).

\textbf{State vs.~cancer type.} Among 63,121 deceased patients with
cancer type annotations, one-way \(\eta^2\) decomposes survival variance
as: hallmark state (8-bit signature) explains 3.89\%, cancer type
explains 16.10\%, and their interaction explains 25.42\%. Cancer type is
\(4.1\times\) more predictive than hallmark state alone. Furthermore, of
22 states observed in \(\geq 3\) cancer types with \(\geq 20\) patients
each, 21 (95\%) show statistically significant survival differences
across cancer types \emph{within the same state} (all \(p < 10^{-9}\)).
Cancer type is not an illusion. However, hallmark state contributes an
incremental 9.33\% of variance on top of cancer type, confirming that
molecular architecture carries independent information. Both covariates
are real; neither is sufficient alone.

\textbf{Random rebinning null model.} A 1,000-permutation test confirms
that cancer-type labels explain far more variance than random labels
(\(\eta^2 = 0.161\) vs.~null mean \(0.001\), \(p < 0.001\)). State
labels similarly survive the null test (\(p < 0.001\)). Both binning
schemes carry genuine signal.

\textbf{Pair structure discovery.} Hierarchical clustering of the
\(8 \times 8\) hallmark co-occurrence matrix (Ward linkage, 4 clusters)
yields clusters \{Immort, Invade, Immune, Genome, Metab\}, \{GrowSup\},
\{Death\}, \{Prolif\} --- a 5-1-1-1 split that matches only Genome+Metab
among the four biological pairs. An exhaustive search over all 105
distinct pairings of 8 hallmarks into 4 pairs, ranked by ANOVA
\(F\)-statistic for pair count \(\rightarrow\) survival, finds the
WINNER\_PAIRS at rank 19 (\(F = 72.9\)). The data-optimal pairing is
(Prolif, Immort) \(|\) (GrowSup, Death) \(|\) (Invade, Immune) \(|\)
(Genome, Metab), with \(F = 188.2\) --- \(2.6\times\) stronger survival
separation. Notably, GrowSup+Death appears in all top-10 pairings,
suggesting that these two hallmarks are more tightly coupled to survival
than the biological pairing (GrowSup, Immort) reflects. The biological
pairs are structurally motivated (Hanahan \& Weinberg); the data-optimal
pairs are survival-motivated. They are not the same.

\textbf{Universal constants from pure states.} Recomputing the three
universal constants from the full patient pool (ignoring cancer type):
\(\beta \cdot \alpha_c^2 = 2.498\) (target 5/2, deviation 0.08\%) and
\(\lambda_1(J)/\pi = 1.497\) (target 3/2, deviation 0.23\%) hold with
extraordinary precision. Bootstrap 95\% CI for \(\lambda_1/\pi\):
\([1.474, 1.519]\). However,
\((1 - f_\mathrm{forb}) \cdot \beta = 0.231\) (target 1, deviation 77\%)
--- this third relation holds only as a per-cancer average, not
globally, because the global forbidden fraction
(\(f_\mathrm{forb} = 0.858\)) is dominated by the 85.8\% of patients in
the ground state.

\textbf{K/J diagnostic in proper units.} Dividing the Cox Hessian by the
coupling matrix element-wise yields the dimensionless thermodynamic
efficiency \(K/J\): the lethality cost per unit of evolutionary effort.
The mean \(K/J = -0.043\) --- pair interactions are net protective. The
most protective pair is barrier \(\times\) enabler (\(K/J = -0.230\));
the only hazard-increasing interaction is immune \(\times\) enabler
(\(K/J = +0.178\)). The spectral ratio (leading eigenvalue of
\(H^{\text{Cox}}\) / leading eigenvalue of \(J\)) is \(+0.135\).

\textbf{Per-cancer Cox Hessian vs.~survival direction.} Fitting Cox PH
models independently for 34 cancer types and correlating with the
balance \(z\)-score: total interaction \(r = +0.068\) (\(p = 0.70\)) ---
dead, consistent with OLS. \(H_{34}\) specifically: \(r = +0.156\)
(\(p = 0.38\)) --- the OLS signal (\(r = -0.52\), \(p = 0.005\)) does
not survive the Cox framework. But per-cancer concordance correlates
with \(z_\mathrm{balance}\) at \(r = +0.557\)
(\(p = 6.2 \times 10^{-4}\)): cancer types where the Cox model fits well
are those where pair completion improves survival. The model's
explanatory power is itself a predictor of the outcome it explains --- a
meta-signal that suggests the Ising framework captures a real but
heterogeneous biological process.

\subsubsection{3.20 The Grating Smasher --- Testing the Pair Assumption
Itself}\label{the-grating-smasher-testing-the-pair-assumption-itself}

The preceding sections parameterize survival using four hallmark pairs
inherited from Hanahan and Weinberg's biological framework. Section 3.22
showed that these pairs rank 19th of 105 possible pairings for survival
separation. We now remove the pair assumption entirely and ask: what
structure does the data actually want?

\textbf{Individual hallmark Cox model.} Fitting a Cox PH model on all 8
hallmarks independently (145,462 patients, 61,537 events) yields
concordance \(C = 0.569\) --- already higher than any pair-based
representation. The hallmark ranking by absolute \(z\)-statistic is:

\begin{center}
\begin{tabular}{r l r r l}
\hline
\textbf{Rank} & \textbf{Hallmark} & \textbf{Coefficient} & $\boldsymbol{z}$ & \textbf{Direction} \\
\hline
1 & Growth Suppression Evasion & $+0.399$ & $+37.1$ & Hazard-increasing \\
2 & Invasion & $-0.249$ & $-21.0$ & Protective \\
3 & Proliferative Signaling & $+0.107$ & $+13.0$ & Hazard-increasing \\
4 & Immune Evasion & $-0.366$ & $-12.3$ & Protective \\
5 & Deregulated Metabolism & $-0.251$ & $-12.2$ & Protective \\
6 & Replicative Immortality & $-0.166$ & $-8.1$ & Protective \\
7 & Genome Instability & $-0.034$ & $-3.8$ & Protective \\
8 & Resisting Cell Death & $+0.031$ & $+3.0$ & Hazard-increasing \\
\hline
\end{tabular}
\end{center}

All eight hallmarks are significant (\(p < 0.003\)). The critical
finding: GrowSup is by far the dominant survival predictor, followed by
Invade and Prolif. Adding all 28 pairwise interactions (36-term model)
raises concordance to \(C = 0.577\), with GrowSup \(\times\) Death as
the strongest interaction (\(z = +13.2\), \(p < 10^{-38}\)). This
interaction reflects the fact that TP53 mutations simultaneously disable
both growth suppression and apoptosis --- the data recovers the known
biology of the single most important tumor suppressor gene.

\textbf{Core--Periphery decomposition.} Spectral analysis of the
\(8 \times 8\) log-odds co-occurrence matrix reveals a single dominant
eigenvalue \(\lambda_1 = 3.22\) capturing 69.1\% of the interaction
variance, with a spectral gap of 2.93 to the next eigenvalue. The
leading eigenvector loads exclusively on five hallmarks: Immort
(\(-0.54\)), Immune (\(-0.51\)), Metab (\(-0.46\)), Genome (\(-0.37\)),
Invade (\(-0.32\)) --- the ``malignancy core.'' The remaining three
hallmarks (Prolif, GrowSup, Death) have near-zero loadings
(\(|v| < 0.05\)), forming a loosely coupled periphery. The participation
ratio (effective dimensionality) is 1.92 --- the hallmark space is
essentially two-dimensional.

Testing Core count (5 hallmarks) vs.~Periphery count (3 hallmarks) as
Cox predictors: Periphery count alone (\(C = 0.543\), \(z = +35.2\)
hazard-increasing) outperforms Core count alone (\(C = 0.529\),
\(z = -29.6\) protective). Combined: \(C = 0.558\). With interaction
term: \(C = 0.560\), where the interaction is strongly negative
(\(z = -22.3\)): having both core and periphery hallmarks is protective
relative to having only periphery, consistent with the domain-wall
lethality mechanism. Both WINNER\_PAIRS (\(C = 0.518\)) and
OPTIMAL\_PAIRS (\(C = 0.534\)) are inferior to the raw Core + Periphery
decomposition.

\textbf{Exhaustive grouping search.} Testing every possible partition of
8 hallmarks into 2 groups (127 partitions), 3 groups (966 partitions),
and 4 groups (105 pairings): the best groupings by survival
\(F\)-statistic scale \emph{inversely} with group count:

\begin{center}
\begin{tabular}{r r l}
\hline
\textbf{Groups} & \textbf{Best $F$} & \textbf{Best partition} \\
\hline
2 & 287.9 & (GrowSup, Death) $\|$ (rest) \\
3 & \textbf{344.7} & (Prolif, Invade, Immune, Genome) $\|$ (GrowSup, Death) $\|$ (Immort, Metab) \\
4 & 188.2 & (Prolif, Immort) $\|$ (GrowSup, Death) $\|$ (Invade, Immune) $\|$ (Genome, Metab) \\
\hline
\end{tabular}
\end{center}

Three groups is optimal (\(F = 344.7\), \(p = 5.3 \times 10^{-222}\)).
The winning three-group partition separates: (a) the TP53 bond (GrowSup
+ Death), (b) the telomere--metabolism axis (Immort + Metab), and (c)
everything else. Notably, GrowSup + Death appears as a distinct group in
\emph{all} top-10 three-group partitions, and the top-1 two-group split
is exactly (GrowSup, Death) vs.~the other six. TP53 is the organizing
principle.

\textbf{LASSO feature selection.} L1-regularized regression on all 92
features (8 main + 28 pairs + 56 triples) retains 46 non-zero
coefficients (\(R^2 = 0.017\) on \(\log(\mathrm{OS})\)). The top two
LASSO features are Prolif \(\times\) Death (\(\beta = +0.143\), longer
survival) and Prolif \(\times\) GrowSup \(\times\) Death
(\(\beta = -0.136\), shorter survival). Of 46 survivors, 27 are
third-order interactions --- the pairwise Ising model misses more than
half the signal that L1 regularization retains.

\textbf{Mutual information.} Nonparametric MI analysis confirms: total
hallmark count (\(\mathrm{MI} = 0.048\)) carries more survival
information than any pair-based count (WINNER \(\mathrm{MI} = 0.020\);
OPTIMAL \(\mathrm{MI} = 0.024\)). Periphery count
(\(\mathrm{MI} = 0.042\)) exceeds Core count (\(\mathrm{MI} = 0.029\)).
Among individual hallmarks, Prolif (\(\mathrm{MI} = 0.023\)) is most
informative; Immune (\(\mathrm{MI} = 0.004\)) least.

\textbf{Re-piping the constants with optimal pairs.} Recomputing the
Ising model with the \(F = 188.2\) optimal pairing: \(\beta = 1.608\),
\(R^2 = 0.919\) (vs.~0.994 for WINNER). The Boltzmann fit degrades
because 61.4\% of patients complete Pair 2 (GrowSup, Death) --- the TP53
bond is so common that the exponential decay assumption breaks. The
coupling matrix becomes asymmetric: \(J_{13} = 2.36\) (Prolif+Immort
\(\times\) Invade+Immune) and \(J_{14} = 1.84\) dominate, while
\(J_{12} = 0.22\) (TP53 bond \(\times\) Invade+Immune) nearly vanishes
--- the TP53 pair decouples from the rest. The universal constants
shift: \(\beta \cdot \alpha_c^2 = 2.466\) (deviation 1.4\%),
\(\lambda_1/\pi = 1.350\) (deviation 10.0\%),
\((1 - f_\mathrm{forb}) \cdot \beta = 1.009\) (deviation 0.9\%). Two of
three constants \emph{worsen}, but the third (the one that failed
globally for WINNER pairs) now hits target. The constants are
\emph{pair-choice-dependent}, not universal --- they characterize the
specific mathematical representation, not the underlying biology.

\textbf{Per-cancer hallmark ranking.} Fitting individual Cox models for
35 cancer types: GrowSup is the top predictor in 8 cancers (including
breast, prostate, head/neck, renal), Death in 7 (lung, pancreas,
prostate, cholangiocarcinoma), Genome in 5, Invade in 5, Prolif in 5. No
single hallmark dominates universally. The glioma signal (concordance
\(C = 0.755\), far above all others) is driven by Metab (\(z = -28.3\)),
confirming that the IDH1 bifurcation --- a metabolic enzyme mutation ---
is the strongest single-hallmark signal in the dataset.

\textbf{The definitive concordance scoreboard:}

\begin{center}
\begin{tabular}{l r r}
\hline
\textbf{Model} & $\boldsymbol{C}$\textbf{-index} & $\boldsymbol{\Delta}$ \textbf{above random} \\
\hline
Random baseline & 0.500 & 0.000 \\
$k_\mathrm{old}$ (WINNER pairs) & 0.518 & +0.018 \\
$k_\mathrm{new}$ (OPTIMAL pairs) & 0.534 & +0.034 \\
$n_\mathrm{total}$ (hallmark count) & 0.521 & +0.021 \\
$n_\mathrm{core}$ (5 hallmarks) & 0.529 & +0.029 \\
$n_\mathrm{core} + n_\mathrm{peri}$ & 0.558 & +0.058 \\
Domain walls (WINNER) & 0.548 & +0.048 \\
All 8 hallmarks (individual) & \textbf{0.569} & \textbf{+0.069} \\
All 36 terms (8 + 28 interactions) & \textbf{0.577} & \textbf{+0.077} \\
\hline
\end{tabular}
\end{center}

\textbf{Verdict.} The pair structure --- \emph{any} pair structure ---
is a lossy compression of the 8-dimensional hallmark space. The data's
natural structure is not four pairs but two populations: a tightly
coupled five-hallmark malignancy core and a loosely coupled
three-hallmark periphery dominated by TP53. The optimal number of groups
is three, not four. The universal constants
\(\beta \cdot \alpha_c^2 = 5/2\) and \(\lambda_1/\pi = 3/2\) hold with
remarkable precision for the WINNER pairing specifically --- and degrade
under the survival-optimal pairing. This means the constants
characterize a \emph{specific mathematical representation} of the
hallmark data, not an intrinsic thermodynamic property of cancer. The
representation is meaningful (it captures real co-occurrence structure),
but it is not unique, and it is not optimal for survival prediction.

\begin{center}\rule{0.5\linewidth}{0.5pt}\end{center}

\emph{What survived the falsification cascade is a three-coordinate
system: domain wall density, tissue oxygenation, and a Genome
Instability toggle. These features emerged from the data after the pair
framework failed to predict survival direction.}

\subsubsection{3.21 Domain Walls --- The Ising Frustration
Energy}\label{domain-walls-the-ising-frustration-energy}

A domain wall in the Ising model is the boundary between two ordered
regions. In the hallmark pair framework, a domain wall occurs when
exactly one hallmark of a pair is active --- the ``half-pair'' state.
This is the frustrated intermediate: neither the vacuum (both hallmarks
inactive) nor the ordered phase (both active), but the energetically
costly boundary between them.

For each patient, we counted the number of broken pairs (0 to 4) ---
pairs where one hallmark is mutated and its partner is not. In 61,537
deceased patients with recorded survival:

{\def\LTcaptype{none} % do not increment counter
\begin{longtable}[]{@{}crrr@{}}
\toprule\noalign{}
Broken pairs & N & Median OS (mo) & Mean OS (mo) \\
\midrule\noalign{}
\endhead
\bottomrule\noalign{}
\endlastfoot
0 & 402 & 15.3 & 33.6 \\
1 & 10,888 & 16.1 & 24.6 \\
2 & 23,157 & 14.7 & 22.3 \\
3 & 21,767 & 13.0 & 19.2 \\
4 & 5,323 & 11.7 & 18.5 \\
\end{longtable}
}

The gradient is monotonic: more domain walls → shorter survival
(\(z = +4.37\), \(p = 1.26 \times 10^{-5}\) for walls \(= 0\) vs.~walls
\(\geq 1\)). Patients with no frustrated pairs survive a median of 15.3
months; patients with all four pairs frustrated survive a median of 11.7
months --- a 30\% reduction.

\textbf{Per-pair domain wall analysis.} Each pair's domain wall was
tested individually (having exactly one hallmark active vs.~having both
or neither):

{\def\LTcaptype{none} % do not increment counter
\begin{longtable}[]{@{}
  >{\raggedright\arraybackslash}p{(\linewidth - 8\tabcolsep) * \real{0.2727}}
  >{\raggedleft\arraybackslash}p{(\linewidth - 8\tabcolsep) * \real{0.1818}}
  >{\raggedleft\arraybackslash}p{(\linewidth - 8\tabcolsep) * \real{0.1818}}
  >{\raggedleft\arraybackslash}p{(\linewidth - 8\tabcolsep) * \real{0.1818}}
  >{\raggedleft\arraybackslash}p{(\linewidth - 8\tabcolsep) * \real{0.1818}}@{}}
\toprule\noalign{}
\begin{minipage}[b]{\linewidth}\raggedright
Pair
\end{minipage} & \begin{minipage}[b]{\linewidth}\raggedleft
Domain wall z
\end{minipage} & \begin{minipage}[b]{\linewidth}\raggedleft
p
\end{minipage} & \begin{minipage}[b]{\linewidth}\raggedleft
\(N_\mathrm{wall}\)
\end{minipage} & \begin{minipage}[b]{\linewidth}\raggedleft
\(N_\mathrm{no\text{-}wall}\)
\end{minipage} \\
\midrule\noalign{}
\endhead
\bottomrule\noalign{}
\endlastfoot
Prolif + Invade & \(-12.73\) & \(3.8 \times 10^{-37}\) & 31,960 &
29,577 \\
GrowSup + Immort & \(-19.34\) & \(2.7 \times 10^{-83}\) & 45,668 &
15,869 \\
Death + Immune & \(-8.93\) & \(4.2 \times 10^{-19}\) & 45,556 &
15,981 \\
Genome + Metab & \(-1.42\) & \(0.15\) & 20,611 & 40,926 \\
\end{longtable}
}

The first three pairs show massive effects: having one hallmark of a
pair active without its partner is significantly worse than having both
or neither. The Growth Suppression + Immortality domain wall
(\(p = 2.7 \times 10^{-83}\)) is the strongest single-variable survival
signal in the entire dataset outside the glioma IDH1 bifurcation. The
Genome + Metabolism pair --- the enabler --- shows no significant domain
wall penalty, consistent with its role as a co-activated module where
one mutation (e.g., IDH1) frequently triggers the partner pathway
without requiring a separate mutational event.

\textbf{Interpretation.} The domain wall is the Ising frustration energy
made clinical. A tumor that has activated proliferative signaling
without invasion capability, or evaded growth suppression without
achieving immortality, is in a thermodynamically unstable state --- it
has paid the energetic cost of activating one hallmark without gaining
the stabilizing benefit of the paired capability. This frustrated state
maps directly to worse patient outcomes. The data establish domain wall
density as a patient-level survival feature computable from routine
mutation panels.

\subsubsection{3.22 Forbidden Boundary
Proximity}\label{forbidden-boundary-proximity}

Of the 256 possible 8-bit hallmark states, 94 are forbidden --- never
observed in 198,862 tumors. For each observed state, we computed the
minimum Hamming distance to any forbidden state: how many
single-hallmark flips separate this patient from an impossible
configuration.

In 61,537 deceased patients: those at Hamming distance 1 from the
forbidden boundary (one flip away from impossibility) had significantly
better survival than those at distance \(\geq 2\) (\(z = +7.36\),
\(p = 1.84 \times 10^{-13}\); median OS 14.3 vs.~13.0 months). This
effect is counterintuitive --- proximity to impossibility is protective
--- but follows from the structure of the forbidden set: forbidden
states tend to be high-hallmark configurations requiring the barrier
pairs (Pair 2, Pair 3). Patients near this boundary have more hallmarks
active but in more structured (paired) configurations, reducing domain
wall density.

At the cancer-type level, mean forbidden distance does not predict
survival direction (\(r = -0.27\), \(p = 0.12\); \(r = -0.09\) excluding
glioma). The forbidden boundary proximity is a patient-level survival
feature, not a cancer-type organizer.

\subsubsection{3.23 Frustrated Proliferation: The Domain Wall as a
Coordinate
System}\label{frustrated-proliferation-the-domain-wall-as-a-coordinate-system}

Section 3.17 established that domain walls --- hallmark pairs with
exactly one member active --- predict worse survival. We now ask: does
the \emph{pattern} of domain walls carry additional information beyond
the count? Specifically, can the wall encoding serve as an alternative
coordinate system for the hallmark state, and does it reveal structure
invisible to the flat 8-hallmark representation?

\textbf{The domain wall encoding.} For each of the three data-validated
thermodynamic pairs (TP53: GrowSup × Death; Aerobic: Prolif × Invade;
Hypoxic: Immort × Metab), we replace the two binary hallmarks
\((h_A, h_B)\) with three domain-wall features: \emph{both active}
(\(h_A \wedge h_B\), the ``coupled'' state), \emph{A-only}
(\(h_A \wedge \neg h_B\)), and \emph{B-only} (\(\neg h_A \wedge h_B\)).
This yields 9 features from 3 pairs plus Genome Instability as a
standalone toggle (Feature 10), for a total of 10 wall features. The
remaining hallmark (Immune Evasion) is dropped: the sensor blindness
audit (§3.31) shows that only JAK1 (2\% prevalence) instruments this
capability, making it pure noise under the 45-gene panel. Thus the
fourth state of each pair (``neither active'') and the Immune dimension
are absorbed into the intercept. The encoding is a lossless
re-parameterization of the informative 7 hallmark bits, reorganized to
expose wall structure.

\textbf{Pan-cancer wall model.} Fitting a Cox PH model with 10 wall
features + 2 O₂ indicators (= 12 covariates) on 48,675 patients yields
\(C = 0.6242\) (cross-validated: \(C_\mathrm{CV} = 0.6241\)), compared
to \(C = 0.6214\) for the flat 8-hallmark + O₂ model. The improvement is
\(\Delta C = +0.0028\) from zero additional information --- merely from
reorganizing the same bits into wall coordinates. By AIC, the wall model
beats the flat model by \(\Delta\mathrm{AIC} = -220\) pan-cancer and
\(-70\) in HIGH O₂ tissues.

\textbf{The frustrated proliferator.} The dominant wall feature is
\(\mathrm{Prolif\_only}\): Proliferative Signaling active without
Invasion capability. Pan-cancer: \(\beta = +0.2487\), HR \(= 1.282\),
\(p = 4.53 \times 10^{-56}\). A tumor that has activated proliferative
signaling but cannot invade faces 28\% higher hazard than the reference.
This signal is not glioma-specific: removing all diffuse glioma
patients, Prolif\_only retains \(\beta = +0.608\), \(p < 10^{-115}\) ---
it \emph{strengthens} without glioma. Prolif\_only is the most robust
single hallmark-derived survival feature in the dataset.

\textbf{O2 routing of the frustrated proliferator.} Stratifying by
tissue oxygenation reveals that frustrated proliferation is massively
O2-dependent:

{\def\LTcaptype{none} % do not increment counter
\begin{longtable}[]{@{}lllll@{}}
\toprule\noalign{}
O2 Tier & Invade Status & Prolif \(\beta\) & HR & \(p\) \\
\midrule\noalign{}
\endhead
\bottomrule\noalign{}
\endlastfoot
LOW & Invade = 0 & \(+0.594\) & 1.812 & \(< 10^{-145}\) \\
LOW & Invade = 1 & \(+0.321\) & 1.379 & \(6 \times 10^{-5}\) \\
MEDIUM & Invade = 0 & \(-0.134\) & 0.875 & --- \\
MEDIUM & Invade = 1 & \(-0.614\) & 0.541 & \(< 10^{-23}\) \\
HIGH & Either & NS & \textasciitilde1.0 & \(> 0.05\) \\
\end{longtable}
}

In low-O2 tissues, proliferation without invasion is catastrophic (HR =
1.81, 81\% excess hazard). But in medium-O2 tissues, proliferation
\emph{with} invasion is massively protective (HR = 0.54, halving the
hazard). Invasion is the escape valve: in well-enough-oxygenated tissue,
a proliferating tumor that can invade escapes to better-vascularized
territory, converting a lethal trajectory into a survivable one.

The formal interaction term
\(\mathrm{Prolif} \times (1 - \mathrm{Invade})\) is significant in
MEDIUM O2 (\(\beta = +0.420\), \(p = 5 \times 10^{-11}\)) and LOW O2
(\(\beta = +0.284\), \(p = 3 \times 10^{-5}\)), but non-significant in
HIGH O2 (\(p = 0.86\)). The interaction activates only under hypoxic
pressure --- in well-oxygenated tissue, the invasion escape is
unnecessary because the tumor already has adequate oxygen supply.

\textbf{What is NOT real: XOR and ``thermodynamic frustration.''} The
XOR encoding --- counting states where exactly one hallmark of a pair is
active --- was tested as a predictor of survival beyond what the
individual hallmarks provide. Under cross-validation, XOR adds
\(\Delta C = +0.0009\) pan-cancer and \(\Delta C = +0.0000\) in LOW O2.
The domain wall's power comes not from the XOR \emph{count} but from the
\emph{identity} of which hallmark is stranded: Prolif without Invade is
lethal; Invade without Prolif is not.

Variance inflation analysis confirms: the set \{Prolif, Invade,
Prolif\_only\} has VIF \(= 5.93\) because Prolif\_only \(\equiv\) Prolif
\$\times (1 - \$ Invade\()\) --- it is exactly rank-deficient. The wall
encoding is a mathematical re-parameterization, not new physics. Its
value lies in parsimony: the 12-parameter wall model achieves the same
predictive power as a 26-parameter interaction model
(\(\Delta C < 0.001\) between them) while being more interpretable.

\textbf{Robustness.} Prolif\_only survives every leave-one-out test
(removing any single cancer type: \(p < 10^{-21}\) always). The wall
model is validated by 5-fold cross-validation
(\(C_\mathrm{CV} = 0.6241\), overfitting \(< 0.002\)). The AIC advantage
persists in HIGH O2 (\(\Delta\mathrm{AIC} = -70\)) and PAN-CANCER
(\(-220\)), though in LOW O2 the flat model wins by
\(\Delta\mathrm{AIC} = +14\) --- the wall encoding's advantage is driven
by aerobic tissue, where the Prolif/Invade axis is most informative.

\textbf{What is DIFG-specific: the Hypoxic equilibrium.} One wall
feature, \(\mathrm{Hypoxic\_coupled}\) (simultaneous activation of
Replicative Immortality and Deregulated Metabolism), appears significant
pan-cancer (\(p < 10^{-51}\)). However, removing diffuse glioma alone
collapses this signal to \(p = 0.79\). The prevalence of the
Hypoxic\_coupled state is 24.1\% in DIFG and \(< 0.5\%\) in all other
LOW O2 cancer types. This feature describes the IDH-mutant glioma
biology discovered in §3.12, not a pan-cancer phenomenon.

\subsubsection{3.24 Escape Routes, Survivor Bias, and the Oxygenation
Axis}\label{escape-routes-survivor-bias-and-the-oxygenation-axis}

The master equation reveals a paradox: Basin A (few hallmarks) has
higher hazard than Basin B (many hallmarks). Claude.AI raised a critical
objection: is this real biology, or survivor bias --- the artifact that
long-surviving patients accumulate more mutations simply because they
live long enough to be sequenced? We tested this directly.

\textbf{The survivor bias verdict: it's complicated.} Overall, Basin B
patients survive 12.8 months longer (median 54.4 vs.~41.6, log-rank
\(p = 2.5 \times 10^{-75}\)). Controlling for hallmark count, Basin B
wins at \(k = 1\) (\(p = 4.2 \times 10^{-16}\)), \(k = 2\)
(\(p = 4.6 \times 10^{-6}\)), and \(k = 3\)
(\(p = 1.3 \times 10^{-86}\)), but \emph{loses} at \(k = 4\)
(\(p = 1.6 \times 10^{-4}\)). In a multivariate Cox model adjusting for
hallmark count, cancer type, and basin membership simultaneously, the
basin coefficient becomes non-significant (\(p = 0.14\)). The survival
advantage of Basin B is largely explained by its constituent hallmarks,
not by basin membership per se. Each additional core hallmark reduces
hazard by 19.4\% (HR = 0.806, \(p = 6.1 \times 10^{-14}\)), while each
additional periphery hallmark increases hazard by 32.0\% (HR = 1.320,
\(p = 1.1 \times 10^{-39}\)).

\textbf{The within-cancer-type test is definitive for some cancers.} In
diffuse glioma, Basin B survives 3\(\times\) longer than Basin A at
matched \(k\) (\(k = 2\): 46.0 vs.~14.4 mo; \(k = 3\): 61.5 vs.~19.9 mo;
both \(p < 10^{-10}\)). In prostate cancer at \(k = 2\): 218.0 vs.~56.0
months (\(p = 0.001\)). These signals are too large to be survivor bias
--- no sequencing delay can explain a 4-fold survival difference within
the same hallmark count. In breast and lung, however, the within-\(k\)
tests are non-significant or reversed. The hazard inversion appears to
be cancer-type-specific: real in glioma and prostate, ambiguous in
breast and lung.

\textbf{The escape routes.} Of 45 single-step transitions from Basin A
to Basin B, three exit hallmarks dominate: Genome (15 routes), Immune
(15 routes), and Metab (15 routes). But the population flow is
overwhelmingly through Genome: 45,483 patients reach Basin B via Genome
acquisition vs.~1,059 via Immune and 3,110 via Metab. The
highest-traffic escape: \(01100000 \to 01100010\) (TP53 \(\to\)
TP53+Genome), carrying 11,040 patients.

Survival impact of each escape from Basin A:

{\def\LTcaptype{none} % do not increment counter
\begin{longtable}[]{@{}lrrrl@{}}
\toprule\noalign{}
Escape hallmark & Median before & Median after & \(\Delta\) & \(p\) \\
\midrule\noalign{}
\endhead
\bottomrule\noalign{}
\endlastfoot
+Genome & 41.6 mo & 46.3 mo & \textbf{+4.7 mo} &
\(7.1 \times 10^{-6}\) \\
+Metab & 41.6 mo & 50.8 mo & \textbf{+9.2 mo} &
\(1.1 \times 10^{-5}\) \\
+Immune & 41.6 mo & 38.1 mo & \(-3.6\) mo & 0.30 (ns) \\
\end{longtable}
}

Acquiring Metabolism gives the largest survival benefit (+9.2 months),
followed by Genome (+4.7 months). Acquiring Immune Evasion does nothing.
The worst escape paths all involve Immune: patients who acquire Immune
Evasion from Basin A reach median survival of 18--25 months. The best
multi-step path from the ground state is loss of a periphery hallmark
plus Metab acquisition, reaching destinations with median survival
exceeding 250 months.

\textbf{The oxygenation axis.} Tony's intuition that tissue oxygenation
matters was tested by classifying 75,372 patients into
high-\(\mathrm{O}_2\) tissues (lung, breast, skin, kidney),
medium-\(\mathrm{O}_2\) (colon, liver, bladder, stomach, uterus), and
low-\(\mathrm{O}_2\) (pancreas, prostate, brain, soft tissue) based on
published physiological \(p\mathrm{O}_2\) values:

{\def\LTcaptype{none} % do not increment counter
\begin{longtable}[]{@{}
  >{\raggedright\arraybackslash}p{(\linewidth - 8\tabcolsep) * \real{0.1579}}
  >{\raggedleft\arraybackslash}p{(\linewidth - 8\tabcolsep) * \real{0.2105}}
  >{\raggedleft\arraybackslash}p{(\linewidth - 8\tabcolsep) * \real{0.2105}}
  >{\raggedleft\arraybackslash}p{(\linewidth - 8\tabcolsep) * \real{0.2105}}
  >{\raggedleft\arraybackslash}p{(\linewidth - 8\tabcolsep) * \real{0.2105}}@{}}
\toprule\noalign{}
\begin{minipage}[b]{\linewidth}\raggedright
\end{minipage} & \begin{minipage}[b]{\linewidth}\raggedleft
High \(\mathrm{O}_2\)
\end{minipage} & \begin{minipage}[b]{\linewidth}\raggedleft
Medium \(\mathrm{O}_2\)
\end{minipage} & \begin{minipage}[b]{\linewidth}\raggedleft
Low \(\mathrm{O}_2\)
\end{minipage} & \begin{minipage}[b]{\linewidth}\raggedleft
H/L ratio
\end{minipage} \\
\midrule\noalign{}
\endhead
\bottomrule\noalign{}
\endlastfoot
Mean \(\langle k \rangle\) & 2.57 & 2.90 & 2.29 & 1.12 \\
Invade prevalence & \textbf{23.3\%} & 24.7\% & \textbf{7.8\%} &
\textbf{3.00\(\times\)} \\
Metab prevalence & 2.6\% & 4.9\% & \textbf{11.5\%} &
\textbf{0.23\(\times\)} \\
Immort prevalence & 4.2\% & 4.8\% & \textbf{9.4\%} &
\textbf{0.45\(\times\)} \\
Prolif prevalence & \textbf{56.3\%} & 49.1\% & \textbf{33.2\%} &
1.69\(\times\) \\
Genome prevalence & 37.2\% & 43.7\% & 32.0\% & 1.16 \\
Median OS & \textbf{49.5 mo} & 29.2 mo & \textbf{30.8 mo} & --- \\
Death rate & \textbf{36.5\%} & 40.2\% & \textbf{50.3\%} & --- \\
\end{longtable}
}

The oxygenation signal is striking. Invasion is 3\(\times\) more
prevalent in high-\(\mathrm{O}_2\) tissues --- tumors in well-oxygenated
tissue can afford the metabolic cost of invasion. Deregulated Metabolism
is \(4.3\times\) more prevalent in low-\(\mathrm{O}_2\) tissues ---
hypoxic tumors must rewire metabolism (Warburg effect) to survive.
Immortality is \(2.2\times\) more common in hypoxic tissue. The hallmark
profile of a tumor is shaped not just by its mutations but by the oxygen
environment of its host tissue.

In the Cox model, tissue oxygenation is the strongest single predictor
of survival: high-\(\mathrm{O}_2\) (HR = 0.593, \(p < 10^{-183}\)),
low-\(\mathrm{O}_2\) (HR = 1.321, \(p < 10^{-54}\)). Even after
adjusting for basin, hallmark count, and Genome status, oxygenation
remains an independent predictor with the largest hazard ratio in the
model.

The survival impact of Genome acquisition reverses by oxygenation: in
high-\(\mathrm{O}_2\) tissues, Genome \emph{worsens} survival (103.0
\(\to\) 90.8 mo, \(p = 0.001\)). In low-\(\mathrm{O}_2\) tissues, Genome
\emph{improves} survival (33.7 \(\to\) 52.5 mo,
\(p = 8.3 \times 10^{-38}\)). The Genomic Rubicon is protective only in
hypoxic environments. Genomic instability may enable metabolic
adaptation in oxygen-starved tumors while serving no compensatory
function in well-oxygenated tissue.

\subsubsection{3.25 The Dimensional Attack: Testing Whether 8 Hallmarks
Are an
Illusion}\label{the-dimensional-attack-testing-whether-8-hallmarks-are-an-illusion}

The oxygenation axis raises a radical question: do 8 independent
hallmarks truly exist, or are they a lower-dimensional system inflated
by mixing aerobic and anaerobic tumor evolution? We tested this by
computing the full \(8 \times 8\) hallmark correlation matrix
independently within each oxygenation regime and performing
eigendecomposition at the patient level (75,372 patients: 31,980
high-\(\mathrm{O}_2\), 26,892 low-\(\mathrm{O}_2\), 16,500
medium-\(\mathrm{O}_2\)).

\textbf{Effective dimensionality is high, not low.} The effective
dimensionality of the correlation matrix is 7.30 (high
\(\mathrm{O}_2\)), 7.04 (medium), and 6.74 (low). By the
Marchenko-Pastur random matrix threshold (\(\lambda_+ = 1.03\)), exactly
3 eigenvalues are statistically significant above noise in both regimes.
Seven components are needed to capture 90\% of the correlation variance.
The 8 hallmarks do not collapse to 2 dimensions within either
oxygenation regime.

\textbf{The dominant axes are partially aligned, not orthogonal.} The
PC1 eigenvector alignment between high-\(\mathrm{O}_2\) and
low-\(\mathrm{O}_2\) is
\(|\mathbf{v}_1^{\mathrm{high}} \cdot \mathbf{v}_1^{\mathrm{low}}| = 0.789\)
--- substantially aligned, not orthogonal. An earlier analysis on the
\(4 \times 4\) pair coupling matrix reported alignment of 0.983; a text
summary of those results claimed 0.116, which was fabricated and
contradicted the script's own output. The true value at the raw hallmark
level is 0.789 --- the principal axis of variation is \emph{similar}
across oxygenation regimes but not identical.

The PC1 loadings reveal a shared structure: in \emph{both} regimes,
Genome (+0.49/+0.55) and Immort (+0.36/+0.47) load positively while
GrowSup (−0.37/−0.37) and Death (−0.40/−0.40) load negatively. The
core-vs-periphery axis (§3.20) is present in both oxygenation regimes.
What differs: Invade loads on PC1 in high-\(\mathrm{O}_2\) (+0.37) but
vanishes in low-\(\mathrm{O}_2\) (+0.04); Metab is stronger in
low-\(\mathrm{O}_2\) (+0.36 vs +0.25).

\textbf{The 2-score model captures oxygenation but loses 40-52\% of
information.} Compressing the 8-bit hallmark signature to a 2-score
system (TP53 axis + O2-adapted axis) reduces 124--139 states to 9 bins,
losing 40\% (high-\(\mathrm{O}_2\)) to 52\% (low-\(\mathrm{O}_2\)) of
information entropy. In survival prediction, the 2-score model achieves
\(C = 0.611\) versus \(C = 0.568\) for flat 8-hallmark (a gain from
adding oxygenation), but the full 8-hallmark + O2 model reaches
\(C = 0.621\), and adding 16 O2-hallmark interaction terms reaches
\(C = 0.634\). The 8 individual hallmarks carry information that the
2-score compression destroys.

\textbf{Only 2 of 8 hallmarks flip sign by oxygenation.} If the
bifurcation hypothesis were correct, most hallmark effects should
reverse direction between oxygenation regimes. In fact, only Invade (HR
1.00 \(\to\) 0.86) and Immune (HR 1.06 \(\to\) 0.72) flip sign. The
other 6 hallmarks maintain their hazard direction across all oxygenation
levels. The system is modulated by oxygen, not bifurcated by it.

\textbf{The mutual information structure is conserved.} The top two MI
pairs are the same in both regimes: GrowSup-Death (MI = 0.103/0.186
bits) and Immort-Genome (MI = 0.067/0.153 bits). Low-\(\mathrm{O}_2\)
adds a third strong pair (Immort-Metab, MI = 0.055) absent in
high-\(\mathrm{O}_2\), consistent with the Warburg effect coupling
metabolism to immortality under hypoxia.

\textbf{Per-O2 concordance reveals the real structure.} Within
high-\(\mathrm{O}_2\) tissues alone, 8-hallmark concordance is
\(C = 0.563\); within low-\(\mathrm{O}_2\), \(C = 0.630\). The compass
is substantially more predictive in hypoxic cancers. The dominant
signals within low-\(\mathrm{O}_2\): Prolif (HR = 1.66,
\(p < 10^{-110}\)), Death (HR = 1.39, \(p < 10^{-30}\)), and Metab (HR =
0.44, \(p < 10^{-85}\)) --- Metabolic reprogramming more than halves the
hazard in hypoxic tissue.

\textbf{Initial verdict.} The 8 hallmarks carry 6.7--7.3 effective
\emph{variance} dimensions, 3 of which are statistically significant
above random matrix noise. The principal axes are 79\% aligned across
oxygenation regimes, and only 2 of 8 hallmarks change sign. However,
this initial analysis confuses variance spread with signal
concentration. The fact that 7 components are needed for 90\% of
\emph{correlation variance} does not mean 7 components are needed for
\emph{survival prediction}. Subsequent analysis (§3.26--3.28) resolves
this: the signal is 3-dimensional even though the noise is
7-dimensional.

\subsubsection{3.26 The σ-Parity Fracture: Genome Instability as the
O2-Gated Coupling
Term}\label{the-ux3c3-parity-fracture-genome-instability-as-the-o2-gated-coupling-term}

The dimensional attack (§3.25) established that exactly 3 eigenvalues
carry signal. The escape path analysis (§3.24) identified Genome
Instability as the Rubicon whose survival effect inverts by oxygenation.
We now test whether Genome operates as a universal σ-parity breaker
(affecting survival through combinatorial symmetry) or as a conditional
coupling term (activated only under specific environmental conditions).

\textbf{Design.} We stratified 48,675 patients with survival data by
tissue oxygenation (high vs.~low \(\mathrm{O}_2\)) and computed a binary
Hypoxic Score (sum of hallmarks enriched in hypoxic tissue: Immort,
Genome, Metab, weighted by low-\(\mathrm{O}_2\) loadings). Within each
O2 stratum, we fit a Cox proportional hazards model with Genome, Hypoxic
Score, and their interaction.

\textbf{The coupling is O2-gated.} In low-\(\mathrm{O}_2\) tissues, the
Genome \(\times\) Hypoxic interaction is strongly significant:
\(\beta = +0.372\), HR \(= 1.450\), \(p = 7.21 \times 10^{-10}\). In
high-\(\mathrm{O}_2\) tissues, the same interaction is non-significant:
\(\beta = -0.108\), \(p = 0.315\). Genome Instability couples to the
hypoxic hallmark axis \emph{only} when the tissue is actually hypoxic.
The coupling turns off in aerobic tissue.

\textbf{Parity counting is a combinatorial artifact.} We tested whether
the 8-hallmark σ-parity (even vs.~odd number of active hallmarks)
carries survival information beyond what individual hallmarks provide.
The parity-count χ² test yields \(\chi^2 = 0.00\) (\(p = 1.00\)). The
fracture is in the \emph{coupling}, not in the combinatorial
distribution --- parity itself is not a biological variable.

\textbf{Patient-level survival confirms the coupling.} In
low-\(\mathrm{O}_2\) tissue, patients with Genome+ and even σ-parity
survive a median 59.4 months, while Genome− patients with odd σ-parity
survive 30.0 months --- a 29.4-month gap (\(p < 10^{-9}\)). In
high-\(\mathrm{O}_2\) tissue, the corresponding gap is only 5.2 months
(\(p > 0.05\)). Genome Instability is prognostically active only under
hypoxic selective pressure.

\textbf{Cross-talk analysis.} Mutual information between Genome and
other hallmarks shifts by oxygenation regime: - Genome
\(\leftrightarrow\) Aerobic hallmarks: MI \(= 0.000\) bits (high
\(\mathrm{O}_2\)) vs.~\(0.018\) bits (low \(\mathrm{O}_2\)) - Genome
\(\leftrightarrow\) TP53: MI \(= 0.012\) bits (high) vs.~\(0.025\) bits
(low)

In aerobic tissue, Genome is informationally \emph{decoupled} from the
aerobic hallmark axis. In hypoxic tissue, the coupling activates. This
is the von Mangoldt analog of cancer: Genome Instability is not a
universal survival predictor --- it is the conditional coupling term
whose activation depends on the environmental parameter (tissue
oxygenation).

\textbf{Per-O2 concordance.} Within-stratum 8-hallmark Cox models yield
\(C = 0.560\) (high \(\mathrm{O}_2\)) vs.~\(C = 0.612\) (low
\(\mathrm{O}_2\)). The compass is substantially more predictive in
hypoxic cancers, consistent with the additional coupling term being
active.

\subsubsection{3.27 The 3-Dimensional Cancer Manifold: Convergence of
Three Independent
Tests}\label{the-3-dimensional-cancer-manifold-convergence-of-three-independent-tests}

Three independent lines of evidence --- spectral analysis, predictive
compression, and conditional coupling --- converge on the same
conclusion: cancer survival in this dataset is governed by a
3-dimensional manifold embedded in 8-dimensional hallmark space.

\textbf{Test 1: Marchenko-Pastur eigenvalue analysis.} The
\(8 \times 8\) hallmark correlation matrix (75,372 O2-classified
patients) yields exactly 3 eigenvalues above the random matrix noise
threshold in all three oxygenation regimes:

{\def\LTcaptype{none} % do not increment counter
\begin{longtable}[]{@{}
  >{\raggedright\arraybackslash}p{(\linewidth - 6\tabcolsep) * \real{0.1159}}
  >{\raggedright\arraybackslash}p{(\linewidth - 6\tabcolsep) * \real{0.0725}}
  >{\raggedright\arraybackslash}p{(\linewidth - 6\tabcolsep) * \real{0.4058}}
  >{\raggedright\arraybackslash}p{(\linewidth - 6\tabcolsep) * \real{0.4058}}@{}}
\toprule\noalign{}
\begin{minipage}[b]{\linewidth}\raggedright
Regime
\end{minipage} & \begin{minipage}[b]{\linewidth}\raggedright
\(n\)
\end{minipage} & \begin{minipage}[b]{\linewidth}\raggedright
\(\lambda_+\) (MP threshold)
\end{minipage} & \begin{minipage}[b]{\linewidth}\raggedright
Eigenvalues above threshold
\end{minipage} \\
\midrule\noalign{}
\endhead
\bottomrule\noalign{}
\endlastfoot
All patients & 75,372 & 1.026 & \(\lambda_1 = 1.564\),
\(\lambda_2 = 1.464\), \(\lambda_3 = 1.180\) \\
High \(\mathrm{O}_2\) & 31,980 & 1.040 & \(\lambda_1 = 1.564\),
\(\lambda_2 = 1.180\) \\
Low \(\mathrm{O}_2\) & 26,892 & 1.045 & \(\lambda_1 = 1.564\),
\(\lambda_2 = 1.464\), \(\lambda_3 = 1.180\) \\
\end{longtable}
}

The effective rank at the 90\% variance threshold is 7 in all regimes
--- variance is spread broadly. But only 3 eigenvalues are statistically
distinguishable from noise. The distinction between variance spread and
signal concentration is critical: 5 of the 8 hallmark dimensions carry
correlated noise, not survival-relevant signal.

\textbf{Test 2: Concordance compression.} Cox proportional hazards
models on \(n = 48{,}675\) patients with survival data:

{\def\LTcaptype{none} % do not increment counter
\begin{longtable}[]{@{}
  >{\raggedright\arraybackslash}p{(\linewidth - 8\tabcolsep) * \real{0.1373}}
  >{\raggedright\arraybackslash}p{(\linewidth - 8\tabcolsep) * \real{0.2157}}
  >{\raggedright\arraybackslash}p{(\linewidth - 8\tabcolsep) * \real{0.0980}}
  >{\raggedright\arraybackslash}p{(\linewidth - 8\tabcolsep) * \real{0.3137}}
  >{\raggedright\arraybackslash}p{(\linewidth - 8\tabcolsep) * \real{0.2353}}@{}}
\toprule\noalign{}
\begin{minipage}[b]{\linewidth}\raggedright
Model
\end{minipage} & \begin{minipage}[b]{\linewidth}\raggedright
Parameters
\end{minipage} & \begin{minipage}[b]{\linewidth}\raggedright
\(C\)
\end{minipage} & \begin{minipage}[b]{\linewidth}\raggedright
\(C_\mathrm{CV}\)
\end{minipage} & \begin{minipage}[b]{\linewidth}\raggedright
Efficiency
\end{minipage} \\
\midrule\noalign{}
\endhead
\bottomrule\noalign{}
\endlastfoot
8 hallmarks (flat) & 8 & 0.568 & 0.567 & --- \\
3-score + O2 & 5 & 0.611 & 0.610 & 96.3\% of full \\
8 hallmarks + O2 (flat) & 10 & 0.621 & 0.621 & --- \\
Wall + O2 & 12 & 0.624 & 0.624 & --- \\
8 hallmarks + O2 + 16 interactions & 25 & 0.634 & 0.634 & 100\%
(reference) \\
\end{longtable}
}

The 3-score model uses: (1) TP53 score (\(\beta = +0.149\),
\(p = 3.1 \times 10^{-57}\)), (2) O2-adaptive score (\(\beta = -0.182\),
\(p = 1.4 \times 10^{-47}\)), and (3) Genome Instability bit
(\(\beta = -0.077\), \(p = 3.3 \times 10^{-7}\)), plus O2 regime
indicators (high: \(p = 5.1 \times 10^{-80}\); low:
\(p = 8.7 \times 10^{-28}\)). All 5 parameters are significant at
\(p < 10^{-6}\). The model achieves 96.3\% of the predictive power of
the full 25-parameter interaction model with 80\% fewer parameters.

\textbf{Test 3: O2-gated coupling (§3.26).} Genome Instability couples
to survival only in low-\(\mathrm{O}_2\) tissue
(\(p = 7.2 \times 10^{-10}\)), not in high-\(\mathrm{O}_2\)
(\(p = 0.32\)). This is the third dimension: a conditional coupling
term, not a universal covariate. The 3 dimensions are:

\begin{enumerate}
\def\labelenumi{\arabic{enumi}.}
\tightlist
\item
  \textbf{The TP53 axis} --- GrowSup and Death, dominated by TP53
  mutations, universally hazard-increasing
\item
  \textbf{The O2-adaptive axis} --- Genome, Metab, Immort, modulated by
  tissue oxygenation
\item
  \textbf{The Genome toggle} --- Genome Instability as a binary switch
  whose survival coupling activates only under hypoxic selective
  pressure
\end{enumerate}

\textbf{What this means for the original framework.} The pair structure
(§3.1--3.18) and Ising model (§3.11) describe the \emph{variance}
landscape accurately: 4 pairs, 6 couplings, cooperative condensation
dynamics. But variance and signal are not the same thing. Five of the 8
hallmark dimensions contribute correlated noise to the variance
structure without contributing independently to survival prediction. The
original 8-hallmark framework remains descriptively valid --- all 8
hallmarks are real biological capabilities --- but for \emph{prognostic}
purposes, the effective model is 3 scores + tissue oxygenation.
Cross-validation (§3.29), domain wall re-parameterization (§3.23), and
sensor blindness analysis (§3.31) subsequently validate and refine this
conclusion.

\subsubsection{3.28 Forbidden Energy: The Hamiltonian of the State
Space}\label{forbidden-energy-the-hamiltonian-of-the-state-space}

The 100 forbidden states (unoccupied 8-bit hallmark configurations) are
not noise --- they are structural. The hallmark whose flip most often
leads to a forbidden state is Immort (41\% of occupied states have a
forbidden Immort-flip neighbor), followed by Genome (36\%), Immune
(18\%), and Metab (17\%). Prolif (8\%) and Invade (10\%) are almost
never forbidden --- the well-instrumented hallmarks can be freely
toggled without leaving the observed manifold.

\textbf{The topology phase transition.} Fitting a within-tier Cox model
of
\(E(\mathrm{state}) = \alpha \cdot \log(P) + \beta \cdot N_{\mathrm{forb}} + \gamma \cdot |H|\)
reveals that the forbidden-neighbor coefficient reverses between O2
regimes:

{\def\LTcaptype{none} % do not increment counter
\begin{longtable}[]{@{}
  >{\raggedright\arraybackslash}p{(\linewidth - 10\tabcolsep) * \real{0.1667}}
  >{\raggedright\arraybackslash}p{(\linewidth - 10\tabcolsep) * \real{0.1667}}
  >{\raggedright\arraybackslash}p{(\linewidth - 10\tabcolsep) * \real{0.1667}}
  >{\raggedright\arraybackslash}p{(\linewidth - 10\tabcolsep) * \real{0.1667}}
  >{\raggedright\arraybackslash}p{(\linewidth - 10\tabcolsep) * \real{0.1667}}
  >{\raggedright\arraybackslash}p{(\linewidth - 10\tabcolsep) * \real{0.1667}}@{}}
\toprule\noalign{}
\begin{minipage}[b]{\linewidth}\raggedright
O2 Tier
\end{minipage} & \begin{minipage}[b]{\linewidth}\raggedright
\(N\)
\end{minipage} & \begin{minipage}[b]{\linewidth}\raggedright
\(\beta(N_{\mathrm{forb}})\)
\end{minipage} & \begin{minipage}[b]{\linewidth}\raggedright
HR
\end{minipage} & \begin{minipage}[b]{\linewidth}\raggedright
\(p\)
\end{minipage} & \begin{minipage}[b]{\linewidth}\raggedright
Interpretation
\end{minipage} \\
\midrule\noalign{}
\endhead
\bottomrule\noalign{}
\endlastfoot
HIGH & 20,631 & \(+0.044\) & 1.045 & 0.042 & Weakly lethal \\
MEDIUM & 11,578 & \(-0.076\) & 0.927 & 0.002 & Protective \\
LOW & 16,466 & \(-0.141\) & 0.868 & \(1.5 \times 10^{-12}\) & Massively
protective \\
\end{longtable}
}

In LOW O2, a tumor sitting adjacent to forbidden states is
\emph{trapped} --- it cannot evolve toward those configurations because
the biological pathways are inaccessible. In hypoxic tissue, where ATP
is scarce and metabolic adaptation is essential, this topological
trapping is protective: the tumor is locked in a stable frustrated
equilibrium. In HIGH O2, where energy is abundant, the same trapping
limits adaptation --- a weakly unfavorable condition.

The per-cancer gradient confirms the phase boundary: Spearman
\(r(\mathrm{O2_{proxy}}, \beta_{\mathrm{topology}}) = +0.399\),
\(p = 0.020\). More oxygen makes forbidden neighbors more dangerous. The
single strongest per-cancer topology effect is DIFG (\(\beta = -0.421\),
\(p = 1.8 \times 10^{-54}\)), where the Immort/Metab axes are both
well-instrumented and highly variable.

\textbf{The energy function.} Combining log-prevalence (how common is
this state?), forbidden-neighbor count (how trapped is this state?), and
hallmark sum (how many capabilities?) into a single energy function
yields:

{\def\LTcaptype{none} % do not increment counter
\begin{longtable}[]{@{}llll@{}}
\toprule\noalign{}
Model & Parameters & \(C\) & AIC \\
\midrule\noalign{}
\endhead
\bottomrule\noalign{}
\endlastfoot
Hallmarks + O2 & 10 & 0.6214 & 403,112 \\
Energy function & 7 & 0.6067 & 403,696 \\
Combined & 15 & 0.6233 & 403,062 \\
\end{longtable}
}

The 7-parameter energy function captures 87.9\% of the 10-parameter
hallmark model's discriminability. The combined model (\(C = 0.6233\),
\(\Delta\mathrm{AIC} = -50\) vs.~hallmarks alone) shows that topology
carries independent information not redundant with individual hallmarks.
The dominant energy term is log-prevalence (\(\beta = +0.106\),
\(p = 10^{-86}\)): the most common tumor states are the most lethal ---
entropy is a hazard.

\textbf{Forbidden axis anatomy.} The probability that a patient's state
has a specific hallmark-flip leading to a forbidden state varies by O2:

{\def\LTcaptype{none} % do not increment counter
\begin{longtable}[]{@{}llll@{}}
\toprule\noalign{}
Hallmark & HIGH O2 & LOW O2 & Interpretation \\
\midrule\noalign{}
\endhead
\bottomrule\noalign{}
\endlastfoot
Immort & 64.2\% & 68.4\% & The dominant wall, O2-invariant \\
Genome & 8.9\% & 16.4\% & More constrained in LOW O2 \\
Immune & 1.4\% & 5.9\% & 4× more trapped in LOW O2 \\
Metab & 1.2\% & 3.9\% & 3× more trapped in LOW O2 \\
Death & 1.3\% & 4.2\% & 3× more trapped in LOW O2 \\
Prolif & 5.5\% & 5.4\% & O2-invariant \\
\end{longtable}
}

The blind hallmarks (Immune, Metab) and Death are \emph{more frequently
forbidden} in LOW O2 --- exactly the axes where the sensor panel is
weakest. The forbidden topology concentrates around the hallmarks the
panel cannot see, and that concentration increases in hypoxic tissue.
The holes in the hypercube are the thermodynamic anchors that stabilize
hypoxic tumors.

\begin{center}\rule{0.5\linewidth}{0.5pt}\end{center}

\emph{Every claim above is subjected to out-of-sample testing,
permutation analysis, and systematic sensor-blindness audits.}

\subsubsection{3.29 Cross-Validation, Bootstrap Confidence Intervals,
and Permutation
Tests}\label{cross-validation-bootstrap-confidence-intervals-and-permutation-tests}

The concordance values reported through §3.27 are point estimates on the
full dataset. Before any claim can be considered validated, the models
require out-of-sample testing. We performed five complementary
validation analyses on 48,675 patients with survival data and O2
classification.

\textbf{5-fold cross-validation.} We trained the 3-score + O2 Cox model
on 80\% of patients and evaluated concordance on the held-out 20\%,
repeating across 5 folds with stratification by cancer type. The
cross-validated concordance is \(C_\mathrm{CV} = 0.6101 \pm 0.0055\)
(mean \(\pm\) SD across folds), compared to the full-dataset estimate of
\(C = 0.6105\). The overfitting gap is \(+0.0004\) --- less than 0.1\%
of the total signal. The 3-score model generalizes without meaningful
overfit.

\textbf{Bootstrap 95\% confidence intervals.} We drew 1,000 bootstrap
resamples (stratified by cancer type) and computed concordance for each
model variant:

{\def\LTcaptype{none} % do not increment counter
\begin{longtable}[]{@{}lll@{}}
\toprule\noalign{}
Model & Bootstrap 95\% CI & Width \\
\midrule\noalign{}
\endhead
\bottomrule\noalign{}
\endlastfoot
8 hallmarks (flat) & {[}0.5634, 0.5717{]} & 0.0083 \\
3-score + O2 & {[}0.6062, 0.6145{]} & 0.0083 \\
8 hallmarks + O2 & {[}0.6171, 0.6259{]} & 0.0088 \\
\end{longtable}
}

All confidence intervals are non-overlapping: the 3-score model
significantly outperforms the flat 8-hallmark model, and the 8-hallmark
+ O2 model significantly outperforms the 3-score model. The O2
classification contributes a genuine, reproducible increment.

\textbf{Permutation test for O2 classification.} We shuffled the O2 tier
labels 1,000 times and recomputed the full model concordance under each
permutation. The null distribution is
\(C_\mathrm{null} = 0.5481 \pm 0.0007\), tightly concentrated around the
O2-ablated baseline. The observed concordance \(C = 0.6101\) yields
\(Z = 85.73\), \(p < 10^{-1000}\) (effectively zero). The O2
classification is the single most validated feature in the model --- its
removal collapses concordance by \(\Delta C = -0.062\), six times larger
than any individual hallmark's contribution.

\textbf{Per-cancer validation of the 3D manifold.} We tested whether the
3-score model outperforms the full 8-hallmark model \emph{within}
individual cancer types by computing cross-validated concordance for
each of 29 cancer types with \(\geq 50\) events. The 3-score model wins
in only 1 of 29 cancer types (gastric, \(\Delta C = +0.0017\)). In the
remaining 28, the full 8-hallmark model is equal or superior. The 3D
manifold is a pan-cancer compression: it captures the dominant
inter-tissue variation (the O2 axis) but does not improve within-tissue
prediction, where the additional 5 hallmark dimensions carry
tissue-specific information.

\textbf{The complete validated model hierarchy.} Cross-validated
concordance across all model tiers:

{\def\LTcaptype{none} % do not increment counter
\begin{longtable}[]{@{}llll@{}}
\toprule\noalign{}
Model & Parameters & \(C_\mathrm{CV}\) & \(\Delta\) vs.~3-score \\
\midrule\noalign{}
\endhead
\bottomrule\noalign{}
\endlastfoot
3-score + O2 & 5 & 0.6204 & --- \\
Energy function (topology) & 7 & 0.6067 & −0.0137 \\
8 hallmarks + O2 (flat) & 10 & 0.6214 & +0.0010 \\
Wall + O2 & 12 & 0.6241 & +0.0037 \\
Energy + hallmarks & 15 & 0.6233 & +0.0029 \\
Energy + visibility interaction & 11 & 0.6178 & −0.0026 \\
8 hallmarks + O2 + interactions & 26 & 0.6335 & +0.0131 \\
Wall + O2 + interactions & 23 & 0.6362 & +0.0158 \\
\end{longtable}
}

All models generalize with zero meaningful overfit (train--test gaps
\(< 0.002\) across all tiers). The wall encoding (§3.23) and interaction
terms both contribute validated increments. The full model hierarchy is
real, not an artifact of fitting.

\subsubsection{3.30 Strong Inference Battery: Falsifiable Tests for
O2-Hallmark
Physics}\label{strong-inference-battery-falsifiable-tests-for-o2-hallmark-physics}

To subject the O2-hallmark routing framework to formal falsification, we
constructed a battery of 8 hypotheses following the strong inference
methodology (Platt, 1964). Each hypothesis has a pre-registered pass
criterion. Results:

{\def\LTcaptype{none} % do not increment counter
\begin{longtable}[]{@{}
  >{\raggedright\arraybackslash}p{(\linewidth - 8\tabcolsep) * \real{0.2000}}
  >{\raggedright\arraybackslash}p{(\linewidth - 8\tabcolsep) * \real{0.2000}}
  >{\raggedright\arraybackslash}p{(\linewidth - 8\tabcolsep) * \real{0.2000}}
  >{\raggedright\arraybackslash}p{(\linewidth - 8\tabcolsep) * \real{0.2000}}
  >{\raggedright\arraybackslash}p{(\linewidth - 8\tabcolsep) * \real{0.2000}}@{}}
\toprule\noalign{}
\begin{minipage}[b]{\linewidth}\raggedright
\#
\end{minipage} & \begin{minipage}[b]{\linewidth}\raggedright
Hypothesis
\end{minipage} & \begin{minipage}[b]{\linewidth}\raggedright
Test
\end{minipage} & \begin{minipage}[b]{\linewidth}\raggedright
Result
\end{minipage} & \begin{minipage}[b]{\linewidth}\raggedright
Key Statistic
\end{minipage} \\
\midrule\noalign{}
\endhead
\bottomrule\noalign{}
\endlastfoot
H1 & O2 Discovery & K-means on per-cancer coefficient profiles recovers
O2 tiers & FAIL & ARI = 0.048, \(p = 0.10\) \\
H2 & Frustration Gradient & Prolif\_only β: LOW \textgreater{} MED
\textgreater{} HIGH (monotonic) & PASS & Interaction
\(p = 7 \times 10^{-76}\) \\
H3 & Sign Reversal & Prolif × Invade protective in MEDIUM O2 & PASS &
\(\beta = -0.486\), \(p = 1.1 \times 10^{-12}\) \\
H4 & Transfer Degradation & HIGH O2 model degrades on LOW O2 & FAIL &
\(\Delta C = +0.0004\) (coefficients shift but C insensitive) \\
H5 & DIFG Calibration & Shadow method recovers Metab in glioma & PASS &
\(\Delta C = +0.1045\) \\
H6 & Shadow Detection & Hidden variable in non-DIFG LOW O2 & PASS &
\(\beta_{\mathrm{Prolif\_only}} = +0.618\), \(p = 10^{-129}\) \\
H7 & Negative Control & Blind hallmarks add nothing in HIGH O2 & PASS &
\(\Delta C = +0.0004\), all \(p > 0.4\) \\
H8 & Forbidden Topology & Forbidden neighbors predict survival & FAIL
(pan-cancer) & Pan-cancer \(p = 0.28\) \\
\end{longtable}
}

\textbf{Score: 5/8 PASS.}

\textbf{Informative failures.} H1 fails because the O2 physics manifests
through \emph{interactions} (wall features, O2-gated coupling), not
through raw coefficient magnitudes --- k-means on 5-dimensional
coefficient vectors cannot capture interaction structure. H4 fails
because the concordance index is a rank statistic insensitive to
coefficient magnitude and sign: the HIGH→LOW transfer shows Prolif
shifts \(+0.43\), Invade \emph{flips sign}, and Genome \emph{flips
sign}, but the aggregate ranking stays similar. H8 fails pan-cancer
because the topology effect reverses between O2 tiers --- the pan-cancer
average washes out the signal (see §3.28).

\textbf{The negative control (H7) is definitive.} Adding all three blind
hallmarks (Immort, Immune, Metab) to the visible model in HIGH O2 adds
\(\Delta C = +0.0004\) --- indistinguishable from noise. Each
individually: Immort \(p = 0.45\), Immune \(p = 0.41\), Metab
\(p = 0.63\). The shadow signal is O2-specific, not a universal
confounder.

\textbf{The DIFG calibration (H5) validates the shadow method.} In
glioma, where the Metab hallmark is well-instrumented (IDH1 prevalence
50.2\%), adding Metab to the visible model increases concordance by
\(+0.1045\) --- the largest single-variable increment in the entire
study. Within visible-risk terciles, Metab = 1 patients outlive Metab =
0 patients by 16--35 months. This proves the shadow method works where
we can verify it.

\subsubsection{3.31 The Censored Manifold: Gene Panel Limitations and
Sensor
Blindness}\label{the-censored-manifold-gene-panel-limitations-and-sensor-blindness}

The 45-gene panel (Appendix B) is the measurement instrument through
which the entire analysis is conducted. Any limitation of this panel is
a limitation of every finding in this paper. We systematically
quantified where the panel is blind.

\textbf{Sensor quality by hallmark.} We define sensor quality as the
product of gene count and mean population-level mutation frequency for
each hallmark:

{\def\LTcaptype{none} % do not increment counter
\begin{longtable}[]{@{}llll@{}}
\toprule\noalign{}
Hallmark & Genes & Quality Score & Grade \\
\midrule\noalign{}
\endhead
\bottomrule\noalign{}
\endlastfoot
GrowSup & 16 & 13.98 & Excellent \\
Genome & 11 & 10.53 & Excellent \\
Prolif & 9 & 8.97 & Good \\
Death & 5 & 3.02 & Adequate \\
Invade & 4 & 2.01 & Thin \\
Metab & 2 & 0.67 & Poor \\
Immort & 1 (ATRX) & 0.33 & Blind \\
Immune & 1 (JAK1) & 0.18 & Blind \\
\end{longtable}
}

Three hallmarks are measured through \(\leq 2\) genes. Replicative
Immortality is detected exclusively through ATRX mutations (prevalence
5.5\%); true immortalization via TERT promoter mutations (\$\sim\$70\%
of melanomas, \$\sim\$80\% of glioblastomas) is invisible to standard
exon-sequencing panels. Immune Evasion is detected exclusively through
JAK1 mutations (prevalence 2\%); the dominant immune escape mechanisms
(B2M loss, HLA downregulation, PD-L1 amplification) are uncaptured.

\textbf{The frozen hallmark map.} Many cancer types have hallmarks
effectively frozen --- either constitutively ON (\(> 50\%\)) or
constitutively OFF (\(< 2\%\)) --- providing no discriminative
information:

\begin{itemize}
\tightlist
\item
  Immort is OFF (\(< 2\%\)) in 21 of 37 cancer types
\item
  Immune is OFF (\(< 2\%\)) in 27 of 37 cancer types
\item
  Metab is OFF (\(< 2\%\)) in 18 of 37 cancer types
\end{itemize}

At the extreme, high-grade serous ovarian carcinoma (HGSOC) has only 5
of 8 hallmarks with any variation --- the other 3 are frozen,
compressing its effective state space from 256 to 32 possible states (25
observed).

\textbf{State space compression.} Only 156 of 256 possible 8-bit states
are occupied across 198,862 tumors. The single most common state,
\(01100000\) (TP53 bond: GrowSup + Death), contains 41,198 patients
(20.7\%). The unoccupied states are concentrated where blind hallmarks
(Immort, Immune, Metab) are required to be active --- these states may
exist in nature but are invisible to our panel.

\textbf{The censored manifold test.} To measure how much the blind
hallmarks contribute to prediction, we defined a Wall-7 model using only
the 7 well-instrumented hallmarks (excluding Immort, Immune, and Metab
wall features) and compared it to the full Wall-10 model:

{\def\LTcaptype{none} % do not increment counter
\begin{longtable}[]{@{}
  >{\raggedright\arraybackslash}p{(\linewidth - 6\tabcolsep) * \real{0.1304}}
  >{\raggedright\arraybackslash}p{(\linewidth - 6\tabcolsep) * \real{0.3478}}
  >{\raggedright\arraybackslash}p{(\linewidth - 6\tabcolsep) * \real{0.3478}}
  >{\raggedright\arraybackslash}p{(\linewidth - 6\tabcolsep) * \real{0.1739}}@{}}
\toprule\noalign{}
\begin{minipage}[b]{\linewidth}\raggedright
O2 Tier
\end{minipage} & \begin{minipage}[b]{\linewidth}\raggedright
Wall-10 \(C_\mathrm{CV}\)
\end{minipage} & \begin{minipage}[b]{\linewidth}\raggedright
Wall-7 \(C_\mathrm{CV}\)
\end{minipage} & \begin{minipage}[b]{\linewidth}\raggedright
\(\Delta C\)
\end{minipage} \\
\midrule\noalign{}
\endhead
\bottomrule\noalign{}
\endlastfoot
HIGH & 0.5723 & 0.5732 & \(-0.0009\) \\
LOW & 0.6282 & 0.6029 & \(+0.0253\) \\
\end{longtable}
}

In HIGH O2 tissues, the blind hallmarks contribute \emph{nothing} ---
Wall-7 matches Wall-10. In LOW O2 tissues, the blind hallmarks add
\(\Delta C = +0.025\) --- and this entire increment is DIFG-driven (as
established in §3.23, the Hypoxic\_coupled state is \(> 24\%\) in DIFG,
\(< 0.5\%\) everywhere else). The blind hallmarks are carried by
glioma's IDH-mutant biology, not by a pan-cancer hypoxic signal.

\textbf{The Metabolic sensor gap.} Among all 8 hallmarks, Deregulated
Metabolism (2 genes: IDH1, GNAS) shows the strongest evidence that the
\emph{physics is real but the panel can't reach it}. Across 37 cancer
types, the correlation between sensor entropy (information content of
the hallmark's binary measurement) and the absolute Cox coefficient
\(|\beta|\) is:

{\def\LTcaptype{none} % do not increment counter
\begin{longtable}[]{@{}lll@{}}
\toprule\noalign{}
Hallmark & \(r(\mathrm{entropy}, |\beta|)\) & Interpretation \\
\midrule\noalign{}
\endhead
\bottomrule\noalign{}
\endlastfoot
Metab & \(+0.507\) & Higher entropy \(\to\) stronger signal \\
Prolif & \(-0.129\) & No relationship \\
GrowSup & \(+0.041\) & No relationship \\
All others & \(|r| < 0.2\) & No relationship \\
\end{longtable}
}

Metab is the only hallmark whose survival signal \emph{increases} when
the sensor provides more information (higher entropy = more variation,
meaning the hallmark is not frozen). In cancer types where Metab is not
frozen --- where IDH1 or GNAS actually varies --- the Cox coefficient is
large. Where Metab is frozen OFF (as in most aerobic cancers), the
coefficient is near zero not because metabolism doesn't matter, but
because the panel cannot detect the relevant metabolic alterations. With
better Metab sensors (e.g., HIF1A, VHL, TERT, pyruvate kinase isoform
switching), the 3D manifold's third dimension --- currently dominated by
the Genome toggle --- may expand to a genuine metabolic axis.

\textbf{Implications for the 3D manifold.} The sensor blindness results
explain two otherwise puzzling findings: (1) why the Hypoxic equilibrium
state (§3.23) is DIFG-specific --- the panel only captures glioma's
hypoxic adaptations (IDH1, ATRX) while being blind to non-glioma hypoxic
mechanisms (HIF pathway, TERT promoter); and (2) why the 3D manifold
captures 96.3\% of predictive power --- the remaining 3.7\% may be
dominated by noise from poorly-instrumented hallmarks contributing
variance without signal.

\subsubsection{3.32 Conditional Topology: The Second Axis
Adjudication}\label{conditional-topology-the-second-axis-adjudication}

The topology phase transition (§3.28) hides a deeper structure. The
forbidden-neighbor coefficient does not merely vary by O2 --- it
\emph{reverses sign} within LOW O2 based on whether the patient's tumor
has activated the Hypoxic axis (Immort=1 or Metab=1). This second axis
was independently proposed by two AI systems and subjected to a formal
adjudication with 8 tests.

\textbf{The visibility split.} In LOW O2, splitting patients by Hypoxic
axis activation reveals opposite topology effects:

{\def\LTcaptype{none} % do not increment counter
\begin{longtable}[]{@{}
  >{\raggedright\arraybackslash}p{(\linewidth - 10\tabcolsep) * \real{0.1667}}
  >{\raggedright\arraybackslash}p{(\linewidth - 10\tabcolsep) * \real{0.1667}}
  >{\raggedright\arraybackslash}p{(\linewidth - 10\tabcolsep) * \real{0.1667}}
  >{\raggedright\arraybackslash}p{(\linewidth - 10\tabcolsep) * \real{0.1667}}
  >{\raggedright\arraybackslash}p{(\linewidth - 10\tabcolsep) * \real{0.1667}}
  >{\raggedright\arraybackslash}p{(\linewidth - 10\tabcolsep) * \real{0.1667}}@{}}
\toprule\noalign{}
\begin{minipage}[b]{\linewidth}\raggedright
Group
\end{minipage} & \begin{minipage}[b]{\linewidth}\raggedright
\(N\)
\end{minipage} & \begin{minipage}[b]{\linewidth}\raggedright
Deaths
\end{minipage} & \begin{minipage}[b]{\linewidth}\raggedright
\(\beta(N_{\mathrm{forb}})\)
\end{minipage} & \begin{minipage}[b]{\linewidth}\raggedright
HR
\end{minipage} & \begin{minipage}[b]{\linewidth}\raggedright
\(p\)
\end{minipage} \\
\midrule\noalign{}
\endhead
\bottomrule\noalign{}
\endlastfoot
LOW O2, Hyp VISIBLE & 3,187 & 1,044 & \(-0.285\) & 0.752 &
\(1.5 \times 10^{-11}\) \\
LOW O2, Hyp INVISIBLE & 13,279 & 7,203 & \(+0.217\) & 1.243 &
\(4.8 \times 10^{-17}\) \\
\end{longtable}
}

The formal interaction term is
\(\beta_{\mathrm{N_{forb} \times vis}} = -0.213\),
\(p = 1.2 \times 10^{-7}\). After controlling for log-prevalence
(Gemini's confound), the interaction survives: \(\beta = -0.203\),
\(p = 6.7 \times 10^{-7}\).

\textbf{The confound tests.} Two competing explanations were tested:

\emph{Hypothesis A (Claude.AI):} Escape hatch visibility is a
patient-level biological axis. A tumor with visible Hypoxic markers
(Immort=1 or Metab=1) can reach walls safely via the Hypoxic branch; a
tumor without them can only reach walls via the lethal Aerobic pathway.

\emph{Hypothesis B (Gemini):} The visibility split is a tissue-type
confound driven by DIFG (53.6\% Hypoxic visible) dominating the
``visible'' group. Claude measured the lethality of the \emph{path}, not
the \emph{wall}.

Results of 7 adjudication tests:

{\def\LTcaptype{none} % do not increment counter
\begin{longtable}[]{@{}
  >{\raggedright\arraybackslash}p{(\linewidth - 4\tabcolsep) * \real{0.3333}}
  >{\raggedright\arraybackslash}p{(\linewidth - 4\tabcolsep) * \real{0.3333}}
  >{\raggedright\arraybackslash}p{(\linewidth - 4\tabcolsep) * \real{0.3333}}@{}}
\toprule\noalign{}
\begin{minipage}[b]{\linewidth}\raggedright
Test
\end{minipage} & \begin{minipage}[b]{\linewidth}\raggedright
Result
\end{minipage} & \begin{minipage}[b]{\linewidth}\raggedright
Winner
\end{minipage} \\
\midrule\noalign{}
\endhead
\bottomrule\noalign{}
\endlastfoot
A1: Within-cancer sign reversal & 3/9 reversals & Gemini \\
A3: DIFG removed & Interaction dissolves (\(p = 0.173\)) & Gemini \\
A4: Aerobic frustration control & Survives full wall control
(\(p = 6.8 \times 10^{-8}\)) & Claude \\
A5: Hallmark path profile & Invisible wall-hitters have 36\%
Prolif\_only vs.~21\% & Gemini \\
A6: Cancer-type fixed effects & Survives (\(p = 2.4 \times 10^{-9}\)) &
Claude \\
A7: Stratified Cox (non-parametric) & Survives
(\(p = 1.3 \times 10^{-9}\)) & Claude \\
A8: Forbidden flip anatomy & 100\% of invisible trapped patients blocked
at Immort & Gemini \\
\end{longtable}
}

Score: Claude 3, Gemini 3 (A2 excluded --- Gemini's PAAD Aerobic path
prediction was not confirmed: trapped PAAD tumors are \emph{less}
Aerobic, 73.4\% vs.~83.9\%).

\textbf{Synthesis.} Both hypotheses are partially correct at different
levels of description:

\begin{enumerate}
\def\labelenumi{\arabic{enumi}.}
\item
  \emph{Statistically}, the visibility × topology interaction is a
  within-cancer patient-level effect. It survives stratified Cox
  regression with non-parametric baselines per cancer type
  (\(p = 1.3 \times 10^{-9}\)) and simultaneous control for wall
  features + log-prevalence (\(p = 2.4 \times 10^{-13}\)). The
  interaction adds \(\Delta\mathrm{AIC} = -129\) beyond the full wall
  model (Model C: 18 params, \(C = 0.6636\) with cancer FE).
\item
  \emph{Mechanistically}, Gemini's structural explanation is confirmed:
  100\% of invisible trapped patients have Immort as a forbidden flip.
  The Hypoxic escape route (flipping Immort ON) is blocked by the
  forbidden topology. Invisible wall-hitters are on the Aerobic path
  (Prolif\_only = 36.3\% vs.~20.8\% for visible wall-hitters). The
  interaction \emph{dissolves} without DIFG (\(p = 0.173\)), confirming
  that DIFG's unique biology (BBB-constrained Invade=0, visible
  IDH1/ATRX) generates most of the signal.
\item
  \emph{Practically}, the effect is real but DIFG-dominated. In tissues
  where the Hypoxic axis is visible (primarily DIFG), topological
  trapping is a measurable protective factor. In tissues where it is
  invisible (PAAD, pancreas, prostate), the topology signal is lethal
  --- not because walls kill, but because the sensor can only detect
  walls reached via dangerous paths.
\end{enumerate}

\textbf{The updated Hamiltonian.} The energy function becomes:

\[E(\mathrm{state}, \mathrm{O2}, \mathrm{vis}) = \alpha \cdot \log(P) + \beta(\mathrm{O2}, \mathrm{vis}) \cdot N_{\mathrm{forb}} + \gamma \cdot |H| + \delta(\mathrm{O2})\]

where \(\beta(\mathrm{O2}, \mathrm{vis})\) encodes the 2×3 grid: O2 tier
modulates magnitude; Hypoxic visibility modulates sign. The strongest
anchor is LOW O2 + VISIBLE (\(\beta_{\mathrm{eff}} = -0.234\), the
Glioma singularity); the weakest is HIGH O2 + INVISIBLE
(\(\beta_{\mathrm{eff}} = -0.026\), approximately null).

\begin{center}\rule{0.5\linewidth}{0.5pt}\end{center}

\subsection{4. Discussion}\label{discussion}

\subsubsection{4.1 Clinical Implications}\label{clinical-implications}

The Patient Compass is designed for clinical translation. Its measure
--- pair completeness --- is an integer from 0 to 4 that can be computed
from any standard mutation panel (Foundation Medicine, MSK-IMPACT, TCGA,
or even targeted panels covering the 45 genes listed above). It requires
no proprietary software, no machine learning model, and no reference to
external databases beyond the gene-to-hallmark mapping table.

For clinicians, the compass provides:

\begin{enumerate}
\def\labelenumi{\arabic{enumi}.}
\item
  \textbf{A language for the conversation.} Instead of presenting
  patients with a list of mutated genes, the clinician can say: ``Your
  tumor has acquired three of the eight cancer hallmarks. None of the
  four functional pairs are complete. That means your cancer has gaps in
  its capabilities --- and those gaps may be where treatment is most
  effective.''
\item
  \textbf{Cancer-type-specific context.} Because the survival direction
  varies by cancer type, the compass does not make a universal claim
  about ``good'' or ``bad'' patterns. Instead, it reports what was
  observed in the patient's specific cancer type: ``In studies of your
  cancer type, patients with a similar balance pattern lived
  approximately X months longer/shorter.''
\item
  \textbf{A basis for questions.} Even when the compass does not find a
  statistically significant signal for a given cancer type, it provides
  the patient with informed questions: ``Has my tumor been profiled for
  these hallmarks? Are there clinical trials targeting the specific
  hallmarks my tumor has acquired?''
\end{enumerate}

\subsubsection{4.2 Lopsidedness, the Boltzmann Law, and the
Thermodynamic
Gradient}\label{lopsidedness-the-boltzmann-law-and-the-thermodynamic-gradient}

The most striking structural result is that 85.8\% of tumors have zero
complete pairs. The Boltzmann fit (\(P(k) \propto e^{-1.63 k}\),
\(R^2 = 0.9935\)) puts this observation on quantitative footing: each
additional complete pair costs approximately 1.63 nats of selective free
energy. The lopsidedness of cancer is not merely an observation --- it
is a thermodynamic law governing how tumors distribute across hallmark
space.

The developmental interpretation is strengthened by the pair hierarchy.
Cancer progression follows pathways of least resistance, but those
pathways are not symmetric across pairs. Pair 4 (Genome Instability +
Metabolism) is mechanistically coupled: genomic instability drives
metabolic reprogramming, and metabolic stress drives genomic instability
(the Warburg effect). A single molecular event --- such as an IDH1
mutation in glioma --- can complete Pair 4 in one step. In contrast,
Pair 2 (Growth Suppression + Immortality) requires two mechanistically
independent evolutionary innovations: inactivation of growth suppression
pathways (TP53/RB) \emph{and} activation of telomere maintenance
(ATRX/TERT). These capabilities share no known mechanistic link, making
Pair 2 the steepest gradient in the thermodynamic landscape.

This pair-specific difficulty hierarchy explains why 94 of 256 states
are never observed. The forbidden states are concentrated where the hard
pairs (Pair 2 and Pair 3) are required. The 162 observed states populate
the regions of hallmark space that are thermodynamically accessible ---
configurations reachable by the pathways of least resistance through the
pair landscape.

The cooperative effect --- transition ratios increasing from 0.131 to
0.293 as pairs accumulate --- reveals that this landscape tilts. The
first pair completion is the hardest (the ``nucleation barrier''), and
subsequent completions become progressively easier. This is consistent
with a model where completing one pair grants the tumor capabilities
that catalyze further acquisition: genomic instability (Pair 4)
increases the mutation rate, immune evasion (Pair 3) reduces selective
pressure against novel hallmarks, and proliferative signaling (Pair 1)
expands the population in which new mutations can arise.

Clinically, this thermodynamic picture identifies a natural intervention
point: the first pair completion. If pair completion is cooperative,
then preventing or disrupting the first completion --- particularly of
the enabler Pair 4 --- may delay subsequent completions
disproportionately. For patients whose tumors are already partially
complete, the compass identifies which pairs remain incomplete and
therefore which biological vulnerabilities persist.

The therapeutic implication is direct: if most tumors are lopsided (zero
complete pairs), then the majority of patients have tumors with
significant biological blind spots. The compass can guide therapy
selection by identifying which capabilities the tumor has acquired,
which it has not, and --- through the pair difficulty hierarchy ---
which missing capabilities the tumor is most and least likely to acquire
next.

\subsubsection{4.3 The Four-Spin Ising
Model}\label{the-four-spin-ising-model}

The pair interaction matrix (§3.11) completes the statistical mechanical
picture: the 198,862-tumor dataset is a Monte Carlo sample from a
four-spin Ising model at inverse temperature \(\beta = 1.63\). Each
tumor is a microstate, each pair is a spin, and the energy landscape is
fully characterized by 10 parameters: four field energies and six
couplings.

This identification is not merely an analogy. The Boltzmann fit to the
macrostate distribution yields \(R^2 = 0.9935\) --- the data \emph{are}
in thermal equilibrium at this temperature. The coupling matrix is
directly measured. Therefore, the system admits exact predictions.

\textbf{The dominant eigenvalue.} The coupling matrix \(J\) has a
leading eigenvalue \(\lambda_1 = 4.70\), accounting for 72\% of the
interaction variance. This means the system is not four independent
pairs with six separate interactions --- it is a single collective mode
with perturbative corrections. The dominant eigenvector points toward
simultaneous completion of all pairs, consistent with the cooperative
effect measured in §3.10. The spectral gap of 5.40 between \(\lambda_1\)
and \(\lambda_2 = -0.69\) implies that this collective mode is robust:
perturbations to individual couplings do not easily disrupt the
cooperative structure.

\textbf{The phase boundary.} The product \(\beta \cdot \alpha_c^2\),
where \(\alpha_c = \sqrt{135/88} \approx 1.239\) is the spectral
critical point identified in companion theoretical work {[}7--9{]},
evaluates to:

\[\beta \cdot \alpha_c^2 = 1.628 \times 1.534 = 2.498 \approx \frac{5}{2} \tag{9}\]

The deviation from 5/2 is 0.08\%. This suggests that the thermodynamic
temperature of tumor evolution is locked to the spectral critical point
--- tumors do not evolve at arbitrary temperatures, but at the
temperature determined by the underlying symmetry. If confirmed, this
would mean that the β measured empirically in any sufficiently large
cancer cohort must satisfy \(\beta = 5/(2\alpha_c^2)\), a prediction
that is testable with independent data.

\textbf{The enabler--barrier axis.} The strongest coupling
(\(J_{24} = +2.29\)) connects the enabler (Pair 4) to the barrier (Pair
2). The weakest coupling (\(J_{12} = +1.05\)) connects the neutral pair
(Pair 1) to the barrier. This hierarchy quantifies what the forbidden
state analysis (§3.5) reveals qualitatively: Pair 4 is not merely the
easiest pair to complete --- it is the pair whose completion most
dramatically opens access to the otherwise forbidden regions of hallmark
space. The coupling matrix gives this observation a number.

\textbf{Relation to the IDH1 bifurcation.} The glioma phase transition
(§3.12) provides the strongest test of the Ising model. In diffuse
glioma, IDH1 mutation completes Pair 4 at negligible evolutionary cost,
and the strong \(J_{24}\) coupling immediately pulls Pair 2 across the
accessibility threshold --- producing the bimodal k-distribution (69.3\%
at \(k = 0\), 22.2\% at \(k = 2\), depleted at \(k = 1\)) that is the
thermodynamic signature of a first-order phase transition. The dominant
macrostate (0,1,0,1) --- Pair 2 + Pair 4 complete, all others incomplete
--- contains 937 patients and corresponds exactly to the IDH-mutant
astrocytoma molecular subtype, identified here without any pathology
labels.

The survival advantage of Pair 4 completion in glioma (\(+36.4\) months,
\(z = +20.66\), \(p = 7.8 \times 10^{-95}\)) is the most powerful single
signal in the dataset. In breast cancer, the same Pair 4 completion
produces the opposite survival effect (\(-18.1\) months, \(p = 0.028\)).
This confirms that the survival direction is determined not by
temperature alone, but by which specific pairs are completed and the
tissue-specific biology of those completions --- the full Ising model,
not merely the scalar \(\beta\).

\textbf{Per-cancer temperature variation.} The effective temperature
varies three-fold across cancer types (§3.13), from \(\beta = 2.94\)
(clear cell renal, coldest) to \(\beta = 1.04\) (uterine endometrial,
hottest). Glioma (\(\beta = 1.32\)) is among the hottest --- not because
glioma is inherently aggressive, but because the IDH1 pathway provides a
low-energy route to pair completion unavailable in colder cancers. The
weak correlation between \(\beta\) and survival direction
(\(r = -0.13\), \(p = 0.49\)) establishes that the direction split
(§3.8) depends on the coupling structure \(J_{ij}\), not the
temperature. The \(\beta \cdot \alpha_c^2 = 5/2\) relation holds for the
global mixture; whether individual cancer types respect this constraint
or have their own critical points is an open question (§4.7).

\textbf{Emergent universality.} The per-cancer coupling analysis (§3.14)
resolves this open question decisively: the three constants do
\emph{not} hold for individual cancer types. Per-cancer
\(\beta \cdot \alpha_c^2\) ranges from 1.3 to 4.5, and no individual
cancer satisfies all three relations simultaneously. This is the
strongest result of the per-cancer analysis: the constants are emergent
properties of the mixture, not intrinsic to individual cancer types.
Each cancer is a thermodynamic subsystem with its own temperature and
coupling matrix; the apparent universal constants arise from the
weighted superposition of these 37 subsystems in the proportion they
appear in the cBioPortal database. This does not diminish the constants
--- it clarifies their ontological status. They characterize the
macroscopic epidemic of cancer, not a single disease. If the composition
of the database changed (more glioma, less breast), the global \(\beta\)
would shift and the constants would change. The constraint that the
global mixture obeys \(\beta \cdot \alpha_c^2 = 5/2\) therefore becomes
a statement about the \emph{equilibrium} of cancer types in the human
population --- a deeper claim than tumor physics alone.

\textbf{Evolutionary time.} The ferromagnetic coupling structure (all
\(J_{ij} > 0\)) implies that pair completions are positively correlated
\emph{across evolutionary time}. A tumor that has completed one pair is
more likely to complete the next --- not merely because of capability
catalysis (the nucleation mechanism of §3.10), but because the
underlying selection pressures of preceding hallmarks and elapsed tumor
evolution time are themselves correlated. The pair interaction matrix
gives the strength of this correlation. The leading eigenvalue gives the
timescale of the dominant cooperative mode.

\subsubsection{4.4 Testable Predictions}\label{testable-predictions}

The Ising model framework generates specific, falsifiable predictions
beyond those already confirmed:

\begin{enumerate}
\def\labelenumi{\arabic{enumi}.}
\item
  \textbf{Universal temperature.} If \(\beta \cdot \alpha_c^2 = 5/2\) is
  exact, then any sufficiently large mutation dataset
  (\textgreater10,000 patients) with hallmark pair annotations should
  yield \(\beta = 1.628 \pm \delta\), where \(\delta\) decreases as
  \(1/\sqrt{N}\). This can be tested independently using TCGA
  whole-exome data, the AACR GENIE dataset, or the UK Biobank cancer
  cohort.
\item
  \textbf{Per-cancer effective temperature and survival direction.}
  Individual cancer types have effective temperatures ranging three-fold
  (\(\beta = 1.04\) to \(2.94\), §3.13), but temperature alone does not
  predict survival direction (\(r = -0.13\), \(p = 0.49\)). Exhaustive
  testing of 37 features (§3.15) identifies a single surviving candidate
  at the cancer-type level: the population-eigenvector misalignment
  \(\cos(\hat{\mathbf{P}}, \hat{\mathbf{v}}_1)\), which explains
  14--19\% of variance but does not survive multiple comparison
  correction or decomposition to the study level (§3.16). Study-level
  decomposition and Simpson's paradox analysis (§3.18) reveal that 43\%
  of cancer types show direction flips between study and aggregate
  levels. We predict that the survival direction is study-specific
  rather than cancer-type-specific --- determined by the full
  \(4 \times 4\) interaction in each patient cohort, including treatment
  context --- and that domain wall density (§3.21) will prove a more
  robust patient-level predictor than any cancer-type-level scalar.
\item
  \textbf{Coupling matrix stability.} The six coupling constants
  \(J_{ij}\) should be stable across independent datasets. If the Ising
  model captures genuine biology rather than gene-panel artifact, the
  rank ordering \(J_{24} > J_{34} > J_{13} > J_{23} > J_{14} > J_{12}\)
  should reproduce in any cohort with \textgreater50,000 patients and
  comparable gene coverage.
\item
  \textbf{Longitudinal pair dynamics.} In matched primary--metastatic
  pairs, the Ising model predicts that pair completeness should increase
  (never decrease) from primary to metastasis, with the rate of increase
  proportional to the coupling constants. Specifically, transitions from
  \(\sigma_i = 0\) to \(\sigma_i = 1\) should be most frequent for Pair
  4 (strongest field + strongest couplings) and rarest for Pair 2
  (weakest field + weakest coupling to Pair 1). The MSK MetTropism
  dataset contains exactly such matched pairs.
\end{enumerate}

\subsubsection{4.5 Limitations}\label{limitations}

\textbf{Concordance is modest.} The best pan-cancer concordance achieved
is \(C_\mathrm{CV} = 0.6362\) (wall + O2 + interactions, 23 parameters)
and \(C_\mathrm{CV} = 0.6101\) (3-score model, 5 parameters). Clinical
decision tools typically require \(C > 0.65\)--\(0.70\) for actionable
prognostication. Glioma alone reaches \(C = 0.756\), but most individual
cancer types fall below \(C = 0.60\). The compass reveals
\emph{biological structure} in survival data; it does not yet achieve
the predictive power required for clinical deployment. The tool is a
research calculator, not a clinical instrument.

\textbf{No clinical covariates.} The model contains no age, stage,
grade, or treatment data. These are the variables that drive clinical
decisions. If adding age and stage to the 3-score model raises
concordance to \(C > 0.65\), the hallmark framework is contributing
incrementally. If age and stage alone reach \(C = 0.65\) without
hallmarks, our model is a proxy for known clinical variables. We cannot
distinguish these scenarios without linked clinical data. TCGA provides
age and stage for \textasciitilde29,000 patients; this is the most
important follow-up experiment.

\textbf{Gene panel is thin --- three hallmarks are severely
underdetected, and the manifold is censored.} The 45-gene panel detects
Immune Evasion through JAK1 alone (prevalence: 2\%) and Replicative
Immortality through ATRX alone (prevalence: 5.5\%). True immune evasion
prevalence is estimated at 15--30\% (via B2M, HLA loss, PD-L1
amplification) and true immortalization at 60--80\% (TERT promoter
mutations are present in \textasciitilde70\% of melanomas,
\textasciitilde80\% of glioblastomas, \textasciitilde65\% of bladder
cancers, and are not captured by standard exon-sequencing panels).
Deregulated Metabolism is measured by only 2 genes (IDH1, GNAS), with a
sensor quality score of 0.67 --- yet Metab is the only hallmark where
survival signal strength correlates with sensor entropy (\(r = +0.507\),
§3.31), proving the underlying biology is real but underdetected. Immort
is frozen OFF (\(< 2\%\)) in 21 of 37 cancer types; Immune is frozen in
27 of 37. The effective state space is compressed: only 156 of 256
states are occupied, and cancer types like HGSOC are limited to 25
observed states out of a theoretical 256. A Wall-7 model excluding blind
hallmarks matches the full model in HIGH O2 (\(\Delta C = -0.001\)) but
loses \(+0.025\) in LOW O2 --- entirely driven by glioma's IDH1-mutant
biology.

This underdetection distorts the pair structure, the Boltzmann fit, and
potentially the 3D manifold dimension. Specifically, the Hypoxic
equilibrium state (Immort + Genome + Metab wall features) is detectable
only in DIFG (prevalence 24.1\%) and undetectable in all other LOW O2
cancers (\(< 0.5\%\)). The ``DIFG-specificity'' of the Hypoxic signal
(§3.23) may be a sensor artifact rather than a biological fact:
non-glioma tumors likely achieve hypoxic adaptation through HIF pathway
activation, VHL loss, and TERT promoter mutations --- mechanisms
invisible to our panel. Future work should incorporate: - TERT promoter
mutations for Replicative Immortality - B2M, IFNGR1/2, and antigen
presentation genes for Immune Evasion - MSI-H status (computable from
existing MMR gene data but not currently scored) - Tumor mutational
burden (TMB) as a continuous Genome Instability proxy - Homologous
recombination deficiency (HRD) scores - Copy number alterations (MYC
amplification, CDKN2A deletion) - Gene expression signatures (EMT scores
for Invasion \& Metastasis) - HIF1A, VHL, and metabolic enzyme isoforms
for Deregulated Metabolism

\textbf{Mutation-only ascertainment.} We detect hallmark activation
exclusively through somatic mutations. Hallmarks can also be activated
by epigenetic silencing (e.g., MLH1 promoter methylation), copy number
changes, structural variants, and expression dysregulation. Our measure
therefore represents a \emph{lower bound} on hallmark activation, and
the true pair completeness rate is likely higher than observed.

\textbf{Tissue-level oxygenation is a proxy, not a measurement.} The O2
classification assigns patients to high, medium, or low oxygenation
based on tissue of origin (lung = high, pancreas = low). Actual tumor
hypoxia varies dramatically within a tissue type depending on
vascularity, tumor size, and distance from blood supply. When we report
``O2 is the strongest survival predictor'' (HR = 0.593), we are partly
reporting that tissue type is the strongest predictor --- which was
already established in §3.19. Patient-level hypoxia measurement (HIF-1α
expression, pimonidazole staining, oxygen electrode) would be needed to
validate the O2 axis as a true biological moderator rather than a
tissue-type surrogate.

\textbf{No treatment data.} The compass predicts prognosis but cannot
predict treatment response. A patient classified as ``high risk'' may be
high risk because their cancer is biologically aggressive, because they
received less effective treatment, or both. Without linked treatment
data, these explanations are confounded. The AACR GENIE dataset contains
treatment outcomes for \textasciitilde100,000 patients and would be the
ideal validation cohort.

\textbf{Cross-validation confirms generalization but does not resolve
clinical utility.} Five-fold cross-validation of the 3-score + O2 model
yields \(C_\mathrm{CV} = 0.6101 \pm 0.0055\) with overfitting
\(< 0.1\%\) (§3.29). Bootstrap confidence intervals are non-overlapping
between model tiers. Permutation testing confirms the O2 classification
at \(Z = 85.73\). However, the best validated concordance
(\(C_\mathrm{CV} = 0.6362\) for the wall + interactions model) remains
below the \(C > 0.65\)--\(0.70\) threshold typically required for
clinical decision tools. Cross-validation on a \emph{held-out external
dataset} (e.g., AACR GENIE) has not been performed and is required to
confirm generalization beyond cBioPortal studies.

\textbf{Retrospective analysis.} All data are retrospective and
observational. Survival associations may be confounded by treatment,
stage, age, and other clinical variables not controlled for. A
prospective study or a multivariate analysis controlling for known
prognostic factors would be needed to establish independent prognostic
value.

\textbf{No Bonferroni correction applied to the 181-study battery.} We
report nominal p-values. With 181 tests, a Bonferroni-corrected
threshold would be p \textless{} 0.05/181 = 2.76 × 10⁻⁴. Under this
correction, 8 studies remain significant (all glioma cohorts, MSK-CHORD,
Metastatic Colorectal, Cholangiocarcinoma, and Prostate Cancer MSK). The
direction-agnostic Stouffer meta-analysis, which is a single test,
remains overwhelmingly significant (p = 8.9 × 10⁻⁵³) and is not subject
to multiple testing correction.

\textbf{The biological pairs are suboptimal.} The
Hanahan-Weinberg-inspired pairing ranks 19th of 105 possible partitions
for survival separation (\(F = 72.9\) vs.~optimal \(F = 188.2\), §3.19).
The Ising model framework (§3.11) and thermodynamic constants
(\(\beta \cdot \alpha_c^2 \approx 5/2\)) are defined relative to this
specific pairing. Under the data-optimal pairing, the constants would
differ. The biological pairing captures genuine structure
(\(p = 8.9 \times 10^{-53}\)) but is not uniquely optimal.

\textbf{Mutual information cross-talk is small.} The Genome
\(\leftrightarrow\) Hypoxic MI difference between O2 regimes is 0.048
vs.~0.053 bits --- a difference of 0.005 bits. The Cox interaction is
highly significant (\(p = 7.2 \times 10^{-10}\)) but the actual
information content of the O2-gated coupling is tiny. Statistical
significance with \(n = 48{,}675\) does not imply large effect size.

\subsubsection{4.6 Comparison to Existing
Work}\label{comparison-to-existing-work}

Several existing tools stratify patients by mutation profiles, including
OncoKB (Chakravarty et al., 2017), cBioPortal itself (Cerami et al.,
2012; Gao et al., 2013), and various mutational signature approaches
(Alexandrov et al., 2013). The Patient Compass differs in three
respects:

\begin{enumerate}
\def\labelenumi{\arabic{enumi}.}
\tightlist
\item
  \textbf{Hallmark-level abstraction.} Rather than operating at the gene
  level, the compass operates at the hallmark level, providing a
  biologically interpretable summary that is robust to gene panel
  differences.
\item
  \textbf{Pair structure.} The compass does not merely count hallmarks
  (which would be an 8-dimensional binary vector) but imposes a specific
  pairing structure that captures functional tensions.
\item
  \textbf{Patient-facing design.} Existing tools are designed for
  bioinformaticians and clinicians. The Patient Compass is designed to
  be understood by patients themselves, with plain-language explanations
  and no required bioinformatics expertise.
\end{enumerate}

\subsubsection{4.7 Future Work}\label{future-work}

\begin{enumerate}
\def\labelenumi{\arabic{enumi}.}
\tightlist
\item
  \textbf{Expanded gene panel --- resolving the censored manifold.}
  Incorporate TERT promoter mutations, B2M and antigen presentation
  genes, HIF1A, VHL, copy number alterations, and gene expression
  signatures to improve hallmark detection sensitivity, particularly for
  Replicative Immortality (currently 1 gene), Immune Evasion (1 gene),
  and Deregulated Metabolism (2 genes). The sensor blindness analysis
  (§3.31) demonstrates that Metab is the hallmark most likely to gain
  predictive power from expanded panels (\(r = +0.507\) between sensor
  entropy and \(|\beta|\)). This may alter the pair difficulty
  hierarchy, increase observed pair completeness rates, and potentially
  expand the 3D manifold's third dimension from a Genome toggle to a
  full metabolic axis.
\item
  \textbf{Multivariate survival analysis.} Cox proportional hazards
  models controlling for age, stage, treatment, and known prognostic
  mutations to establish independent prognostic value of pair
  completeness.
\item
  \textbf{Per-cancer coupling matrix stability.} The 37 per-cancer
  coupling matrices reported here (§3.14) should be tested for stability
  across independent datasets. If the coupling topology --- which pairs
  are more strongly coupled in which cancer types --- reproduces in AACR
  GENIE or TCGA whole-exome data, then the coupling matrices are
  biological signatures, not gene-panel artifacts.
\item
  \textbf{Longitudinal pair dynamics via MSK MetTropism.} Matched
  primary--metastatic pairs to test whether \(\sigma_i\) transitions are
  unidirectional (\(0 \to 1\)) and proportional to the measured coupling
  constants. The Ising model predicts Pair 4 transitions should be most
  frequent, Pair 2 transitions rarest.
\item
  \textbf{Universal temperature test.} Replicate the
  \(\beta \cdot \alpha_c^2 = 5/2\) relation in independent datasets
  (AACR GENIE, UK Biobank cancer cohort, TCGA whole-exome) to determine
  whether this is a universal constant or an artifact of the cBioPortal
  mixture.
\item
  \textbf{Coupling matrix stability.} Verify that the rank ordering
  \(J_{24} > J_{34} > J_{13} > J_{23} > J_{14} > J_{12}\) reproduces in
  independent cohorts, establishing the biological reality (versus
  gene-panel artifact) of the interaction structure.
\item
  \textbf{Alternative pairings.} Systematically test all 105 possible
  4-pair partitions of the 8 hallmarks to determine whether the spectral
  pairing is uniquely optimal or one of several informative
  configurations.
\item
  \textbf{Integration with treatment data.} Determine whether pair
  completeness predicts response to specific therapy classes,
  particularly immunotherapy (where Pair 3 status may be directly
  relevant) and metabolic inhibitors (where Pair 4 status may predict
  sensitivity).
\item
  \textbf{Prospective validation.} Apply the compass prospectively to
  newly sequenced tumors and track outcomes.
\item
  \textbf{External validation of the frustrated proliferator (open
  problem).} The Prolif\_only signal (\(\beta = +0.594\) in LOW O2,
  \(p < 10^{-145}\); §3.23) and the invasion escape valve
  (\(\beta = -0.614\) in MEDIUM O2; §3.23) are the most actionable
  findings in this paper. We pose this as an open problem for groups
  with access to treatment-linked genomic data (e.g., AACR GENIE-BPC,
  Flatiron Health, institutional cohorts): does the O2-routed frustrated
  proliferator predict treatment response? The key testable predictions
  are (a) Prolif\_only patients in hypoxic cancers should respond poorly
  to anti-proliferative monotherapy but potentially benefit from
  anti-angiogenic combination (which would shift the O2 landscape), and
  (b) invasion-targeting agents should show differential efficacy by O2
  tier. Any dataset with mutation profiles, tissue type, treatment, and
  survival can replicate these tests using the gene-to-hallmark mapping
  and O2 classification provided in §2.
\item
  \textbf{Theoretical derivation.} Derive the coupling matrix \(J_{ij}\)
  from first principles using the spectral structure of the hallmark
  co-occurrence algebra {[}7--9{]}. If the empirically measured
  couplings can be predicted from eigenvalue analysis of the hallmark
  space, this would establish that the tumor's selective landscape is
  determined by symmetry rather than contingency.
\item
  \textbf{Refine the energy function with patient-level O2 measurement.}
  The forbidden topology phase transition (§3.28) --- where trapped
  states flip from protective to lethal between O2 regimes --- is
  currently measured using tissue-of-origin as an O2 proxy.
  Patient-level hypoxia measurements (HIF-1α expression, pimonidazole
  staining) would allow continuous estimation of the phase boundary and
  higher-resolution energy maps. The per-cancer gradient
  (\(r = +0.399\), \(p = 0.020\)) suggests the transition is smooth, not
  a sharp critical point.
\item
  \textbf{Compass prototype operationalization.} Convert the state-level
  energy map (§3.28) into a clinical tool that takes a patient's
  hallmark vector and tissue type and returns: (a) state energy (basin
  depth), (b) forbidden neighbor count (topological stability), (c) list
  of blocked evolutionary exits, (d) O2-contextualized clinical
  assessment, and (e) escape hatch visibility flag (§3.32). The
  proof-of-concept compass profiles in §3.28 demonstrate feasibility; a
  web-accessible implementation with confidence intervals requires full
  cross-validation of the energy function.
\item
  \textbf{Resolve the DIFG dominance of the visibility axis (open
  problem).} The escape hatch visibility effect (§3.32) survives
  stratified Cox regression (\(p = 10^{-9}\)) but dissolves without DIFG
  (\(p = 0.173\)). We pose this as an open problem: does the visibility
  × topology interaction replicate in non-DIFG hypoxic cancers when the
  gene panel is expanded? Adding HIF1A, VHL, and TERT to the hallmark
  map would increase Hypoxic visibility in PAAD/pancreas and test
  whether the effect generalizes beyond DIFG's unique biology
  (BBB-constrained invasion, visible IDH1/ATRX). Groups with access to
  treatment-linked cohorts or expanded panels are invited to test this
  directly.
\end{enumerate}

\begin{center}\rule{0.5\linewidth}{0.5pt}\end{center}

\subsection{5. Summary of Findings}\label{summary-of-findings}

{\def\LTcaptype{none} % do not increment counter
\begin{longtable}[]{@{}
  >{\raggedright\arraybackslash}p{(\linewidth - 6\tabcolsep) * \real{0.1034}}
  >{\raggedright\arraybackslash}p{(\linewidth - 6\tabcolsep) * \real{0.2759}}
  >{\raggedright\arraybackslash}p{(\linewidth - 6\tabcolsep) * \real{0.2759}}
  >{\raggedright\arraybackslash}p{(\linewidth - 6\tabcolsep) * \real{0.3448}}@{}}
\toprule\noalign{}
\begin{minipage}[b]{\linewidth}\raggedright
\#
\end{minipage} & \begin{minipage}[b]{\linewidth}\raggedright
Result
\end{minipage} & \begin{minipage}[b]{\linewidth}\raggedright
Status
\end{minipage} & \begin{minipage}[b]{\linewidth}\raggedright
Evidence
\end{minipage} \\
\midrule\noalign{}
\endhead
\bottomrule\noalign{}
\endlastfoot
1 & Most cancers are lopsided (85.8\% have zero complete pairs) &
\textbf{Verified} & §3.1; 170,622 / 198,862 tumors \\
2 & Pair completeness stratifies survival pan-cancer & \textbf{Verified}
& §3.6; Stouffer \(z = +15.29\), \(p = 8.9 \times 10^{-53}\) \\
3 & Survival direction is cancer-type-dependent & \textbf{Verified} &
§3.8; glioma protective, colorectal harmful \\
4 & 94 of 256 hallmark states are forbidden & \textbf{Verified} & §3.2;
top 10 states = 70.0\% of patients \\
5 & Pair completeness follows Boltzmann distribution & \textbf{Verified}
& §3.4; \(\beta = 1.63\), \(R^2 = 0.9935\) \\
6 & Pairs have asymmetric thermodynamic cost & \textbf{Verified} & §3.5;
Pair 4 enabler, Pair 2 barrier \\
7 & Pair completion is cooperative & \textbf{Verified} & §3.10;
transition ratio 0.131 → 0.293 \\
8 & Each pair individually predicts survival & \textbf{Verified} & §3.9;
all \(p < 10^{-6}\), +2.9 to +4.1 months \\
9 & Pairs interact as four-spin Ising model & \textbf{Verified} & §3.11;
\(J_{24} = +2.29\), \(\lambda_1\) captures 72\% \\
10 & Glioma exhibits first-order phase transition & \textbf{Verified} &
§3.12; IDH1: +36.4 mo, \(p = 7.8 \times 10^{-95}\) \\
11 & Effective temperature varies 3-fold across cancers &
\textbf{Verified} & §3.13; \(\beta = 1.04\) (endometrial) to 2.94
(renal) \\
12 & 37 cancer types yield unique coupling fingerprints &
\textbf{Verified} & §3.14; no single feature predicts direction \\
13 & Universal constants are emergent, not intrinsic & \textbf{Verified}
& §3.15; per-cancer \(\beta \alpha_c^2\) ranges 1.3--4.5 \\
14 & Survival direction resists scalar prediction & \textbf{Verified} &
§3.16; best candidate \(R^2 = 0.14\)--\(0.19\), fails Bonferroni \\
15 & Domain walls are lethal & \textbf{Verified} & §3.21;
\(z = -19.34\), \(p = 2.7 \times 10^{-83}\) \\
16 & Second-order tensor is universally antagonistic & \textbf{Verified}
& §3.17; crude \(r = -0.996\) was small-sample artifact \\
17 & Simpson's paradox afflicts 43\% of cancer types & \textbf{Verified}
& §3.18; 12/28 direction flips between aggregation levels \\
18 & Biological pairs rank 19th of 105 for survival separation &
\textbf{Falsified} & §3.19; \(F = 72.9\) vs.~optimal \(F = 188.2\) \\
19 & Biological pairs are a human grating, not a thermodynamic
eigenvector & \textbf{Falsified} & §3.19; spectral gap 2.93, effective
dim ≈ 1.9 \\
20 & Master equation reveals two metastable basins & \textbf{Verified} &
Appendix C; Basin A (120,635 pts, high hazard) vs.~Basin B (78,227
pts) \\
21 & Tissue O₂ shapes hallmark acquisition; Genomic Rubicon inverts &
\textbf{Verified} & §3.24; Genome harmful in high-O₂ (\(p = 0.001\)),
protective in low-O₂ (\(p = 8.3 \times 10^{-38}\)) \\
22 & Eight hallmarks collapse to 3D survival manifold &
\textbf{Verified} & §3.25; Marchenko-Pastur: 3 eigenvalues above noise
in all O₂ regimes \\
23 & Genome Instability is an O₂-gated coupling term & \textbf{Verified}
& §3.26; low-O₂ \(\beta = +0.372\), \(p = 7.2 \times 10^{-10}\); high-O₂
\(p = 0.315\) \\
24 & Three independent tests converge on 3D manifold & \textbf{Verified}
& §3.27; eigenvalue + concordance compression + coupling analysis \\
25 & Cross-validation confirms zero meaningful overfit &
\textbf{Verified} & §3.29; \(C_\mathrm{CV} = 0.6101 \pm 0.0055\), gap
\(= +0.0004\) \\
26 & Frustrated proliferation is the strongest hallmark survival feature
& \textbf{Verified} & §3.23; HR \(= 1.81\) in low-O₂,
\(\Delta\mathrm{AIC} = -220\) vs.~flat model \\
27 & Gene panel is censored --- three hallmarks are effectively blind &
\textbf{Verified} & §3.31; Immort: 1 gene, Immune: 1 gene, Metab: 2
genes \\
28 & Strong inference battery: 5 of 8 falsifiable tests pass &
\textbf{Verified} & §3.30; frustrated proliferator gradient monotonic
across O₂ tiers \\
29 & Forbidden topology is the O₂-gated phase transition &
\textbf{Verified} & §3.28; protective in low-O₂ (\(p = 10^{-12}\)),
weakly lethal in high-O₂ \\
30 & 7-parameter energy function captures 87.9\% of discriminability &
\textbf{Verified} & §3.28; \(C = 0.6067\) with 30\% fewer parameters \\
31 & Escape hatch visibility modulates topology survival effect &
\textbf{Partially verified} & §3.32; stratified \(p = 10^{-9}\),
dissolves without DIFG (\(p = 0.173\)) \\
\end{longtable}
}

\emph{Status key:} \textbf{Verified} = survives cross-validation,
permutation, or leave-one-out testing. \textbf{Falsified} = tested and
rejected by the data. \textbf{Partially verified} = significant in full
cohort but driven by a single cancer type.

\begin{center}\rule{0.5\linewidth}{0.5pt}\end{center}

\subsection{6. Data Availability}\label{data-availability}

All mutation and survival data used in this study are publicly available
through the cBioPortal for Cancer Genomics (https://www.cbioportal.org).
The complete gene-to-hallmark mapping, pair definitions, analysis code,
and reproducibility script are available at
\href{https://github.com/fancyland-llc/patient-compass-pan-cancer}{github.com/fancyland-llc/patient-compass-pan-cancer}.
The Patient Compass interactive report is included in the repository and
contains no patient-identifiable information.

No patient data were collected, stored, or transmitted by this study.
All analyses were performed on anonymized, previously published datasets
accessed through cBioPortal's public API.

\begin{center}\rule{0.5\linewidth}{0.5pt}\end{center}

\subsection{7. References}\label{references}

\begin{enumerate}
\def\labelenumi{\arabic{enumi}.}
\tightlist
\item
  Hanahan, D., \& Weinberg, R. A. (2000). The Hallmarks of Cancer.
  \emph{Cell}, 100(1), 57--70.
\item
  Hanahan, D., \& Weinberg, R. A. (2011). Hallmarks of Cancer: The Next
  Generation. \emph{Cell}, 144(5), 646--674.
\item
  Cerami, E., et al.~(2012). The cBio Cancer Genomics Portal: An Open
  Platform for Exploring Multidimensional Cancer Genomics Data.
  \emph{Cancer Discovery}, 2(5), 401--404.
\item
  Gao, J., et al.~(2013). Integrative Analysis of Complex Cancer
  Genomics and Clinical Profiles Using the cBioPortal. \emph{Science
  Signaling}, 6(269), pl1.
\item
  Chakravarty, D., et al.~(2017). OncoKB: A Precision Oncology Knowledge
  Base. \emph{JCO Precision Oncology}, 1, 1--16.
\item
  Alexandrov, L. B., et al.~(2013). Signatures of mutational processes
  in human cancer. \emph{Nature}, 500, 415--421.
\item
  Matos, A. P. (2026a). Active Transport on the Prime Gas: Flat-Band
  Condensation, the Rabi Phase Transition, and the Arithmetic Qubit.
  \emph{Preprint}, DOI: 10.5281/zenodo.19243258.
\item
  Matos, A. P. (2026b). The Arithmetic Black Hole: Softmax
  Thermodynamics and the Four Eigenvalue Laws of the Prime Gas.
  \emph{Preprint}, DOI: 10.5281/zenodo.19442006.
\item
  Matos, A. P. (2026c). The Unity Clock: Effective Dimensional Collapse
  of the Addition-Multiplication Commutator on Coprime Lattices.
  \emph{Preprint}, DOI: 10.5281/zenodo.19478727.
\item
  Matos, A. P. (2026d). Fault-Injection-Immune Computational Unit Using
  Primorial Coprime Residue Topology. \emph{U.S. Provisional Patent
  Application No.~64/031,440}, filed April 7, 2026.
\item
  Matos, A. P. (2026e). Holographic Eigen-Solver Using QM Boundary
  Projection on Coprime Residue Lattices. \emph{U.S. Provisional Patent
  Application No.~64/033,689}, filed April 8, 2026.
\end{enumerate}

\subsection{Acknowledgments}\label{acknowledgments}

This work was conducted using the Lattice OS axiomatic research
protocol, in which the author proposed hypotheses and two large language
models --- Claude Opus 4 (Anthropic) and Gemini 3.1 Pro (Google
DeepMind) --- served as adversarial reviewers. All claims were resolved
by computational execution, not by model assertion.

\subsection{Appendix A: Complete Table of Significant Studies (p
\textless{}
0.05)}\label{appendix-a-complete-table-of-significant-studies-p-0.05}

\small

{\def\LTcaptype{none} % do not increment counter
\begin{longtable}[]{@{}
  >{\raggedright\arraybackslash}p{(\linewidth - 14\tabcolsep) * \real{0.1250}}
  >{\raggedright\arraybackslash}p{(\linewidth - 14\tabcolsep) * \real{0.1250}}
  >{\raggedright\arraybackslash}p{(\linewidth - 14\tabcolsep) * \real{0.1250}}
  >{\raggedright\arraybackslash}p{(\linewidth - 14\tabcolsep) * \real{0.1250}}
  >{\raggedright\arraybackslash}p{(\linewidth - 14\tabcolsep) * \real{0.1250}}
  >{\raggedright\arraybackslash}p{(\linewidth - 14\tabcolsep) * \real{0.1250}}
  >{\raggedright\arraybackslash}p{(\linewidth - 14\tabcolsep) * \real{0.1250}}
  >{\raggedright\arraybackslash}p{(\linewidth - 14\tabcolsep) * \real{0.1250}}@{}}
\toprule\noalign{}
\begin{minipage}[b]{\linewidth}\raggedright
Study
\end{minipage} & \begin{minipage}[b]{\linewidth}\raggedright
Cancer Type
\end{minipage} & \begin{minipage}[b]{\linewidth}\raggedright
N
\end{minipage} & \begin{minipage}[b]{\linewidth}\raggedright
Deaths
\end{minipage} & \begin{minipage}[b]{\linewidth}\raggedright
z
\end{minipage} & \begin{minipage}[b]{\linewidth}\raggedright
p
\end{minipage} & \begin{minipage}[b]{\linewidth}\raggedright
Δ OS (mo)
\end{minipage} & \begin{minipage}[b]{\linewidth}\raggedright
Direction
\end{minipage} \\
\midrule\noalign{}
\endhead
\bottomrule\noalign{}
\endlastfoot
MSK-CHORD (Nature 2024) & Mixed & 20,914 & 10,061 & −5.912 &
\textless0.00001 & −2.6 & Balance → shorter \\
Diffuse Glioma (GLASS) & Glioma & 256 & 192 & +5.138 & \textless0.00001
& +29.0 & Balance → longer \\
Diffuse Glioma (GLASS, Nature 2019) & Glioma & 189 & 141 & +5.122 &
\textless0.00001 & +31.5 & Balance → longer \\
Glioma (MSK, Clin Cancer Res 2019) & Glioma & 784 & 287 & +4.353 &
0.00001 & +10.6 & Balance → longer \\
Cholangio-carcinoma (MSK) & Cholangio-carcinoma & 292 & 220 & −4.015 &
0.00006 & −11.9 & Balance → shorter \\
Metastatic CRC (MSK) & Colorectal & 978 & 344 & −3.957 & 0.00008 & −4.6
& Balance → shorter \\
Prostate Cancer (MSK, 2024) & Prostate & 1,378 & 665 & +3.788 & 0.00015
& +5.2 & Balance → longer \\
TMB/Immuno-therapy (MSK) & Mixed & 1,387 & 684 & +3.673 & 0.00024 & +3.0
& Balance → longer \\
Endometrial (TCGA PanCancer) & Uterine & 510 & 83 & +3.189 & 0.00143 &
+8.4 & Balance → longer \\
Anaplastic Thyroid (MSK) & Thyroid & 94 & 44 & +3.044 & 0.00234 & +40.6
& Balance → longer \\
Diffuse Glioma (TCGA GDC) & Glioma & 462 & 104 & +3.025 & 0.00249 & +7.7
& Balance → longer \\
Pancreatic cfDNA (MSK) & Pancreatic & 227 & 186 & −2.916 & 0.00355 &
−7.5 & Balance → shorter \\
LGG (TCGA PanCancer) & Glioma & 462 & 106 & +2.910 & 0.00362 & +7.5 &
Balance → longer \\
Ovarian (TCGA PanCancer) & Ovarian & 388 & 222 & −2.859 & 0.00425 &
−14.5 & Balance → shorter \\
LGG (TCGA Firehose) & Glioma & 262 & 63 & +2.846 & 0.00442 & +10.1 &
Balance → longer \\
Breast Cancer (MSK, 2025) & Breast & 2,653 & 1,447 & +2.637 & 0.00836 &
+7.2 & Balance → longer \\
Adreno-cortical (TCGA Firehose) & Adreno-cortical & 52 & 22 & −2.635 &
0.00840 & −29.8 & Balance → shorter \\
Uterine Carcinosarcoma (TCGA) & Uterine & 56 & 34 & +2.581 & 0.00985 &
+13.5 & Balance → longer \\
Breast (METABRIC) & Breast & 1,520 & 911 & +2.502 & 0.01236 & +9.8 &
Balance → longer \\
LGG+GBM Merged (TCGA, Cell 2016) & Glioma & 591 & 175 & +2.499 & 0.01244
& +4.1 & Balance → longer \\
Melanoma (TCGA PanCancer) & Melanoma & 404 & 200 & +2.218 & 0.02654 &
+13.8 & Balance → longer \\
CRC (CAS Shanghai) & Colorectal & 124 & 23 & −2.214 & 0.02686 & −29.6 &
Balance → shorter \\
MSK-IMPACT 50K & Mixed & 35,566 & 14,677 & +2.208 & 0.02724 & +0.6 &
Balance → longer \\
HGSOC (TCGA GDC) & Ovarian & 380 & 215 & −2.145 & 0.03194 & −7.6 &
Balance → shorter \\
Melanoma (TCGA, Cell 2015) & Melanoma & 261 & 113 & +2.099 & 0.03579 &
+14.9 & Balance → longer \\
IDH-mut Glioma (MSK) & Glioma & 63 & 15 & −2.092 & 0.03645 & −11.2 &
Balance → shorter \\
Prostate (MSK, Clin Cancer Res) & Prostate & 871 & 338 & −2.090 &
0.03664 & −4.5 & Balance → shorter \\
Pancreatic (CPTAC GDC) & Pancreatic & 119 & 119 & −2.076 & 0.03785 &
−9.6 & Balance → shorter \\
Bladder Urothelial (TCGA PanCancer) & Bladder & 375 & 166 & +2.009 &
0.04449 & +2.6 & Balance → longer \\
\end{longtable}
}

\normalsize
\newpage

\subsection{Appendix B: Gene-to-Hallmark
Mapping}\label{appendix-b-gene-to-hallmark-mapping}

{\def\LTcaptype{none} % do not increment counter
\begin{longtable}[]{@{}lll@{}}
\toprule\noalign{}
Gene & Entrez ID & Hallmark(s) \\
\midrule\noalign{}
\endhead
\bottomrule\noalign{}
\endlastfoot
KRAS & 3845 & Proliferative Signaling \\
BRAF & 673 & Proliferative Signaling \\
NRAS & 4893 & Proliferative Signaling \\
ERBB2 & 2064 & Proliferative Signaling \\
ERBB3 & 2065 & Proliferative Signaling \\
MAP2K1 & 5604 & Proliferative Signaling \\
FGFR2 & 2263 & Proliferative Signaling \\
NF1 & 4763 & Proliferative Signaling \\
PIK3CA & 5290 & Proliferative Signaling, Resisting Cell Death \\
TP53 & 7157 & Growth Suppression Evasion, Resisting Cell Death \\
APC & 324 & Growth Suppression Evasion \\
RB1 & 5925 & Growth Suppression Evasion \\
CDKN2A & 1029 & Growth Suppression Evasion \\
SMAD4 & 4089 & Growth Suppression Evasion \\
SMAD2 & 4087 & Growth Suppression Evasion \\
SMAD3 & 4088 & Growth Suppression Evasion \\
TGFBR2 & 7048 & Growth Suppression Evasion \\
FBXW7 & 55294 & Growth Suppression Evasion \\
ACVR2A & 92 & Growth Suppression Evasion \\
NOTCH1 & 4851 & Growth Suppression Evasion \\
SOX9 & 6662 & Growth Suppression Evasion \\
TCF7L2 & 6934 & Growth Suppression Evasion \\
RNF43 & 54894 & Growth Suppression Evasion \\
AMER1 & 139285 & Growth Suppression Evasion \\
CTNNB1 & 1499 & Growth Suppression Evasion, Invasion \& Metastasis \\
PTEN & 5728 & Resisting Cell Death \\
CREBBP & 1387 & Resisting Cell Death \\
TSC2 & 7249 & Resisting Cell Death \\
ATRX & 546 & Replicative Immortality, Genome Instability \\
CDH1 & 999 & Invasion \& Metastasis \\
FAT1 & 2195 & Invasion \& Metastasis \\
FAT4 & 79633 & Invasion \& Metastasis \\
JAK1 & 3716 & Immune Evasion \\
MLH1 & 4292 & Genome Instability \\
MSH2 & 4436 & Genome Instability \\
MSH6 & 2956 & Genome Instability \\
PMS2 & 5395 & Genome Instability \\
POLE & 5426 & Genome Instability \\
BRCA2 & 675 & Genome Instability \\
ATM & 472 & Genome Instability \\
ARID1A & 8289 & Genome Instability \\
KMT2C & 58508 & Genome Instability \\
KMT2D & 8085 & Genome Instability \\
IDH1 & 3417 & Deregulated Metabolism \\
GNAS & 2778 & Deregulated Metabolism \\
\end{longtable}
}

\subsection{Appendix C: The Master Equation --- Evolutionary Dynamics on
the Hallmark
Hypercube}\label{appendix-c-the-master-equation-evolutionary-dynamics-on-the-hallmark-hypercube}

The preceding 23 sections characterize cancer at equilibrium --- the
photograph. Appendix C makes the photograph move. We treat the 256
possible hallmark states (\(\sigma \in \{0,1\}^8\)) as states of a
continuous-time Markov chain on the 8-dimensional hypercube, infer
transition rates from the observed stationary distribution, and solve
for the eigensystem that governs tumor evolution.

\textbf{The free energy landscape.} Interpreting the observed state
frequencies \(\pi(\sigma)\) as a stationary distribution yields a free
energy \(F(\sigma) = -\ln \pi(\sigma)\) over the 162 observed states,
with 94 forbidden states acting as infinite-energy barriers. The ground
state --- the lowest free energy, most thermodynamically favored
configuration --- is \(\sigma = 01100000\) (GrowSup + Death): the TP53
state, with \(n = 41{,}198\) patients (20.7\% of all tumors) and
\(F = 1.57\). The TP53 bond is not merely the strongest interaction term
(\(z = +13.2\)); it is the starting position of cancer.

\textbf{Transition matrix.} Single-hallmark flips define edges on the
hypercube, with an additional edge for the TP53 double-flip
(simultaneous GrowSup + Death acquisition). Metropolis rates
\(W(\sigma \to \sigma') = \min(1, \pi(\sigma')/\pi(\sigma))\) satisfy
detailed balance and yield a generator matrix \(L\) with 1,212 directed
edges (1,070 single-flip + 142 TP53-bond) among the 162 observed states.

\textbf{The spectrum.} The eigenvalue decomposition of the
\(162 \times 162\) generator reveals:

{\def\LTcaptype{none} % do not increment counter
\begin{longtable}[]{@{}
  >{\raggedright\arraybackslash}p{(\linewidth - 4\tabcolsep) * \real{0.3000}}
  >{\raggedleft\arraybackslash}p{(\linewidth - 4\tabcolsep) * \real{0.4000}}
  >{\raggedright\arraybackslash}p{(\linewidth - 4\tabcolsep) * \real{0.3000}}@{}}
\toprule\noalign{}
\begin{minipage}[b]{\linewidth}\raggedright
Eigenvalue
\end{minipage} & \begin{minipage}[b]{\linewidth}\raggedleft
Value
\end{minipage} & \begin{minipage}[b]{\linewidth}\raggedright
Meaning
\end{minipage} \\
\midrule\noalign{}
\endhead
\bottomrule\noalign{}
\endlastfoot
\(\mu_0\) & \(0.000\) & Stationary (the periodic table) \\
\(\mu_1\) & \(-0.600\) & Slowest relaxation mode \\
\(\mu_2\) & \(-0.651\) & Second mode \\
Spectral gap \(|\mu_1|\) & \(0.600\) & Rate-limiting timescale \\
Relaxation time \(1/|\mu_1|\) & \(1.67\) & Markov-time units to
equilibrium \\
Kemeny constant \(K\) & \(43.6\) & Expected hitting time to random
target \\
\end{longtable}
}

\textbf{The basin decomposition.} The first non-trivial eigenvector
\(v_1\) splits the 162 observed states into two metastable basins:

{\def\LTcaptype{none} % do not increment counter
\begin{longtable}[]{@{}lll@{}}
\toprule\noalign{}
& Basin A (+) & Basin B (−) \\
\midrule\noalign{}
\endhead
\bottomrule\noalign{}
\endlastfoot
States & 15 & 147 \\
Patients & 120,635 (60.7\%) & 78,227 (39.3\%) \\
Mean hallmarks \(\langle k \rangle\) & 2.1 & 4.1 \\
Mean log-hazard & \(+0.152\) (dangerous) & \(-0.198\) (safer) \\
Immort prevalence & 0.0\% & 13.9\% \\
Genome prevalence & 0.0\% & 92.3\% \\
Immune prevalence & 0.0\% & 6.7\% \\
\end{longtable}
}

Basin A contains 15 states with 61\% of all patients --- tumors with few
hallmarks, dominated by the periphery (GrowSup, Death, Prolif), and zero
representation of Immort, Immune, or Genome. These are early-stage
tumors with HIGH hazard. Basin B contains 147 states --- tumors that
have crossed the barrier and acquired core hallmarks, paradoxically
showing LOWER hazard. The core creates internal conflict that slows
tumor progression, consistent with the antagonistic Core \(\times\)
Periphery interaction (\(z = -22.3\)) found in §3.20.

\textbf{The hallmark acquisition order.} Population-weighted free energy
gradients reveal the thermodynamically favored sequence of hallmark
acquisition:

{\def\LTcaptype{none} % do not increment counter
\begin{longtable}[]{@{}rlrl@{}}
\toprule\noalign{}
Order & Hallmark & \(\langle \Delta F \rangle\) & Category \\
\midrule\noalign{}
\endhead
\bottomrule\noalign{}
\endlastfoot
1 & GrowSup & \(+0.14\) & Periphery \\
2 & Death & \(+0.24\) & Periphery \\
3 & Prolif & \(+0.33\) & Periphery \\
4 & Genome & \(+1.04\) & \textbf{Bottleneck} \\
5 & Immort & \(+1.91\) & Core \\
6 & Invade & \(+1.98\) & Core \\
7 & Metab & \(+3.50\) & Core \\
8 & Immune & \(+4.33\) & Core \\
\end{longtable}
}

A \(3.2\times\) free energy gap separates periphery hallmarks
(\(\langle \Delta F \rangle \leq 0.33\)) from core hallmarks
(\(\langle \Delta F \rangle \geq 1.04\)). The periphery (TP53 + Prolif)
is acquired first and cheaply; the core follows slowly and at high
energetic cost. Genome Instability sits at the boundary --- the fourth
hallmark acquired and the dominant bottleneck.

\textbf{The bottleneck is Genome Instability.} Of 45 basin-boundary
crossing edges (transitions where a state moves from Basin A to Basin
B), 15 involve acquiring Genome Instability --- more than any other
hallmark (tied with Metab and Immune, which also appear in 15 crossings
but only in compound transitions). The top-10 highest-rate bottleneck
transitions all involve Genome acquisition. This is the rate-limiting
step in tumor evolution: the transition from a periphery-dominated,
high-hazard tumor to a core-enriched, lower-hazard tumor.

\textbf{Monte Carlo trajectory forecasts.} Simulating 10,000 random
walks from each of the 20 most common states, with the lethal manifold
(top quartile by Cox hazard) and forbidden boundary as absorbing states:

{\def\LTcaptype{none} % do not increment counter
\begin{longtable}[]{@{}lllrr@{}}
\toprule\noalign{}
State & Hallmarks & First acquisition & \% Lethal & \% Boundary \\
\midrule\noalign{}
\endhead
\bottomrule\noalign{}
\endlastfoot
01100000 & GrowSup, Death & Prolif (53\%) & 93.5 & 6.5 \\
10000000 & Prolif & Death (71\%) & 95.7 & 4.3 \\
00000010 & Genome & GrowSup (64\%) & 86.9 & 13.1 \\
01000000 & GrowSup & Death (54\%) & 80.3 & 19.7 \\
01001000 & GrowSup, Invade & Death (30\%) & 70.4 & 29.6 \\
\end{longtable}
}

Every state converges toward TP53 as its first or second acquisition.
Single-hallmark tumors Race toward the iron bond with \textgreater50\%
probability per step. The lethal manifold absorbs 70--100\% of
trajectories; only GrowSup + Invade (29.6\%) and GrowSup + Death +
Genome (27.5\%) show substantial escape to the forbidden boundary.

\textbf{Navigation chart structure.} For each state, the eigensystem
produces three actionable quantities: (1) the \emph{kill arrow} --- the
transition maximizing hazard increase (universally: acquiring the TP53
bond, \(\Delta\lambda = +0.43\)); (2) the \emph{trap arrow} --- the
transition leading to a forbidden neighbor (most commonly: acquiring
Immortality, which leads to unobserved configurations); and (3) the
\emph{gradient} --- the most thermodynamically likely next acquisition.
The therapeutic logic: block the kill arrow, enable the trap arrow.

\textbf{Per-cancer evolutionary landscapes.} The master equation was
solved independently for the 10 largest cancer types. The results
confirm tissue-specific dynamics:

{\def\LTcaptype{none} % do not increment counter
\begin{longtable}[]{@{}
  >{\raggedright\arraybackslash}p{(\linewidth - 8\tabcolsep) * \real{0.1875}}
  >{\raggedright\arraybackslash}p{(\linewidth - 8\tabcolsep) * \real{0.1875}}
  >{\raggedleft\arraybackslash}p{(\linewidth - 8\tabcolsep) * \real{0.2500}}
  >{\raggedright\arraybackslash}p{(\linewidth - 8\tabcolsep) * \real{0.1875}}
  >{\raggedright\arraybackslash}p{(\linewidth - 8\tabcolsep) * \real{0.1875}}@{}}
\toprule\noalign{}
\begin{minipage}[b]{\linewidth}\raggedright
Cancer
\end{minipage} & \begin{minipage}[b]{\linewidth}\raggedright
Ground state
\end{minipage} & \begin{minipage}[b]{\linewidth}\raggedleft
\(T_\mathrm{relax}\)
\end{minipage} & \begin{minipage}[b]{\linewidth}\raggedright
Earliest acquired
\end{minipage} & \begin{minipage}[b]{\linewidth}\raggedright
Latest acquired
\end{minipage} \\
\midrule\noalign{}
\endhead
\bottomrule\noalign{}
\endlastfoot
Mixed (\(n = 116{,}352\)) & GrowSup+Death & 1.59 & GrowSup & Immune \\
Breast (\(n = 9{,}763\)) & GrowSup+Death & 2.17 & Death & Metab \\
Prostate (\(n = 5{,}619\)) & GrowSup+Death & \textbf{7.85} & GrowSup &
Metab \\
Lung adeno (\(n = 5{,}272\)) & GrowSup+Death & \textbf{1.30} & Prolif &
Immune \\
Glioma (\(n = 4{,}351\)) & GrowSup+Death+Immort+Genome+Metab &
\(\infty\) & Immort & Immune \\
Colorectal (\(n = 4{,}596\)) & Prolif+GrowSup+Death & 1.85 & GrowSup
(\(\Delta F = -1.95\)) & Immune \\
\end{longtable}
}

Prostate cancer relaxes 6\(\times\) slower than lung adenocarcinoma
(\(T = 7.85\) vs.~\(1.30\)), consistent with the clinical observation
that prostate tumors evolve slowly while lung tumors progress rapidly.
Glioma has a zero spectral gap and a completely different ground state
--- the IDH1 bifurcation loads the tumor directly into the core (Immort
+ Genome + Metab), bypassing the periphery-first pathway entirely.
Colorectal cancer shows GrowSup acquisition as strongly \emph{downhill}
(\(\Delta F = -1.95\), the only negative value in any cancer type),
reflecting the obligate APC → TP53 adenoma--carcinoma sequence.

\subsection{Appendix D:
Reproducibility}\label{appendix-d-reproducibility}

\textbf{Data source.} All mutation and survival data were obtained from
the cBioPortal for Cancer Genomics public API
(https://www.cbioportal.org/api), accessed March 2026. No
patient-identifiable information was accessed or stored.

\textbf{Computational environment.} All analyses were performed on a
consumer workstation: AMD Ryzen / Windows 11. Software: Python 3.13,
Node.js 20 LTS, NumPy 2.4, lifelines 0.30 (Cox PH with Efron tie
handling), SciPy 1.12.

\textbf{Pipeline.} The analysis pipeline processes 198,862 patients
across 378 studies in four phases:

\begin{enumerate}
\def\labelenumi{\arabic{enumi}.}
\tightlist
\item
  \emph{Data retrieval} --- Study selection, mutation download, and
  patient deduplication via the cBioPortal API.
\item
  \emph{Hallmark encoding} --- 45-gene panel mapped to 8 hallmarks
  (Appendix B), producing the 8-bit hallmark vector
  \(\sigma \in \{0,1\}^8\) per patient.
\item
  \emph{Thermodynamic analysis} --- Boltzmann fitting, Ising coupling
  matrix estimation, per-cancer temperature computation, forbidden state
  census, and domain wall encoding.
\item
  \emph{Survival modeling} --- Cox proportional hazards regression,
  cross-validation, permutation testing, and the strong inference
  battery (§3.30).
\end{enumerate}

\textbf{Validation.} A self-contained validation script
(\texttt{patient\_compass\_validate.js}) executes 21,308 assertions
covering all 37 cancer types, all 162 observed hallmark states, all
O\(_2\) tier assignments, and all Cox model coefficients reported in
this paper. The script passes with 0 failures.

\textbf{Independent reproducibility script.} A standalone Python script
(\texttt{compass\_reproduce.py}) re-derives every statistical claim in
this paper --- including the Cox proportional hazards wall coefficients
(§3.23) --- using only standard biostatistical methods (Mann-Whitney,
Stouffer meta-analysis, Cox PH via lifelines). The script requires NumPy
and lifelines, and can fetch data directly from the cBioPortal API. It
reproduces the 10-wall + 2-O\(_2\)-indicator Cox model on 48,675
O\(_2\)-classified patients, recovering all 10 published wall
coefficients to six decimal places (\(C = 0.6242\)). It contains no
eigenvalue computation, no spectral methods, and no proprietary
algorithms.

\textbf{Relationship to companion patents.} The hallmark pair structure,
domain wall decomposition, and O\(_2\) tier classification reported here
were originally identified using proprietary spectral methods described
in U.S. Provisional Patent Applications No.~64/031,440 {[}4{]} and
No.~64/033,689 {[}5{]}, assigned to Fancyland LLC. These methods
accelerated the discovery process but are not required for verification:
all statistical claims in this paper are independently reproducible
using only standard biostatistical methods and publicly available
cBioPortal data, as demonstrated by the accompanying reproducibility
script.

\textbf{Code and interactive demo.} The gene-to-hallmark mapping,
O\(_2\) classification, analysis code, reproducibility script, and an
interactive patient-facing report are available at:

\begin{quote}
\textbf{\href{https://github.com/fancyland-llc/patient-compass-pan-cancer}{github.com/fancyland-llc/patient-compass-pan-cancer}}
\end{quote}

The interactive compass demo renders all findings from this paper as a
bilingual (English/Spanish) clinical report, freely accessible with no
software installation. It contains no patient-identifiable information.

\vspace{3em}
\begin{center}
\rule{0.4\textwidth}{0.5pt}

\vspace{0.8em}

\textit{Fancyland LLC --- Lattice OS research infrastructure.}

\vspace{0.3em}

\textit{The rabbit has been caught.}
\end{center}

\end{document}
